% \iffalse meta-comment
%
% Copyright (C) 2019 by Zangwei Zheng <zhengzangw@gmail.com>
%
% This file may be distributed and/or modified under the conditions of
% the LaTeX Project Public License, either version 1.3c of this license
% or (at your option) any later version. The latest version of this 
% license is in:
%
%   http://www.latex-project.org/lppl.txt
%
% and version 1.3c or later is part of all distributions of LaTeX
% version 2005/12/01 or later.
%
% \fi
%
% \iffalse
%<*driver>
\ProvidesFile{njurepo.dtx}[2019/01/29 1.1.2 Nanjing University Report Template]
\documentclass{ltxdoc}
\usepackage{dtx-style}
    \EnableCrossrefs
    \CodelineIndex
    \RecordChanges
\begin{document}
    \DocInput{njurepo.dtx}
    \PrintChanges
    \PrintIndex
\end{document}
%</driver>
% \fi
%
% \CheckSum{0}
%
% \CharacterTable
%  {Upper-case    \A\B\C\D\E\F\G\H\I\J\K\L\M\N\O\P\Q\R\S\T\U\V\W\X\Y\Z
%   Lower-case    \a\b\c\d\e\f\g\h\i\j\k\l\m\n\o\p\q\r\s\t\u\v\w\x\y\z
%   Digits        \0\1\2\3\4\5\6\7\8\9
%   Exclamation   \!     Double quote  \"     Hash (number) \#
%   Dollar        \$     Percent       \%     Ampersand     \&
%   Acute accent  \'     Left paren    \(     Right paren   \)
%   Asterisk      \*     Plus          \+     Comma         \,
%   Minus         \-     Point         \.     Solidus       \/
%   Colon         \:     Semicolon     \;     Less than     \<
%   Equals        \=     Greater than  \>     Question mark \?
%   Commercial at \@     Left bracket  \[     Backslash     \\
%   Right bracket \]     Circumflex    \^     Underscore    \_
%   Grave accent  \`     Left brace    \{     Vertical bar  \|
%   Right brace   \}     Tilde         \~}
%
% \DoNotIndex{\newenvironment,\@bsphack,\@empty,\@esphack,\sfcode}
% \DoNotIndex{\addtocounter,\label,\let,\linewidth,\newcounter}
% \DoNotIndex{\noindent,\normalfont,\par,\parskip,\phantomsection}
% \DoNotIndex{\providecommand,\ProvidesPackage,\refstepcounter}
% \DoNotIndex{\RequirePackage,\setcounter,\setlength,\string,\strut}
% \DoNotIndex{\textbackslash,\texttt,\ttfamily,\usepackage}
% \DoNotIndex{\begin,\end,\begingroup,\endgroup,\par,\\}
% \DoNotIndex{\if,\ifx,\ifdim,\ifnum,\ifcase,\else,\or,\fi}
% \DoNotIndex{\let,\def,\xdef,\edef,\newcommand,\renewcommand}
% \DoNotIndex{\expandafter,\csname,\endcsname,\relax,\protect}
% \DoNotIndex{\Huge,\huge,\LARGE,\Large,\large,\normalsize}
% \DoNotIndex{\small,\footnotesize,\scriptsize,\tiny}
% \DoNotIndex{\normalfont,\bfseries,\slshape,\sffamily,\interlinepenalty}
% \DoNotIndex{\textbf,\textit,\textsf,\textsc}
% \DoNotIndex{\hfil,\par,\hskip,\vskip,\vspace,\quad}
% \DoNotIndex{\centering,\raggedright,\ref}
% \DoNotIndex{\c@secnumdepth,\@startsection,\@setfontsize}
% \DoNotIndex{\ ,\@plus,\@minus,\p@,\z@,\@m,\@M,\@ne,\m@ne}
% \DoNotIndex{\@@par,\DeclareOperation,\RequirePackage,\LoadClass}
% \DoNotIndex{\AtBeginDocument,\AtEndDocument}
%
% \changes{v1.0.0}{2019/01/22}{Initial version}
% \changes{v1.0.1}{2019/01/29}{Add more ability}
% \changes{v1.1.0}{2019/01/29}{Stable version}
% \changes{v1.1.1}{2019/02/20}{Fix little bugs}
% \changes{v1.1.2}{2019/03/20}{Fix some typo and add more predined macro. Polish manual}
%
% \GetFileInfo{\jobname.dtx}
%
% \def\indexname{索引}
% \def\glossaryname{修改记录}
% \IndexPrologue{\section{\indexname}}
% \GlossaryPrologue{\section{\glossaryname}}

% \title{\bfseries\color{violet}\njurepo： 南京大学本科生泛用报告}
% \author{郑奘巍 \\[5pt]\texttt{zhengzangw@gmail.com}}
% \date{\fileversion\ (\filedate)}
% \maketitle\thispagestyle{empty}
%
% \begin{abstract}\noindent
% 此宏包旨在建立一个免于配置的、指令相对简单的南京大学作业、实验报告通用模板。
% \end{abstract}
%
%
% \vskip2cm
% \def\abstractname{免责声明}
% \begin{abstract}
% \noindent
% \begin{enumerate}
% \item 本模板的发布遵守 \LaTeX\ Project Public License，使用前请认真阅读协议内
%   容。
% \item \textbf{本模板为作者自己通常使用的报告模板，与南京大学官方没有任何关系}。任何使用该宏包进行实验报告制作时，请\textbf{务必根据课程要求进行写作}。由于使用本模板而引起的作业验收问题，均与本模板作者无关。
% \item 本模板借鉴\thuthesis{}宏包的许多内容，需要稳定模板的同学也可以选择使用清华大学的\thuthesis{}宏包并自己进行配置。
% \item 任何个人或组织以本模板为基础进行修改、扩展而生成的新的专用模板，请严格遵
%   守 \LaTeX\ Project Public License 协议。由于违犯协议而引起的任何纠纷争端均与
%   本模板作者无关。
% \end{enumerate}
% \end{abstract}
%
% \clearpage
% \pagestyle{fancy}
% \begin{multicols}{2}[
%   \setlength{\columnseprule}{.4pt}
%   \setlength{\columnsep}{18pt}]
%   \tableofcontents
% \end{multicols}
% \clearpage
%
% \section{模板介绍}
% \njurepo\ (\textbf{N}an\textbf{jing} \textbf{U}niversity \LaTeX\ Versatile \textbf{Repo}rt Template)是为南京大学本科生设计的一个免于配置的作业、实验报告模板。希望它可以使你的作业/实验报告不会因形式上的缺陷导致评分的下降。
%
% 该文档格式基于 ctexbook, 主要完成了除了主体内容以外的几乎\textbf{全部}工作。同时，通过使用 Github 版本宏包，你还可以更好的管理自己的 $\LaTeX$ 文档。
%
% 本文档将尽量完整的介绍模板的使用方法，如有不清楚之处可以参考示例文档或者根据第 3.1 节说明提问，有兴趣者都可以参与完善此手册，也非常欢迎对代码的贡献。
% 
% \section{安装}
% \label{sec:installation}
% \subsection{CTAN}
% 本宏包已被收纳于 CTAN 中，凡安装完整版 Texlive 用户可直接使用：
% \begin{latex} 
% \usepackage{njurepo} 
% \end{latex}
%
% \subsection{Github}
% 想获得最新版本的 \njurepo 请前往 Github 主页下载：https://github.com/zhengzangw/njurepo 
% 
% 同时，面向 Github 版本的说明主要在 README.md 中，如果使用的是 Github 版本，请阅读 README.md
%
% \subsection{模板的组成}
% 下表列出了\njurepo 的主要文件及其功能介绍：
%
% \begin{longtable}{l|p{8cm}}
% \toprule
% {\heiti 文件（夹）} & {\heiti 功能描述}\\\midrule
% \endfirsthead
% \midrule
% {\heiti 文件（夹）} & {\heiti 功能描述}\\\midrule
% \endhead
% \endfoot
% \endlastfoot
% njurepo.ins & \textsc{DocStrip} 驱动文件（开发用） \\
% njurepo.dtx & \textsc{DocStrip} 源文件（开发用）\\\midrule
% resources/ & 资源路径\\
% resources/logos/ & 示例文档资源路径\\
% parts/ & 具体内容\\
% parts/examples/ & 示例文档具体内容\\
% ref/ & 参考文献和参考文献样式文件\\
% njurepo.cls & 模板类文件\\
% Makefile & 自动运行脚本 \\
% util.py & 实用管理工具 \\
% README.md & 说明文件 \\
% \textbf{njurepo.pdf} & 用户手册（本文档）\\ \bottomrule
% \end{longtable}
%
% \subsection{生成模板}
% 使用Makefile或\XeLaTeX 生成模板文件
% \begin{shell}
% make cls
% \end{shell}
%
% \subsection{生成论文}
% \subsubsection{latexmk}
% latexmk 命令支持全自动生成\LaTeX{}编写的文档，并且支持使用不同的工具链来进行生成，它会自动运行多次工具直到交叉引用都被解决。下面给出了一个用 latexmk 调用 xelatex 生成最终文档的示例:
% \begin{shell}
% latexmk -xelatex main
% \end{shell}
%
% \subsection{升级}
% 在github上下载最新版，运行:
% \begin{shell}
% make cls
% \end{shell}
%
% 生成新的类文件和配置文件即可。也可以直接拷贝 njurepo.cls，免去上面命令的执行。
% 
%
% \section{使用说明}
% \subsection{示例文件}
% 推荐从模板自带的示例文档入手，其中包括了论文写作用到的所有命令及其使用方法，只需要用自己的内容进行相应替换就可以。对于不清楚的命令可以查阅本手册。下面的例子描述了模板中章节的组织形式，来自于示例文档，具体内容可以参考模板附带的 main.tex 和 parts/examples/。
% \begin{latex}
% \documentclass[language=english,open=any]{njurepo}
% \begin{document}
% \frontmatter
% \njusetup{
  %******************************
  % 注意：
  %   1. 配置里面不要出现空行
  %   2. 不需要的配置信息可以删除
  %******************************
  %
  %=========
  % 中文信息
  %=========
  ctitle={清华大学学位论文 \LaTeX\ 模板\\使用示例文档},
  csubtitle={通用报告模板},
  cdegree={工学硕士},
  cdepartment={计算机科学与技术系},
  cmajor={计算机科学与技术},
  cauthor={薛瑞尼},
  stdid={12345678},
  csupervisor={郑纬民教授},
  cassosupervisor={陈文光教授}, % 副指导老师
  ccosupervisor={某某某教授}, % 联合指导老师
  % 日期自动使用当前时间，若需指定按如下方式修改：
  % cdate={超新星纪元},
  %
  %=========
  % 英文信息
  %=========
  etitle={An Introduction to \LaTeX{} Thesis Template of Tsinghua University v\version},
  esubtitle={subtile here},
  %esubsubtitle={Sub},
  edegree={Doctor of Engineering},
  emajor={Computer Science and Technology},
  eauthor={Xue Ruini},
  esupervisor={Professor Zheng Weimin},
  eassosupervisor={Chen Wenguang},
  % 日期自动生成，若需指定按如下方式修改：
  % edate={December, 2005}
  %
  % 关键词用“英文逗号”分割
  ckeywords={\TeX, \LaTeX, CJK, 模板, 论文},
  ekeywords={\TeX, \LaTeX, CJK, template, thesis}
}


% % 定义中英文摘要和关键字
\begin{cabstract}
    论文的摘要是对论文研究内容和成果的高度概括。摘要应对论文所研究的问题及其研究目
    的进行描述，对研究方法和过程进行简单介绍，对研究成果和所得结论进行概括。摘要应
    具有独立性和自明性，其内容应包含与论文全文同等量的主要信息。使读者即使不阅读全
    文，通过摘要就能了解论文的总体内容和主要成果。
  
    论文摘要的书写应力求精确、简明。切忌写成对论文书写内容进行提要的形式，尤其要避
    免“第 1 章……；第 2 章……；……”这种或类似的陈述方式。
  
    本文介绍清华大学论文模板 \thuthesis{} 的使用方法。本模板符合学校的本科、硕士、
    博士论文格式要求。
  
    本文的创新点主要有：
    \begin{itemize}
      \item 用例子来解释模板的使用方法；
      \item 用废话来填充无关紧要的部分；
      \item 一边学习摸索一边编写新代码。
    \end{itemize}
  
    关键词是为了文献标引工作、用以表示全文主要内容信息的单词或术语。关键词不超过 5
    个，每个关键词中间用分号分隔。（模板作者注：关键词分隔符不用考虑，模板会自动处
    理。英文关键词同理。）
  \end{cabstract}
  
  \ckeywords{\TeX, \LaTeX, CJK, 模板, 论文}
  
  \begin{eabstract}
     An abstract of a dissertation is a summary and extraction of research work
     and contributions. Included in an abstract should be description of research
     topic and research objective, brief introduction to methodology and research
     process, and summarization of conclusion and contributions of the
     research. An abstract should be characterized by independence and clarity and
     carry identical information with the dissertation. It should be such that the
     general idea and major contributions of the dissertation are conveyed without
     reading the dissertation.
  
     An abstract should be concise and to the point. It is a misunderstanding to
     make an abstract an outline of the dissertation and words ``the first
     chapter'', ``the second chapter'' and the like should be avoided in the
     abstract.
  
     Key words are terms used in a dissertation for indexing, reflecting core
     information of the dissertation. An abstract may contain a maximum of 5 key
     words, with semi-colons used in between to separate one another.
  \end{eabstract}
  
  \ekeywords{\TeX, \LaTeX, CJK, template, thesis}
  
% \maketitlepage
% \makecover
% \makeabstract
% \tableofcontents
% \begin{denotation}[3cm]
\item[HPC] 高性能计算 (High Performance Computing)
\item[cluster] 集群
\item[Itanium] 安腾
\item[SMP] 对称多处理
\item[API] 应用程序编程接口
\item[PI] 聚酰亚胺
\item[MPI] 聚酰亚胺模型化合物，N-苯基邻苯酰亚胺
\item[PBI] 聚苯并咪唑
\item[MPBI] 聚苯并咪唑模型化合物，N-苯基苯并咪唑
\item[PY] 聚吡咙
\item[PMDA-BDA]	均苯四酸二酐与联苯四胺合成的聚吡咙薄膜
\item[$\Delta G$] 活化自由能 (Activation Free Energy)
\item[$\chi$] 传输系数 (Transmission Coefficient)
\item[$E$] 能量
\item[$m$] 质量
\item[$c$] 光速
\item[$P$] 概率
\item[$T$] 时间
\item[$v$] 速度
\item[劝学] 君子曰：学不可以已。青，取之于蓝，而青于蓝；冰，水为之，而寒于水。木
  直中绳。輮以为轮，其曲中规。虽有槁暴，不复挺者，輮使之然也。故木受绳则直，金就
  砺则利，君子博学而日参省乎己，则知明而行无过矣。吾尝终日而思矣，不如须臾之所学
  也；吾尝跂而望矣，不如登高之博见也。登高而招，臂非加长也，而见者远；顺风而呼，
  声非加疾也，而闻者彰。假舆马者，非利足也，而致千里；假舟楫者，非能水也，而绝江
  河，君子生非异也，善假于物也。积土成山，风雨兴焉；积水成渊，蛟龙生焉；积善成德，
  而神明自得，圣心备焉。故不积跬步，无以至千里；不积小流，无以成江海。骐骥一跃，
  不能十步；驽马十驾，功在不舍。锲而舍之，朽木不折；锲而不舍，金石可镂。蚓无爪牙
  之利，筋骨之强，上食埃土，下饮黄泉，用心一也。蟹六跪而二螯，非蛇鳝之穴无可寄托
  者，用心躁也。—— 荀况
\end{denotation}



% % 也可以使用 nomencl 宏包：

% \printnomenclature[3cm]

% \nomenclature{HPC}{高性能计算 (High Performance Computing)}
% \nomenclature{cluster}{集群}
% \nomenclature{Itanium}{安腾}
% \nomenclature{SMP}{对称多处理}
% \nomenclature{API}{应用程序编程接口}
% \nomenclature{PI}{聚酰亚胺}
% \nomenclature{MPI}{聚酰亚胺模型化合物，N-苯基邻苯酰亚胺}
% \nomenclature{PBI}{聚苯并咪唑}
% \nomenclature{MPBI}{聚苯并咪唑模型化合物，N-苯基苯并咪唑}
% \nomenclature{PY}{聚吡咙}
% \nomenclature{PMDA-BDA}{均苯四酸二酐与联苯四胺合成的聚吡咙薄膜}
% \nomenclature{$\Delta G$}{活化自由能 (Activation Free Energy)}
% \nomenclature{$\chi$}{传输系数 (Transmission Coefficient)}
% \nomenclature{$E$}{能量}
% \nomenclature{$m$}{质量}
% \nomenclature{$c$}{光速}
% \nomenclature{$P$}{概率}
% \nomenclature{$T$}{时间}
% \nomenclature{$v$}{速度}

% \mainmatter
% \maketitle
% \section*{\scshape Part I\ \ \ Homework}
\begin{psproblem}[26.1.1]
    We will show that each
\end{psproblem}
\begin{pssolution}
    We will show that each flow in G' is mapping to a flow of same value in G and vice versa. Thus the maximum flow in G' has the same value as a maximum flow in G.

    \marginpar{this is marginpar}
    For a flow in G, consider the flow f, in which there exist antiparallel edge $v_i->v_j$ and $v_j->v_i$. Suppose the flow on $v_i->v_j$ is $F$ and $F'$ in another direction. The flow f' in G' has same flow in each edge except that $v_i->v_j$ is replaced by $v_i->v'$ and $v'->v_j$ and the flow is $F$. It is obvious that Capacity constraint and Flow conservation meets. The value of f' is 
        $$|f'|=\sum_{v\in V'}f'(s,v)-\sum{v\in V'}f'(v,s)$$ 
    If there is no antiparallel edges incident with s, the value is same. Otherwise, $f(s,v)=f'(s,v')$ the value is still the same.

    For the other side, due to Flow conservation, $f'(v_1,v')=f'(v',v_2)$. By deleting $v'$ and add edge $v_1\rightarrow v_2$ with $f(v_1,v_2)=f'(v_1,v')$ we can get the graph G. Thus we have built a bijection $G:\mathbb{F}\rightarrow\mathbb{F'}$ where $\mathbb{F}$ denotes all possible flow and with the same value. Hence we have proved this problem.
\end{pssolution}

\begin{psproblem}[26.1.2]
\end{psproblem}
\begin{pssolution}
    Deleting the source and the sink, we get a unique multiple-source, multiple-sink flow network.

    For the other side, for a multiple-source, multiple-sink flow network, we add two kinds of edges $s\rightarrow s_i$ and $t\rightarrow t_i$. We will prove the flow on each kind of each is unique.
    
    For the first kind, by flow conservation , $f(s, s_i)=\sum f(s_i, v)$, where v is vertices incident with $s_i$ except $s$. Thus the value is unique. The same proof applied for the second kind.

    Thus we have proved the problem.
\end{pssolution}

\begin{psproblem}[26.1.6]
\end{psproblem}
\begin{pssolution}
    Build the graph in which corners are vertices and edges (u,v) and (v,u) are added if there is a road connecting them. Constrain the edge to have capacity 1, and the problem is reverted to find a maximum flow and check if it can be 2 and with integer flow value in each edge.

    However, this leads to a problem. The two may take a road in different direction, namely one pass (u,v) and the other (v,u). A direct algorithm is to exhaust all possible direction of each edge and run the above procedure, namely for one road connecting u, v, adding edge (u,v) in one time and (v,u) in another. However, this costs a lot of time, I wonder if there is a better way.
\end{pssolution}

\begin{psproblem}[26.1.7]
\end{psproblem}
\begin{pssolution}
    Build network G' in the following way: for a vertex $v\in G$ with vertex constrant $C$, add two vertices $v_1,v_2$ into G', all edge connected from $u$ to $v$ will lead to edge $u\rightarrow v_1$ in G' and all edge connected from $v$ to $u$ will lead to edge $v_2\rightarrow u$ in G'. Two edges $(v_1,v_2),(v_2,v_1)$ with capacity $C$ is also added, which has flow $\sum_{u}f(u,v)-\sum_{u}f(v,u)$. In this way, the constraint of edge between $v_1$ and $v_2$ serves as constraint of $v$.

    The graph G' has $2\abs{V}$ vertices and $\abs{E}+\abs{V}$ edges.
\end{pssolution}

\begin{psproblem}[26.2.2]
\end{psproblem}
\begin{pssolution}
    The flow is $11+1-4+7+4=19$, the capacity of this cut is $16+4+7+4=31$
\end{pssolution}

\begin{psproblem}[26.2.6]
\end{psproblem}
\begin{pssolution}
    We construct a new flow network by adding a super source s with edge $(s,s_i)$ of capacity $p_i$ and a super sink t with edge $(t_i,t)$ of capacity $q_i$. We will prove the original problem is to find the maximum flow in the new network and check whether it can be $sum_{p_i}$

    First, by conservation of flow in $s_i$, we have $\sum_{v\in V}f(s_i,v)=f(s,s_i)=p_i$. The same proof applies for each $t_i$. Sum over all flow in and out of $V-\{s,\}$, we get $\sum_{i}f(s,s_i)=\sum_{i}f_{t_i,j}$. In the case we find the maximum flow in the new network and check whether it can be $sum_{p_i}$, we can have $\sum_ip_i=\sum_jq_j$. So the convertion is right.
\end{pssolution}

\begin{psproblem}[26.2.8]
\end{psproblem}
\begin{pssolution}
    To prove this problem, we will prove that in the while loop, if we find the path p from s to t will never contain edge $(v,s)$. Otherwise, since the path begins at s and reaches s again, thus the path is not a simple path which contradict the definition of augment path.

    Since $(v,s)$ will not appear in any augment path, the path finded each time by Ford-Fulkerson method remains the same. Thus this procedure is not influenced and can get the right result.

\end{pssolution}

\begin{psproblem}[26.2.10]
\end{psproblem}
\begin{pssolution}
    Suppose we already find a maximum flow f. Make a new network where the capacity of each edge equals the flow in the original network. We only need to prove by $|E|$ times of finding augment path we can find the maximum flow of the new network. Note that the maximum flow of the new network is full on every edge.

    Run Ford-Fulkerson algorithm on the new network with one modification that each time when we find the minimal capacity of n argument path, instead of update the corresponding edges in residual network, we just delete (u,v). This will not influence the augment path finded afterwards for it is has the same value of the final flow. Otherwise, there is an augment path contain (v,u), then another augment path containing (u,v) must exist. Denote the first as $s\rightarrow v\rightarrow u \rightarrow t$ and the other $s\rightarrow u\rightarrow v\rightarrow t$, we can construct two augment path $s\rightarrow t$ without traversing (u,v).

    Since each time an edge is deleted, at most $|E|$ augment path can be finded. In this way, we can get at most $|E|$ augment path to get a maximum flow.
\end{pssolution}

\begin{psproblem}[26.2.12]
\end{psproblem}
\begin{pssolution}
    Since there is an edge (v,s) that f(v,s)=1, we hope to find a circle from s to s that there is flow on each edge. Since the capacity is all integer, if $f(v,u)>0$ the n$f(v,u)\geq 1$. Thus by deleting 1 from flow on the edges of the circle, we get a flow f' with the same value of f. In a flow, each vertex is connected with s, so is v. Thus there exists circle $s\sim v\rightarrow s$.

    To get f', we use a deep-first search on f starting at s. When we find this circle, we can perform above operations to get the desired flow f'.
\end{pssolution}

\begin{psproblem}[26.2.13]
\end{psproblem}
\begin{pssolution}
    Add $\frac{1}{2|E|}$ to each edge to form a network G'. To prove the minimal cut in G' is the minimal cut in G with the minimal edges, we will prove that the minimal cut in G' has corresponding minimal value of cut in G and in all minimal value cut it contains the least edges.

    For the first thing, for all the cut $f(S,T)$ in G, the corresponding cut in G' is $f(S,T)+\frac{1}{2|E|}E(S,T)<f(S,T)+1$. Since the flow is integer, two flow with different value differ at least one. Thus the minimal one remains minimal.

    For the second thing, the minimal cut in G' has minimal value $f(S,T)+\frac{1}{2|E|}E(S,T)$. Since $f(S,T)$ is the same for minimal cut in G, thus E(S,T) is minimal, namely it contains minimal edges in G.
\end{pssolution}

\begin{psproblem}[26.3.3]
\end{psproblem}
\begin{pssolution}
    The augment path is a simple path. Since the graph is a partite graph, the path must go like: $s\rightarrow l_1\rightarrow r_1\dots \rightarrow t$. Thus the over bounding of the augment path is $2\min{(V(L),V(R))}+1$.
\end{pssolution}

\begin{psproblem}[26-1]
\end{psproblem}
\begin{pssolution}
    (a). We have already discussed this problem in 26.1.7. The graph G' has $2\abs{V}$ vertices and $\abs{E}+\abs{V}$ edges.

    (b) Construct the flow network as follows: add super source s and edges connecting s to all initial points. If two node in the grid can be reached from each other, then add edges in the flow. Finally, add edges from boundary point to super sink t. All edges and vertices have unit capacity. Then the problem converted to determine whether the maximum flow can be m.
    
    With the relabeled-to-front algorithm, the running time is $O(V^3)=O(n^6)$.
\end{pssolution}

\begin{psproblem}[26-2]
\end{psproblem}
\begin{pssolution}
    (a). Construct the network G' as suggested and let all edges have capacity 1. This network solves a matching problem and we have the partite graph $X,Y$. We can construct a path cover from a flow as follows:
    \begin{itemize}
        \item select a node not in path unless all nodes are in the path.
        \item if there is no $f(x_i,y_j)>0$, put it into path cover as a path of length.
        \item else put $x_i,y_j$ into path, and select $x_j$ back to step two.
    \end{itemize}

    In this way, we construct a path for each flow. Actually, each edge in the path cover is an edge in the matching of the partite network. Thus we get $|V|=\sum_{p\in P}{v(p)}=\sum_{p\in P}{1+e(p)}=P+\sum_{p}{e(p)}=|P|+M=|P|+|f|$(Note the graph is DAG). When $|f|$ is maximal, we will get the minimal path cover.

    (b) If the graph contains circle, the deduction $\sum_{p\in P}{v(p)}=\sum_{p\in P}{1+e(p)}$ will not hold.
    
    An counter-example is graph with $E=\{(1,2),(2,3),(3,4),(4,1)\}$. With our algorithm, $(x_1,y_2)$ and $(x_4,y_1)$ can be either chosen. If the latter one is chosed, then a false answer of 2 is generated. Thus in this case we cannot guarantee the algorithm to get the right answer.
\end{pssolution}


%%%%% Correction %%%%%
%section*{Correction}

%%%%% Feedback %%%%%
\section*{Part II: Feedback}
\begin{itemize}
    \item Question in 26.1.6
\end{itemize}
% 
\section{出题}
话说当今计算机业事，机器学习当道，码农趋之若鹜。遇投机者，以数月速学之，可夸夸其谈不下数个时辰，然则凡细节、原理之处，一概不通。更有好事者炒作概念，一时间机器学习从业者良莠不齐。为探求职者功力之深，G公司取一巧计，出题如下：

线性方程组：
\[
	\begin{cases}
		a_{11}x_1+a_{12}x_2+\cdots+a_{1n}x_n & =b_1 \\
		a_{21}x_1+a_{22}x_2+\cdots+a_{2n}x_n & =b_2 \\
		\cdots                               &      \\
		a_{m1}x_1+a_{m2}x_2+\cdots+a_{mn}x_n & =b_m \\
	\end{cases}
\]

可矩阵表示为 $\boldsymbol Ax=\boldsymbol b$ 其中 $\boldsymbol A\in \boldsymbol M_{m\times n}$，$x,\boldsymbol b\in \mathbb{R}^n$。在 $m>n$ 时，上述方程组多数情况下无解。但如果允许引入残差，将方程求解问题转化为最小化残差和问题，则仍可解。

引入定义在 $\mathbb{R}^n$ 上的函数：残差平方和 $S(x)=\left\|\boldsymbol Ax-\boldsymbol b\right\|^2$。由此给出最小二乘解的形式化定义：对于线性方程组 $\boldsymbol Ax=\boldsymbol b$，求最小二乘解 $\hat x=\mathrm{argmin}(S(x))$，即 使 $S(x)$ 最小的 $x$。 这便是经典的最小二乘问题。
\begin{problem}
试计算最小二乘解，并考虑提供算法的实际应用情况
\end{problem}

\begin{solution}
	这里只考虑较为简单的情况。首先设 $\boldsymbol A:$ 行向量分别为 $a_1,a_2,\cdots,a_m$，则有
	\begin{align*}
		S(x) & =\left\|\boldsymbol Ax-\boldsymbol b\right\|^2                      \\
		     & =\sum_{i=1}^m (a_ix-b_i)^2                  \\
		     & =\mathrm{tr}\left((\boldsymbol Ax-\boldsymbol b)^\top (\boldsymbol Ax-\boldsymbol b)\right)
	\end{align*}

	将求导运算推广到矩阵的迹上：下证对矩阵 $\boldsymbol A\in \boldsymbol M_{m\times n},\boldsymbol B\in \boldsymbol M_{n\times m}$，有

	$$\frac{\partial \mathrm{tr}(\boldsymbol A\boldsymbol B)}{\partial \boldsymbol A}=\frac{\partial \mathrm{tr}(\boldsymbol B\boldsymbol A)}{\partial \boldsymbol A}=\boldsymbol B^\top $$

	首先，有 $\mathrm{tr}(\boldsymbol A\boldsymbol B)=\sum_{i=1}^m \sum_{j=1}^n a_{ij}b_{ji}$，进而 $\frac{\partial \mathrm{tr}(\boldsymbol A\boldsymbol B)}{\partial a_{ij} }=b_{ji}$。从而可得上述命题。
	并且类似地，我们有
	$$\frac{\partial \mathrm{tr}(\boldsymbol A^\top \boldsymbol B)}{\partial a_{ij}}=\frac{\partial \mathrm{tr}(\boldsymbol B\boldsymbol A^\top )}{\partial a_{ij}}=\boldsymbol B$$

	进而

	\begin{align*}
		\frac{\partial S(x)}{\partial x} & =\frac{\partial\mathrm{tr}\left((\boldsymbol Ax-\boldsymbol b)^\top (\boldsymbol Ax-\boldsymbol b)\right)}{\partial x} \\
		                                 & =\frac{\partial\mathrm{tr}(x^\top \boldsymbol A^\top \boldsymbol Ax-\boldsymbol b^\top \boldsymbol Ax-\boldsymbol bx^\top \boldsymbol A^\top +\boldsymbol b^\top \boldsymbol b)}{\partial x} \\
		                                 & =2(\boldsymbol A^\top \boldsymbol Ax-\boldsymbol A^\top \boldsymbol b)
	\end{align*}

	令导数为 $0$，得到 $\boldsymbol A^\top \boldsymbol Ax-\boldsymbol A^\top \boldsymbol b=0$，$\boldsymbol A^\top \boldsymbol A$ 可逆时，有
	$$x=(\boldsymbol A^\top \boldsymbol A)^{-1}\boldsymbol A^\top \boldsymbol b$$
\end{solution}

\section{首试}
三人之中，最先面试者为贾有才。其方见此题，心中窃喜，遂写下MATLAB代码：
\begin{solution}
	\hspace{1cm}$x = inv(A'*A)*A'*b; \mathrm{//MATLAB\, codes}$
\end{solution}

贾抬头，却见面试官微微摇头。贾大恐，遂解释其中数学推导云云，皆与上同，不复再言。然面试官不以为意。贾面试之结果可知也。

贾之解法固然正确，何不得面试官之心耶？

可以看到，在这种直接计算最小二乘解的过程中，我们进行了3次矩阵乘法，和一次矩阵求逆。矩阵乘法利用$O(n^3)$的暴力做法已足够快，即使利用$O(n^{2.8})$算法却因常数巨大、实现复杂且无法降低整体复杂度而常不被采用。虽然在矩阵乘法的优化空间小，矩阵求逆却有比最简单的高斯消元法更好的选择。

回顾高斯消元求逆，假设矩阵为n阶方阵，一般我们先进行下三角消元。对于第$i$个主元，进行下三角消元时我们需要进行$(n-i)*(n+1-i)$次乘法和加法（消$n-i$行，每一行$n+1-i$个非0元素），所以将一次乘法和加法作为单位时间后，进行下三角消元消耗的时间为
\[\sum_{i=1}^{n}(n-i)*(n+1-i)=\sum_{i=0}^{n-1}i*(i+1)=\frac{(n+1)n(n-1)}{3}\]
而类似的，进行上三角消元时，我们要进行n次行变换。但由于此时每行只有一个非0元，所以总消耗为
\[\sum_{i=0}^{n-1}i=\frac{(n-1)*n}{2}\]
但是计算矩阵逆时同时要对$I_n$进行相同的行变换，由于操作完全对称，所以将最后时间乘二，我们就得到了高斯消元求逆的一个比较精确的时间相对值：
\[T(n)=(\frac{(n+1)n(n-1)}{3}+\frac{(n-1)*n}{2})*2=\frac{2n^3+3n^2-5n}{3}=\frac{2}{3}n^3+O(n^2)\]

而用$A^{-1}$求矩阵逆采用是最朴素的方法，其相对时间和上述计算相差无几。实际上存在更好的方法。

\section{突破}
其后面试者为甄汇。审题有饷，方落笔
\begin{solution}
	\hspace{1cm}$x=(A'*A)\backslash (A'*b); \mathrm{//MATLAB\, codes}$
\end{solution}

面试官微颔首。此句何意也？

在$MATLAB$中，$\backslash$是矩阵左除的意思。即左乘矩阵的逆。在早期的MATLAB版本中，左除判断矩阵是否能够$LU$分解，若可以则使用$LU$分解求逆。

该算法的思路为，将待求逆矩阵$A$分解成下三角矩阵$L$与$U$的乘积，再通过求解两个上（下）三角形方程组进行求解：

\[\boldsymbol L\boldsymbol U\boldsymbol x=\boldsymbol b\Rightarrow\begin{cases}\boldsymbol L\boldsymbol y&=\boldsymbol b \\ \boldsymbol U\boldsymbol x&=\boldsymbol y\end{cases}\]

对于下三角形方程组，易知有$O(n^2)$的计算方法$x_i=\frac{b_i-\sum_{i=1}^{m-1}l_{m,i}x_i}{l_{m,m}}$。若$LU$分解能在$O(n^3)$完成，则渐进时间复杂度与高斯消元相同。这看起来好像是画蛇添足，但实际不然。

首先有如下命题：
\begin{proposition}
	设方阵A的各阶顺序主子式都非0，则存在下三角矩阵$\boldsymbol L$与上三角矩阵$\boldsymbol U$使得$\boldsymbol A=\boldsymbol L*\boldsymbol U$。
\end{proposition}

为了证明这个命题，我们有如下引理

\begin{lemma}
	若$\boldsymbol A$与$\boldsymbol B$都是n阶上（下）三角方阵，则$\boldsymbol A*\boldsymbol B$也是n阶上（下）三角方阵。
\end{lemma}

此处略去证明，因此引理只需用矩阵乘法的定义验证即可，即验证$\boldsymbol A*\boldsymbol B$的下方（上方）元素都为0，而这是显然的。这样我们就可以给出命题的证明
\begin{proof}
	考虑对$\boldsymbol A$进行高斯消元。注意到一个事实：进行下三角消元时，不需要用到行交换。也就是说，每一行在向下消元后，下一行的主元只会向右移一位。假设第$i$次消元将下一行主元右移了两位以上，则考虑A的$i+1$顺序主子式，其行向量线性相关，为0，矛盾。（这里同时也意味着$\boldsymbol A_{1,1}\not=0$）

	由此知此高斯消元相当于左乘上三角初等阵$\boldsymbol P_i$，而$\boldsymbol P_i$为下三角初等阵。于是$\boldsymbol A=\boldsymbol P_i^{-1}*\boldsymbol P_i*\boldsymbol A=\boldsymbol P_i^{-1}*\boldsymbol A^\top$，此为消元过程。经过n次消元后，$\boldsymbol A$变为上三角阵，左乘了$n$个下三角初等阵，所以$\boldsymbol A=\boldsymbol P*\boldsymbol A^\top$，即$\boldsymbol A=\boldsymbol L*\boldsymbol U$。\hspace{3cm}证毕
\end{proof}

如何快速计算$\boldsymbol L$与$\boldsymbol U$呢？假设L的对角元素都为1（对应杜立特分解），下面给出一种递推算法：

首先由$\boldsymbol L*\boldsymbol U=\boldsymbol A$，设$\boldsymbol L=(l_{i,j}),\boldsymbol U=(u_{i,j}),\boldsymbol A=(a_{i,j})$，显然有$u_{1,j}=a_{1,j}$，从而$l_{i,1}=\frac{a_{i,1}}{u_{1,j}}$。得到这两个初始条件后，如下递推计算剩下元素：
\[u_{r,j}=\frac{a_{r,j}-\sum_{k=1}^{r-1}l_{i,k}u_{k,j}}{l_{r,r}}=a_{r,j}-\sum_{k=1}^{r-1}l_{i,k}*u_{k,j}\]
\[l_{i,r}=\frac{a_{i,r}-\sum_{k=1}^{r-1}l_{i,k}u_{k,r}}{u_{r,r}}\]
此法复杂度如何？注意到对于$l_{ir}$,需要最多$r$次乘除运算。$u_{rj}$计算的时间复杂度与$l_{ir}$相当。总时间复杂度为
\[2*\sum_{i=1}^{n}\sum^{r=1}_{i-1}r=2*\sum_{i=1}^{n}\frac{i(i-1)}{2}=\sum_{i=1}^{n}i^2+O(n^2)=\frac{1}{3}n^3+O(n^2)\]

再者是是两个三角形矩阵对应的方程求解，如上分析，时间为$O(n^2)$将所有步骤复杂度相加，可以发现相对于普通的高斯消元而言，其最高项系数确实减少了：
\[T(n)=\frac{1}{3}n^3+O(n^2)\]

在实际应用中，$LU$分解的高处甚于此，更有
\begin{itemize}
	\item 在MATLAB代码的具体实现中，$LU$分解（可以分解的情况下）解方程组的理论时间复杂度系数是高斯消元的$\frac{1}{3}$。当$n=1000$时，平均每次高斯消元耗时$30s$, $LU$分解法耗时$22s$。（机器性能类于学校机房机）
	\item 节省内存开销。根据$LU$分解时的递推方法，我们不需要同时保存所有信息在内存中，求出$u_{rj}$后$a_{rj},j\geq r$的储存位置不再需要，求出$l_{ir}$后$a_{ir},i\geq r+1$的储存位置不再需要。对于超大矩阵，这可以节省大量内存空间。
	\item 对于求解$\boldsymbol Ax=\boldsymbol b$其中$\boldsymbol A$固定，$b$取多个值的问题，同过一次$O(n^3)$的$LU$分解，每个$b$只需$O(n^2)$的三角矩阵求解即可。当$b$的数量$m$足够大时，时间复杂度为$O(mn^2)$优于$O(mn^3)$。
\end{itemize}

\section{真神技也！}
最后登场者为超勇也。其人气宇轩昂，仪表不凡，洋洋洒洒落笔写道：
\begin{solution}
	$\mathrm{//MATLAB\,code}$
	\begin{center}
		$[Q, R] = qr(A);$ \\ $x=R\backslash(Q'*b);$
	\end{center}
\end{solution}

面试官见此代码，有面露疑色而待解释者，有舒容展眉而称赞者。何以如此也？此非习线代者所不能得之方法也。可见此人定非碌碌之辈，或将效大力于公司。

这首行代码，用的是$QR$分解，又称正交三角分解。它是把一个实可逆矩阵$\boldsymbol A$化成正交矩阵$\boldsymbol Q$和实非奇异上三角矩阵$\boldsymbol R$的乘积，$\boldsymbol A=\boldsymbol Q\boldsymbol R$。对于稀疏矩阵，可以使用Givens方法，对于一般矩阵，可以使用Householder方法求解。对于$n$阶矩阵，其时间复杂度为$O(n^3)$。篇幅有限，这里不展开讨论，有兴趣者可自行搜索相关资料。

对于正交矩阵$\boldsymbol Q$，我们有$\boldsymbol Q^\top =\boldsymbol Q^{-1}$，故问题可转换为求$\boldsymbol Rx=\boldsymbol Q^\top \boldsymbol b$。又$\boldsymbol R$为上三角阵，故只需$O(n^2)$即可得到解。对于最小二乘问题，由$\boldsymbol Q^\top \boldsymbol Q=1$，我们更有惊喜的发现：
\[\boldsymbol A^\top \boldsymbol Ax=\boldsymbol A^\top \boldsymbol b\Rightarrow (\boldsymbol Q\boldsymbol R)^\top \boldsymbol Q\boldsymbol R=(\boldsymbol Q\boldsymbol R)^\top \boldsymbol b\Rightarrow \boldsymbol R^\top \boldsymbol Rx=\boldsymbol R^\top \boldsymbol Q^\top \boldsymbol b\Rightarrow \boldsymbol Rx=\boldsymbol Q^\top \boldsymbol b\]

如此一来，我们首先绕过了矩阵的乘法。为了了解其进一步的优点，我们引入一个新的概念：条件数
\begin{definition}
	对任意一个矩阵 $\boldsymbol A$ ，我们可以对它进行SVD分解 $\boldsymbol A = \boldsymbol U \boldsymbol \Sigma \boldsymbol  V^\top$ ，其中奇异值 $\sigma_1, \sigma_2, \ldots$  是 $\boldsymbol \Sigma$ 矩阵的对角线上的元素。定义矩阵 $\boldsymbol A$ 的条件数为
	\[\mathrm{cond}\, \boldsymbol A = \frac{\max_i\sigma_i}{\min_i \sigma_i} \,. \]
\end{definition}

这里我们不给出条件数的推导和性质，只介绍其直观含义。当$\mathrm{cond}\, \boldsymbol A$越大时，计算的误差也就越大。想象高斯消元中，一个超大值(e.g. $10^{16}$)除以一个小值(e.g. $1^{-16}$)将会导致数值溢出，反之则会导致精度丢失。故保证计算时有较小的条件数是重要的。

对于前两种方法，都计算了$A^\top A$，这导致$\mathrm{cond}\, \boldsymbol A^\top \boldsymbol A = \mathrm{cond}\boldsymbol V\boldsymbol\Sigma^2\boldsymbol V^\top = (\mathrm{cond} \boldsymbol A)^2$。 由此一来，条件数被平方了。

而超氏的方法中，没有计算改矩阵乘法，避免了条件数被平方。配合以数值稳定的$Household$方法计算$QR$分解，将得到数值稳定的答案。

\section{面试之外：MATLAB的反除函数}
事实上，作为一款专业（付费）软件，MATLAB的函数一直在完善升级。当前的反除$\backslash$对大型矩阵运算时过程相当复杂。实际上，反除是使用一系列判断确定矩阵的特有性质（类型）选择最优的分解方法后进行求解。具体的判断如图所示。

简单的概括，反除法下，MATLAB先判断矩阵是否为方阵，否则用$QR$分解，是则判断是否是三角阵或交换行的位置后构成三角阵。再判断是否（共轭）对称，对于对称阵，若可以三角分解，则使用Cholesky分解，其计算量约为$LU$分解的一半，否则使用LDL分解。若不对称，当矩阵是Hessenberg矩阵时（若 $a_{ij}=0， i>j+1$，則 $\boldsymbol A$ 称为上 Hessenberg 矩阵），可以针对其使用Hessenberg分解加速，否则才使用我们上述的$LU$分解。

注意到对于$n$阶方阵，MATLAB并不会去选择$QR$分解，而对于最小二乘问题，$QR$分解有其特殊的好处。故超勇的方法确实比甄汇之方法优。

而在MATLAB R2017b中，通过引入函数decomposition再次对反除法进行了优化。使用$\mathbb{tA=decomposition(A);\,x=tA\backslash b}$，可以先对$\boldsymbol A$进行合适的分解，再进行计算。这可以方便的实现前面提到的$LU$分解优点中的第三点。实践检验，对于$\boldsymbol A$不变$\boldsymbol b$改变的情况下，仅$1000$次运算$n=100$的矩阵，普通的反除效率和先使用decomposition的反除效率相乘达9倍。

希望以上内容也能帮助你更好的使用MATLAB的函数。


\section{结语}
无名氏评曰：当今世事，群雄逐职，一位难求。非有扎实学问，超人之能，难以致臻高手之境，一招一回间，马脚尽露。数学者，计算机之基础也；线代、矩阵者，机器学习之基石也。仅一最小二乘问题，三法共解，高下立断。从计算机者，因以此为鉴，既应通过学习好的数学工具提升算法，也应不满足于仅于数学中得解，而因考虑何以快速的计算答案。此双剑合璧，可冠一方也。不学可乎？不精学可乎？

\section*{参考资料}
\begin{itemize}
    \item[1] 方保镕，周继东，李医民 《矩阵论》
    \item[2] 张晧 https://www.zhihu.com/question/62482926/answer/526304253
    \item[3] 菡父 https://zhuanlan.zhihu.com/p/30958676
    \item[4] 线代启示录 https://ccjou.wordpress.com
\end{itemize}

% \chapter{带 English 的标题}
\label{cha:intro}

这是 NJUrepo 的示例文档，基本上覆盖了模板中所有格式的设置。建议大家在使用模板之前，除了阅读《NJUrepo用户手册》，这个示例文档也最好能看一看。

小老鼠偷吃热凉粉；短长虫环绕矮高粱\footnote{韩愈（768-824），字退之，河南河阳（
  今河南孟县）人，自称郡望昌黎，世称韩昌黎。幼孤贫刻苦好学，德宗贞元八年进士。曾
  任监察御史，因上疏请免关中赋役，贬为阳山县令。后随宰相裴度平定淮西迁刑部侍郎，
  又因上表谏迎佛骨，贬潮州刺史。做过吏部侍郎，死谥文公，故世称韩吏部、韩文公。是
  唐代古文运动领袖，与柳宗元合称韩柳。诗力求险怪新奇，雄浑重气势。}。


\section{封面相关}
封面的例子请参看 \texttt{cover.tex}。主要符号表参看 \texttt{denotation.tex}，附录和
个人简历分别参看 \texttt{appendix01.tex} 和 \texttt{resume.tex}。里面的命令都很直
观，一看即会\footnote{你说还是看不懂？怎么会呢？}。

\section{字体命令}
\label{sec:first}

苏轼（1037-1101），北宋文学家、书画家。字子瞻，号东坡居士，眉州眉山（今属四川）人
。苏洵子。嘉佑进士。神宗时曾任祠部员外郎，因反对王安石新法而求外职，任杭州通判，
知密州、徐州、湖州。后以作诗“谤讪朝廷”罪贬黄州。哲宗时任翰林学士，曾出知杭州、
颖州等，官至礼部尚书。后又贬谪惠州、儋州。北还后第二年病死常州。南宋时追谥文忠。
与父洵弟辙，合称“三苏”。在政治上属于旧党，但也有改革弊政的要求。其文汪洋恣肆，
明白畅达，为“唐宋八大家”之一。  其诗清新豪健，善用夸张比喻，在艺术表现方面独具
风格。少数诗篇也能反映民间疾苦，指责统治者的奢侈骄纵。词开豪放一派，对后代很有影
响。《念奴娇·赤壁怀古》、《水调歌头·丙辰中秋》传诵甚广。

{\kaishu 坡仙擅长行书、楷书，取法李邕、徐浩、颜真卿、杨凝式，而能自创新意。用笔丰腴
  跌宕，有天真烂漫之趣。与蔡襄、黄庭坚、米芾并称“宋四家”。能画竹，学文同，也喜
  作枯木怪石。论画主张“神似”，认为“论画以形似，见与儿童邻”；高度评价“诗中有
  画，画中有诗”的艺术造诣。诗文有《东坡七集》等。存世书迹有《答谢民师论文帖》、
  《祭黄几道文》、《前赤壁赋》、《黄州寒食诗帖》等。  画迹有《枯木怪石图》、《
  竹石图》等。}

{\fangsong 易与天地准，故能弥纶天地之道。仰以观於天文，俯以察於地理，是故知幽明之故。原
  始反终，故知死生之说。精气为物，游魂为变，是故知鬼神之情状。与天地相似，故不违。
  知周乎万物，而道济天下，故不过。旁行而不流，乐天知命，故不忧。安土敦乎仁，故
  能爱。范围天地之化而不过，曲成万物而不遗，通乎昼夜之道而知，故神无方而易无体。}

% 非本科生一般用不到幼圆与隶书字体。需要的同学请查看 ctex 文档。
{\ifcsname youyuan\endcsname\youyuan\else[无 \cs{youyuan} 字体。]\fi 有天地，然后
  万物生焉。盈天地之间者，唯万物，故受之以屯；屯者盈也，屯者物之始生也。物生必蒙，
  故受之以蒙；蒙者蒙也，物之穉也。物穉不可不养也，故受之以需；需者饮食之道也。饮
  食必有讼，故受之以讼。讼必有众起，故受之以师；师者众也。众必有所比，故受之以比；
  比者比也。比必有所畜也，故受之以小畜。物畜然后有礼，故受之以履。}

{\heiti 履而泰，然后安，故受之以泰；泰者通也。物不可以终通，故受之以否。物不可以终
  否，故受之以同人。与人同者，物必归焉，故受之以大有。有大者不可以盈，故受之以谦。
  有大而能谦，必豫，故受之以豫。豫必有随，故受之以随。以喜随人者，必有事，故受
  之以蛊；蛊者事也。}

{\ifcsname lishu\endcsname\lishu\else[无 \cs{lishu} 字体。]\fi 有事而后可大，故受
  之以临；临者大也。物大然后可观，故受之以观。可观而后有所合，故受之以噬嗑；嗑者
  合也。物不可以苟合而已，故受之以贲；贲者饰也。致饰然后亨，则尽矣，故受之以剥；
  剥者剥也。物不可以终尽，剥穷上反下，故受之以复。复则不妄矣，故受之以无妄。}

{\songti 有无妄然后可畜，故受之以大畜。物畜然后可养，故受之以颐；颐者养也。不养则不
  可动，故受之以大过。物不可以终过，故受之以坎；坎者陷也。陷必有所丽，故受之以
  离；离者丽也。}

\section{表格样本}
\label{chap1:sample:table}

\subsection{基本表格}
\label{sec:basictable}

模板中关于表格的宏包有三个：\pkg{booktabs}、\pkg{array} 和 \pkg{longtabular}，命
令有一个 \cs{hlinewd}。三线表可以用 \pkg{booktabs} 提供
的 \cs{toprule}、\cs{midrule} 和 \cs{bottomrule}。它们与 \pkg{longtable} 能很好的
配合使用。如果表格比较简单的话可以直接用命令 \cs{hlinewd}\marg{width} 控制。
\begin{table}[htb]
  \centering
  \begin{minipage}[t]{0.8\linewidth} % 如果想在表格中使用脚注，minipage是个不错的办法
  \caption[模板文件]{模板文件。如果表格的标题很长，那么在表格索引中就会很不美
    观，所以要像 chapter 那样在前面用中括号写一个简短的标题。这个标题会出现在索
    引中。}
  \label{tab:template-files}
    \begin{tabularx}{\linewidth}{lX}
      \toprule[1.5pt]
      {\heiti 文件名} & {\heiti 描述} \\\midrule[1pt]
      thuthesis.ins & \LaTeX{} 安装文件，\textsc{DocStrip}\footnote{表格中的脚注} \\
      thuthesis.dtx & 所有的一切都在这里面\footnote{再来一个}。\\
      thuthesis.cls & 模板类文件。\\
      thuthesis.cfg & 模板配置文。cls 和 cfg 由前两个文件生成。\\
      thuthesis-numeric.bst    & 参考文献 BIB\TeX\ 样式文件。\\
      thuthesis-author-year.bst    & 参考文献 BIB\TeX\ 样式文件。\\
      thuthesis.sty   & 常用的包和命令写在这里，减轻主文件的负担。\\
      \bottomrule[1.5pt]
    \end{tabularx}
  \end{minipage}
\end{table}

首先来看一个最简单的表格。表 \ref{tab:template-files} 列举了本模板主要文件及其功
能。请大家注意三线表中各条线对应的命令。这个例子还展示了如何在表格中正确使用脚注。
由于 \LaTeX{} 本身不支持在表格中使用 \cs{footnote}，所以我们不得不将表格放在
小页中，而且最好将表格的宽度设置为小页的宽度，这样脚注看起来才更美观。

\subsection{复杂表格}
\label{sec:complicatedtable}

我们经常会在表格下方标注数据来源，或者对表格里面的条目进行解释。前面的脚注是一种
不错的方法，如果不喜欢脚注，可以在表格后面写注释，比如表~\ref{tab:tabexamp1}。
\begin{table}[htbp]
  \centering
  \caption{复杂表格示例 1。这个引用 \cite{tex} 不会导致编号混乱。}
  \label{tab:tabexamp1}
  \begin{minipage}[t]{0.8\textwidth}
    \begin{tabularx}{\linewidth}{|l|X|X|X|X|}
      \hline
      \multirow{2}*{\diagbox[width=5em]{x}{y}} & \multicolumn{2}{c|}{First Half} & \multicolumn{2}{c|}{Second Half}\\\cline{2-5}
      & 1st Qtr &2nd Qtr&3rd Qtr&4th Qtr \\ \hline
      East$^{*}$ &   20.4&   27.4&   90&     20.4 \\
      West$^{**}$ &   30.6 &   38.6 &   34.6 &  31.6 \\ \hline
    \end{tabularx}\\[2pt]
    \footnotesize 注：数据来源《NJUrepo{} 使用手册》。\\
    *：东部\\
    **：西部
  \end{minipage}
\end{table}

此外，表~\ref{tab:tabexamp1} 同时还演示了另外两个功能：1）通过 \pkg{tabularx} 的
 \texttt{|X|} 扩展实现表格自动放大；2）通过命令 \cs{diagbox} 在表头部分
插入反斜线。

为了使我们的例子更接近实际情况，我会在必要的时候插入一些“无关”文字，以免太多图
表同时出现，导致排版效果不太理想。第一个出场的当然是我的最爱：风流潇洒、骏马绝尘、
健笔凌云的{\heiti 李太白}了。

李白，字太白，陇西成纪人。凉武昭王暠九世孙。或曰山东人，或曰蜀人。白少有逸才，志
气宏放，飘然有超世之心。初隐岷山，益州长史苏颋见而异之，曰：“是子天才英特，可比
相如。”天宝初，至长安，往见贺知章。知章见其文，叹曰：“子谪仙人也。”言于明皇，
召见金銮殿，奏颂一篇。帝赐食，亲为调羹，有诏供奉翰林。白犹与酒徒饮于市，帝坐沉香
亭子，意有所感，欲得白为乐章，召入，而白已醉。左右以水颒面，稍解，援笔成文，婉丽
精切。帝爱其才，数宴见。白常侍帝，醉，使高力士脱靴。力士素贵，耻之，摘其诗以激杨
贵妃。帝欲官白，妃辄沮止。白自知不为亲近所容，恳求还山。帝赐金放还。乃浪迹江湖，
终日沉饮。永王璘都督江陵，辟为僚佐。璘谋乱，兵败，白坐长流夜郎，会赦得还。族人阳
冰为当涂令，白往依之。代宗立，以左拾遗召，而白已卒。文宗时，诏以白歌诗、裴旻剑舞、
张旭草书为三绝云。集三十卷。今编诗二十五卷。\hfill —— 《全唐诗》诗人小传

浮动体的并排放置一般有两种情况：1）二者没有关系，为两个独立的浮动体；2）二者隶属
于同一个浮动体。对表格来说并排表格既可以像图~\ref{tab:parallel1}、
图~\ref{tab:parallel2} 使用小页环境，也可以如图~\ref{tab:subtable} 使用子表格来做。
图的例子参见第~\ref{sec:multifig} 节。

\begin{table}[htbp]
\noindent\begin{minipage}{0.5\textwidth}
\centering
\caption{第一个并排子表格}
\label{tab:parallel1}
\begin{tabular}{p{2cm}p{2cm}}
\toprule[1.5pt]
111 & 222 \\\midrule[1pt]
222 & 333 \\\bottomrule[1.5pt]
\end{tabular}
\end{minipage}%
\begin{minipage}{0.5\textwidth}
\centering
\caption{第二个并排子表格}
\label{tab:parallel2}
\begin{tabular}{p{2cm}p{2cm}}
\toprule[1.5pt]
111 & 222 \\\midrule[1pt]
222 & 333 \\\bottomrule[1.5pt]
\end{tabular}
\end{minipage}
\end{table}

然后就是忧国忧民，诗家楷模杜工部了。杜甫，字子美，其先襄阳人，曾祖依艺为巩令，因
居巩。甫天宝初应进士，不第。后献《三大礼赋》，明皇奇之，召试文章，授京兆府兵曹参
军。安禄山陷京师，肃宗即位灵武，甫自贼中遁赴行在，拜左拾遗。以论救房琯，出为华州
司功参军。关辅饥乱，寓居同州同谷县，身自负薪采梠，餔糒不给。久之，召补京兆府功曹，
道阻不赴。严武镇成都，奏为参谋、检校工部员外郎，赐绯。武与甫世旧，待遇甚厚。乃于
成都浣花里种竹植树，枕江结庐，纵酒啸歌其中。武卒，甫无所依，乃之东蜀就高適。既至
而適卒。是岁，蜀帅相攻杀，蜀大扰。甫携家避乱荆楚，扁舟下峡，未维舟而江陵亦乱。乃
溯沿湘流，游衡山，寓居耒阳。卒年五十九。元和中，归葬偃师首阳山，元稹志其墓。天宝
间，甫与李白齐名，时称李杜。然元稹之言曰：“李白壮浪纵恣，摆去拘束，诚亦差肩子美
矣。至若铺陈终始，排比声韵，大或千言，次犹数百，词气豪迈，而风调清深，属对律切，
而脱弃凡近，则李尚不能历其藩翰，况堂奥乎。”白居易亦云：“杜诗贯穿古今，  尽工尽
善，殆过于李。”元、白之论如此。盖其出处劳佚，喜乐悲愤，好贤恶恶，一见之于诗。而
又以忠君忧国、伤时念乱为本旨。读其诗可以知其世，故当时谓之“诗史”。旧集诗文共六
十卷，今编诗十九卷。

\begin{table}[htbp]
\centering
\caption{并排子表格}
\label{tab:subtable}
\subcaptionbox{第一个子表格}
{
\begin{tabular}{p{2cm}p{2cm}}
\toprule[1.5pt]
111 & 222 \\\midrule[1pt]
222 & 333 \\\bottomrule[1.5pt]
\end{tabular}
}
\hskip2cm
\subcaptionbox{第二个子表格}
{
\begin{tabular}{p{2cm}p{2cm}}
\toprule[1.5pt]
111 & 222 \\\midrule[1pt]
222 & 333 \\\bottomrule[1.5pt]
\end{tabular}
}
\end{table}

不可否认 \LaTeX{} 的表格功能没有想象中的那么强大，不过只要足够认真，足够细致，
同样可以排出来非常复杂非常漂亮的表格。请参看表~\ref{tab:tabexamp2}。
\begin{table}[htbp]
  \centering\dawu[1.3]
  \caption{复杂表格示例 2}
  \label{tab:tabexamp2}
  \begin{tabular}[c]{|m{1.5cm}|c|c|c|c|c|c|}\hline
    \multicolumn{2}{|c|}{Network Topology} & \# of nodes &
    \multicolumn{3}{c|}{\# of clients} & Server \\\hline
    GT-ITM & Waxman Transit-Stub & 600 &
    \multirow{2}{1em}{2\%}&
    \multirow{2}{1.5em}{10\%}&
    \multirow{2}{1.5em}{50\%}&
    \multirow{2}{1.2in}{Max. Connectivity}\\\cline{1-3}
    \multicolumn{2}{|c|}{Inet-2.1} & 6000 & & & &\\\hline
    \multirow{2}{1.5cm}{Xue} & Rui  & Ni &\multicolumn{4}{c|}{\multirow{2}*{NJUrepo}}\\\cline{2-3}
    & \multicolumn{2}{c|}{ABCDEF} &\multicolumn{4}{c|}{} \\\hline
\end{tabular}
\end{table}

最后就是清新飘逸、文约意赅、空谷绝响的王大侠了。王维，字摩诘，河东人。工书画，与
弟缙俱有俊才。开元九年，进士擢第，调太乐丞。坐累为济州司仓参军，历右拾遗、监察御
史、左补阙、库部郎中，拜吏部郎中。天宝末，为给事中。安禄山陷两都，维为贼所得，服
药阳喑，拘于菩提寺。禄山宴凝碧池，维潜赋诗悲悼，闻于行在。贼平，陷贼官三等定罪，
特原之，责授太子中允，迁中庶子、中书舍人。复拜给事中，转尚书右丞。维以诗名盛于开
元、天宝间，宁薛诸王驸马豪贵之门，无不拂席迎之。得宋之问辋川别墅，山水绝胜，与道
友裴迪，浮舟往来，弹琴赋诗，啸咏终日。笃于奉佛，晚年长斋禅诵。一日，忽索笔作书
数纸，别弟缙及平生亲故，舍笔而卒。赠秘书监。宝应中，代宗问缙：“朕常于诸王坐闻维
乐章，今存几何？”缙集诗六卷，文四卷，表上之。敕答云，卿伯氏位列先朝，名高希代。
抗行周雅，长揖楚辞。诗家者流，时论归美。克成编录，叹息良深。殷璠谓维诗词秀调雅，
意新理惬。在泉成珠，著壁成绘。苏轼亦云：“维诗中有画，画中有诗也。”今编诗四卷。

要想用好论文模板还是得提前学习一些 \TeX/\LaTeX{}的相关知识，具备一些基本能力，掌
握一些常见技巧，否则一旦遇到问题还真是比较麻烦。我们见过很多这样的同学，一直以来
都是使用 Word 等字处理工具，以为 \LaTeX{}模板的用法也应该类似，所以就沿袭同样的思
路来对待这种所见非所得的排版工具，结果被折腾的焦头烂额，疲惫不堪。

如果您要排版的表格长度超过一页，那么推荐使用 \pkg{longtable} 或者 \pkg{supertabular}
宏包，模板对 \pkg{longtable} 进行了相应的设置，所以用起来可能简单一些。
表~\ref{tab:performance} 就是 \pkg{longtable} 的简单示例。
\begin{longtable}[c]{c*{6}{r}}
\caption{实验数据}\label{tab:performance}\\
\toprule[1.5pt]
 测试程序 & \multicolumn{1}{c}{正常运行} & \multicolumn{1}{c}{同步} & \multicolumn{1}{c}{检查点} & \multicolumn{1}{c}{卷回恢复}
& \multicolumn{1}{c}{进程迁移} & \multicolumn{1}{c}{检查点} \\
& \multicolumn{1}{c}{时间 (s)}& \multicolumn{1}{c}{时间 (s)}&
\multicolumn{1}{c}{时间 (s)}& \multicolumn{1}{c}{时间 (s)}& \multicolumn{1}{c}{
  时间 (s)}&  文件（KB）\\\midrule[1pt]
\endfirsthead
\multicolumn{7}{c}{续表~\thetable\hskip1em 实验数据}\\
\toprule[1.5pt]
 测试程序 & \multicolumn{1}{c}{正常运行} & \multicolumn{1}{c}{同步} & \multicolumn{1}{c}{检查点} & \multicolumn{1}{c}{卷回恢复}
& \multicolumn{1}{c}{进程迁移} & \multicolumn{1}{c}{检查点} \\
& \multicolumn{1}{c}{时间 (s)}& \multicolumn{1}{c}{时间 (s)}&
\multicolumn{1}{c}{时间 (s)}& \multicolumn{1}{c}{时间 (s)}& \multicolumn{1}{c}{
  时间 (s)}&  文件（KB）\\\midrule[1pt]
\endhead
\hline
\multicolumn{7}{r}{续下页}
\endfoot
\endlastfoot
CG.A.2 & 23.05 & 0.002 & 0.116 & 0.035 & 0.589 & 32491 \\
CG.A.4 & 15.06 & 0.003 & 0.067 & 0.021 & 0.351 & 18211 \\
CG.A.8 & 13.38 & 0.004 & 0.072 & 0.023 & 0.210 & 9890 \\
CG.B.2 & 867.45 & 0.002 & 0.864 & 0.232 & 3.256 & 228562 \\
CG.B.4 & 501.61 & 0.003 & 0.438 & 0.136 & 2.075 & 123862 \\
CG.B.8 & 384.65 & 0.004 & 0.457 & 0.108 & 1.235 & 63777 \\
MG.A.2 & 112.27 & 0.002 & 0.846 & 0.237 & 3.930 & 236473 \\
MG.A.4 & 59.84 & 0.003 & 0.442 & 0.128 & 2.070 & 123875 \\
MG.A.8 & 31.38 & 0.003 & 0.476 & 0.114 & 1.041 & 60627 \\
MG.B.2 & 526.28 & 0.002 & 0.821 & 0.238 & 4.176 & 236635 \\
MG.B.4 & 280.11 & 0.003 & 0.432 & 0.130 & 1.706 & 123793 \\
MG.B.8 & 148.29 & 0.003 & 0.442 & 0.116 & 0.893 & 60600 \\
LU.A.2 & 2116.54 & 0.002 & 0.110 & 0.030 & 0.532 & 28754 \\
LU.A.4 & 1102.50 & 0.002 & 0.069 & 0.017 & 0.255 & 14915 \\
LU.A.8 & 574.47 & 0.003 & 0.067 & 0.016 & 0.192 & 8655 \\
LU.B.2 & 9712.87 & 0.002 & 0.357 & 0.104 & 1.734 & 101975 \\
LU.B.4 & 4757.80 & 0.003 & 0.190 & 0.056 & 0.808 & 53522 \\
LU.B.8 & 2444.05 & 0.004 & 0.222 & 0.057 & 0.548 & 30134 \\
EP.A.2 & 123.81 & 0.002 & 0.010 & 0.003 & 0.074 & 1834 \\
EP.A.4 & 61.92 & 0.003 & 0.011 & 0.004 & 0.073 & 1743 \\
EP.A.8 & 31.06 & 0.004 & 0.017 & 0.005 & 0.073 & 1661 \\
EP.B.2 & 495.49 & 0.001 & 0.009 & 0.003 & 0.196 & 2011 \\
EP.B.4 & 247.69 & 0.002 & 0.012 & 0.004 & 0.122 & 1663 \\
EP.B.8 & 126.74 & 0.003 & 0.017 & 0.005 & 0.083 & 1656 \\
\bottomrule[1.5pt]
\end{longtable}

\subsection{其它}
\label{sec:tableother}
如果不想让某个表格或者图片出现在索引里面，请使用命令 \cs{caption*}。
这个命令不会给表格编号，也就是出来的只有标题文字而没有“表~XX”，“图~XX”，否则
索引里面序号不连续就显得不伦不类，这也是 \LaTeX{} 里星号命令默认的规则。

有这种需求的多是本科同学的英文资料翻译部分，如果觉得附录中英文原文中的表格和图
片显示成“表”和“图”不协调的话，一个很好的办法就是用 \cs{caption*}，参数
随便自己写，比如不守规矩的表~1.111 和图~1.111 能满足这种特殊需要（可以参看附录部
分）。
\begin{table}[ht]
  \begin{minipage}{0.4\linewidth}
    \centering
    \caption*{表~1.111\quad 这是一个手动编号，不出现在索引中的表格。}
    \label{tab:badtabular}
      \framebox(150,50)[c]{NJUrepo}
  \end{minipage}%
  \hfill%
  \begin{minipage}{0.4\linewidth}
    \centering
    \caption*{Figure~1.111\quad 这是一个手动编号，不出现在索引中的图。}
    \label{tab:badfigure}
    \framebox(150,50)[c]{薛瑞尼}
  \end{minipage}
\end{table}

如果的确想让它编号，但又不想让它出现在索引中的话，目前模板上不支持。

最后，虽然大家不一定会独立使用小页，但是关于小页中的脚注还是有必要提一下。请看下
面的例子。

\begin{minipage}[t]{\linewidth-2\parindent}
  柳宗元，字子厚（773-819），河东（今永济县）人\footnote{山西永济水饺。}，是唐代
  杰出的文学家，哲学家，同时也是一位政治改革家。与韩愈共同倡导唐代古文运动，并称
  韩柳\footnote{唐宋八大家之首二位。}。
\end{minipage}

唐朝安史之乱后，宦官专权，藩镇割据，土地兼并日渐严重，社会生产破坏严重，民不聊生。柳宗
元对这种社会现实极为不满，他积极参加了王叔文领导的“永济革新”，并成为这一
运动的中坚人物。他们革除弊政，打击权奸，触犯了宦官和官僚贵族利益，在他们的联合反
扑下，改革失败了，柳宗元被贬为永州司马。

\section{定理环境}
\label{sec:theorem}

给大家演示一下各种和证明有关的环境：

\begin{assumption}
待月西厢下，迎风户半开；隔墙花影动，疑是玉人来。
\begin{align}
  \label{eq:eqnxmp}
  c & = a^2 - b^2 \\
    & = (a+b)(a-b)
\end{align}
\end{assumption}

千辛万苦，历尽艰难，得有今日。然相从数千里，未曾哀戚。今将渡江，方图百年欢笑，如
何反起悲伤？（引自《杜十娘怒沉百宝箱》）

\begin{definition}
子曰：「道千乘之国，敬事而信，节用而爱人，使民以时。」
\end{definition}

千古第一定义！问世间、情为何物，只教生死相许？天南地北双飞客，老翅几回寒暑。欢乐趣，离别苦，就中更有痴儿女。
君应有语，渺万里层云，千山暮雪，只影向谁去？

横汾路，寂寞当年箫鼓，荒烟依旧平楚。招魂楚些何嗟及，山鬼暗谛风雨。天也妒，未信与，莺儿燕子俱黄土。
千秋万古，为留待骚人，狂歌痛饮，来访雁丘处。

\begin{proposition}
 曾子曰：「吾日三省吾身 —— 为人谋而不忠乎？与朋友交而不信乎？传不习乎？」
\end{proposition}

多么凄美的命题啊！其日牛马嘶，新妇入青庐，奄奄黄昏后，寂寂人定初，我命绝今日，
魂去尸长留，揽裙脱丝履，举身赴清池，府吏闻此事，心知长别离，徘徊庭树下，自挂东南
枝。

\begin{remark}
天不言自高，水不言自流。
\begin{gather*}
\begin{split}
\varphi(x,z)
&=z-\gamma_{10}x-\gamma_{mn}x^mz^n\\
&=z-Mr^{-1}x-Mr^{-(m+n)}x^mz^n
\end{split}\\[6pt]
\begin{align} \zeta^0&=(\xi^0)^2,\\
\zeta^1 &=\xi^0\xi^1,\\
\zeta^2 &=(\xi^1)^2,
\end{align}
\end{gather*}
\end{remark}

天尊地卑，乾坤定矣。卑高以陈，贵贱位矣。 动静有常，刚柔断矣。方以类聚，物以群分，
吉凶生矣。在天成象，在地成形，变化见矣。鼓之以雷霆，润之以风雨，日月运行，一寒一
暑，乾道成男，坤道成女。乾知大始，坤作成物。乾以易知，坤以简能。易则易知，简则易
从。易知则有亲，易从则有功。有亲则可久，有功则可大。可久则贤人之德，可大则贤人之
业。易简，而天下矣之理矣；天下之理得，而成位乎其中矣。

\begin{axiom}
两点间直线段距离最短。
\begin{align}
x&\equiv y+1\pmod{m^2}\\
x&\equiv y+1\mod{m^2}\\
x&\equiv y+1\pod{m^2}
\end{align}
\end{axiom}

《彖曰》：大哉乾元，万物资始，乃统天。云行雨施，品物流形。大明始终，六位时成，时
乘六龙以御天。乾道变化，各正性命，保合大和，乃利贞。首出庶物，万国咸宁。

《象曰》：天行健，君子以自强不息。潜龙勿用，阳在下也。见龙再田，德施普也。终日乾
乾，反复道也。或跃在渊，进无咎也。飞龙在天，大人造也。亢龙有悔，盈不可久也。用九，
天德不可为首也。 　　

\newcommand\dif{\mathop{}\!\mathrm{d}}
\begin{lemma}
《猫和老鼠》是我最爱看的动画片。
\begin{multline*}%\tag*{[a]} % 这个不出现在索引中
\int_a^b\biggl\{\int_a^b[f(x)^2g(y)^2+f(y)^2g(x)^2]
 -2f(x)g(x)f(y)g(y)\dif x\biggr\}\dif y \\
 =\int_a^b\biggl\{g(y)^2\int_a^bf^2+f(y)^2
  \int_a^b g^2-2f(y)g(y)\int_a^b fg\biggr\}\dif y
\end{multline*}
\end{lemma}

行行重行行，与君生别离。相去万余里，各在天一涯。道路阻且长，会面安可知。胡马依北
风，越鸟巢南枝。相去日已远，衣带日已缓。浮云蔽白日，游子不顾返。思君令人老，岁月
忽已晚。  弃捐勿复道，努力加餐饭。

\begin{theorem}\label{the:theorem1}
犯我强汉者，虽远必诛\hfill —— 陈汤（汉）
\end{theorem}
\begin{subequations}
\begin{align}
y & = 1 \\
y & = 0
\end{align}
\end{subequations}
道可道，非常道。名可名，非常名。无名天地之始；有名万物之母。故常无，欲以观其妙；
常有，欲以观其徼。此两者，同出而异名，同谓之玄。玄之又玄，众妙之门。上善若水。水
善利万物而不争，处众人之所恶，故几于道。曲则全，枉则直，洼则盈，敝则新，少则多，
多则惑。人法地，地法天，天法道，道法自然。知人者智，自知者明。胜人者有力，自胜
者强。知足者富。强行者有志。不失其所者久。死而不亡者寿。

\begin{proof}
燕赵古称多感慨悲歌之士。董生举进士，连不得志于有司，怀抱利器，郁郁适兹土，吾
知其必有合也。董生勉乎哉？

夫以子之不遇时，苟慕义强仁者，皆爱惜焉，矧燕、赵之士出乎其性者哉！然吾尝闻
风俗与化移易，吾恶知其今不异于古所云邪？聊以吾子之行卜之也。董生勉乎哉？

吾因子有所感矣。为我吊望诸君之墓，而观于其市，复有昔时屠狗者乎？为我谢
曰：“明天子在上，可以出而仕矣！” \hfill —— 韩愈《送董邵南序》
\end{proof}

\begin{corollary}
  四川话配音的《猫和老鼠》是世界上最好看最好听最有趣的动画片。
\begin{alignat}{3}
V_i & =v_i - q_i v_j, & \qquad X_i & = x_i - q_i x_j,
 & \qquad U_i & = u_i,
 \qquad \text{for $i\ne j$;}\label{eq:B}\\
V_j & = v_j, & \qquad X_j & = x_j,
  & \qquad U_j & u_j + \sum_{i\ne j} q_i u_i.
\end{alignat}
\end{corollary}

迢迢牵牛星，皎皎河汉女。
纤纤擢素手，札札弄机杼。
终日不成章，泣涕零如雨。
河汉清且浅，相去复几许。
盈盈一水间，脉脉不得语。

\begin{example}
  大家来看这个例子。
\begin{equation}
\label{ktc}
\left\{\begin{array}{l}
\nabla f({\mbox{\boldmath $x$}}^*)-\sum\limits_{j=1}^p\lambda_j\nabla g_j({\mbox{\boldmath $x$}}^*)=0\\[0.3cm]
\lambda_jg_j({\mbox{\boldmath $x$}}^*)=0,\quad j=1,2,\cdots,p\\[0.2cm]
\lambda_j\ge 0,\quad j=1,2,\cdots,p.
\end{array}\right.
\end{equation}
\end{example}

\begin{exercise}
  请列出 Andrew S. Tanenbaum 和 W. Richard Stevens 的所有著作。
\end{exercise}

\begin{conjecture} \textit{Poincare Conjecture} If in a closed three-dimensional
  space, any closed curves can shrink to a point continuously, this space can be
  deformed to a sphere.
\end{conjecture}

\begin{problem}
 回答还是不回答，是个问题。
\end{problem}

如何引用定理~\ref{the:theorem1} 呢？加上 \cs{label} 使用 \cs{ref} 即可。妾发
初覆额，折花门前剧。郎骑竹马来，绕床弄青梅。同居长干里，两小无嫌猜。 十四为君妇，
羞颜未尝开。低头向暗壁，千唤不一回。十五始展眉，愿同尘与灰。常存抱柱信，岂上望夫
台。 十六君远行，瞿塘滟滪堆。五月不可触，猿声天上哀。门前迟行迹，一一生绿苔。苔深
不能扫，落叶秋风早。八月蝴蝶来，双飞西园草。感此伤妾心，坐愁红颜老。

\section{参考文献}
\label{sec:bib}
当然参考文献可以直接写 \cs{bibitem}，虽然费点功夫，但是好控制，各种格式可以自己随意改
写。

本模板推荐使用 BIB\TeX，分别提供数字引用（\texttt{thuthesis-numeric.bst}）和作
者年份引用（\texttt{thuthesis-author-year.bst}）样式，基本符合学校的参考文献格式
（如专利等引用未加详细测试）。看看这个例子，关于书的~\cite{tex, companion,
  ColdSources}，还有这些~\cite{Krasnogor2004e, clzs, zjsw}，关于杂志
的~\cite{ELIDRISSI94, MELLINGER96, SHELL02}，硕士论文~\cite{zhubajie,
  metamori2004}，博士论文~\cite{shaheshang, FistSystem01}，标准文
件~\cite{IEEE-1363}，会议论文~\cite{DPMG,kocher99}，技术报告~\cite{NPB2}，电子文
献~\cite{chuban2001,oclc2000}。
若使用著者-出版年制，中文参考文献~\cite{cnarticle}应增加
\texttt{key=\{pinyin\}} 字段，以便正确进行排序~\cite{cnproceed}。
另外，如果对参考文献有不如意的地方，请手动修改 \texttt{bbl} 文件。

有时候不想要上标，那么可以这样~\inlinecite{shaheshang}，这个非常重要。

有时候一些参考文献没有纸质出处，需要标注 URL。缺省情况下，URL 不会在连字符处断行，
这可能使得用连字符代替空格的网址分行很难看。如果需要，可以将模板类文件中
\begin{verbatim}
\RequirePackage{hyperref}
\end{verbatim}
一行改为：
\begin{verbatim}
\PassOptionsToPackage{hyphens}{url}
\RequirePackage{hyperref}
\end{verbatim}
使得连字符处可以断行。更多设置可以参考 \texttt{url} 宏包文档。

\section{公式}
\label{sec:equation}
\renewcommand\vec{\symbf}
\newcommand\mat{\symbf}
贝叶斯公式如式~(\ref{equ:chap1:bayes})，其中 $p(y|\vec{x})$ 为后验；
$p(\vec{x})$ 为先验；分母 $p(\vec{x})$ 为归一化因子。
\begin{equation}
\label{equ:chap1:bayes}
p(y|\vec{x}) = \frac{p(\vec{x},y)}{p(\vec{x})}=
\frac{p(\vec{x}|y)p(y)}{p(\vec{x})}
\end{equation}

论文里面公式越多，\TeX{} 就越 happy。再看一个 \pkg{amsmath} 的例子：
\newcommand{\envert}[1]{\left\lvert#1\right\rvert}
\begin{equation}\label{detK2}
\det\mat{K}(t=1,t_1,\dots,t_n)=\sum_{I\in\vec{n}}(-1)^{\envert{I}}
\prod_{i\in I}t_i\prod_{j\in I}(D_j+\lambda_jt_j)\det\vec{A}
^{(\lambda)}(\overline{I}|\overline{I})=0.
\end{equation}

前面定理示例部分列举了很多公式环境，可以说把常见的情况都覆盖了，大家在写公式的时
候一定要好好看 \pkg{amsmath} 的文档，并参考模板中的用法：
\begin{multline*}%\tag{[b]} % 这个出现在索引中的
\int_a^b\biggl\{\int_a^b[f(x)^2g(y)^2+f(y)^2g(x)^2]
 -2f(x)g(x)f(y)g(y)\dif x\biggr\}\dif y \\
 =\int_a^b\biggl\{g(y)^2\int_a^bf^2+f(y)^2
  \int_a^b g^2-2f(y)g(y)\int_a^b fg\biggr\}\dif y
\end{multline*}

其实还可以看看这个多级规划：
\begin{equation}\label{bilevel}
\left\{\begin{array}{l}
\max\limits_{{\mbox{\footnotesize\boldmath $x$}}} F(x,y_1^*,y_2^*,\cdots,y_m^*)\\[0.2cm]
\mbox{subject to:}\\[0.1cm]
\qquad G(x)\le 0\\[0.1cm]
\qquad(y_1^*,y_2^*,\cdots,y_m^*)\mbox{ solves problems }(i=1,2,\cdots,m)\\[0.1cm]
\qquad\left\{\begin{array}{l}
    \max\limits_{{\mbox{\footnotesize\boldmath $y_i$}}}f_i(x,y_1,y_2,\cdots,y_m)\\[0.2cm]
    \mbox{subject to:}\\[0.1cm]
    \qquad g_i(x,y_1,y_2,\cdots,y_m)\le 0.
    \end{array}\right.
\end{array}\right.
\end{equation}
这些跟规划相关的公式都来自于刘宝碇老师《不确定规划》的课件。

% \chapter{中华人民共和国}
\label{cha:china}

\section{其它例子}
\label{sec:other}

在第~\ref{cha:intro} 章中我们学习了贝叶斯公式~(\ref{equ:chap1:bayes})，这里我们复
习一下：
\begin{equation}
\label{equ:chap2:bayes}
p(y|\vec{x}) = \frac{p(\vec{x},y)}{p(\vec{x})}=
\frac{p(\vec{x}|y)p(y)}{p(\vec{x})}
\end{equation}

\subsection{绘图}
\label{sec:draw}

本模板不再预先装载任何绘图包（如 \pkg{pstricks，pgf} 等），完全由用户来决定。
个人觉得 \pkg{pgf} 不错，不依赖于 Postscript。此外还有很多针对 \LaTeX{} 的
 GUI 作图工具，如 XFig(jFig), WinFig, Tpx, Ipe, Dia, Inkscape, LaTeXPiX,
jPicEdt, jaxdraw 等等。

\subsection{插图}
\label{sec:graphs}

强烈推荐《\LaTeXe\ 插图指南》！关于子图形的使用细节请参看 \pkg{subcaption} 宏包的说明文档。

\subsubsection{一个图形}
\label{sec:onefig}
一般图形都是处在浮动环境中。之所以称为浮动是指最终排版效果图形的位置不一定与源文
件中的位置对应\footnote{This is not a bug, but a feature of \LaTeX!}，这也是刚使
用 \LaTeX{} 同学可能遇到的问题。如果要强制固定浮动图形的位置，请使用 \pkg{float} 宏包，
它提供了 \texttt{[H]} 参数，比如图~\ref{fig:xfig1}。
\begin{figure}[H] % use float package if you want it here
  \centering
  \includegraphics{njulogo0.pdf}
  \caption{利用 Xfig 制图}
  \label{fig:xfig1}
\end{figure}

大学之道，在明明德，在亲民，在止于至善。知止而后有定；定而后能静；静而后能安；安
而后能虑；虑而后能得。物有本末，事有终始。知所先后，则近道矣。古之欲明明德于天
下者，先治其国；欲治其国者，先齐其家；欲齐其家者，先修其身；欲修其身者，先正其心；
欲正其心者，先诚其意；欲诚其意者，先致其知；致知在格物。物格而后知至；知至而后
意诚；意诚而后心正；心正而后身 修；身修而后家齐；家齐而后国治；国治而后天下
平。自天子以至于庶人，壹是皆以修身为本。其本乱而未治者 否矣。其所厚者薄，而其所
薄者厚，未之有也！

\hfill —— 《大学》


\subsubsection{多个图形}
\label{sec:multifig}

如果多个图形相互独立，并不共用一个图形计数器，那么
用 \texttt{minipage} 或者\texttt{parbox} 就可以。否则，请参看
图~\ref{fig:big1-subcaptionbox}，它包含两个小图，分别是图~\ref{fig:subfig1}和
图~\ref{fig:subfig2}。推荐使用 \cs{subcaptionbox}，因为可以像
图~\ref{fig:big1-subcaptionbox} 那样对齐子图的标题，也可以使用 \pkg{subcaption}
宏包的 \cs{subcaption}（放在 minipage中，用法同\cs{caption}）或
是 \pkg{subfigure} 、\pkg{subtable}环境，像图~\ref{fig:big1-subfigure}，不要再
用 \cs{subfloat}、\cs{subfigure} 和 \cs{subtable}。

\begin{figure}[h]
  \centering%
  \subcaptionbox{第一个小图形\label{fig:subfig1}}[3cm] %标题的长度，超过则会换行，如下一个小图。
    {\includegraphics[height=3cm]{njulogo0.pdf}}%
  \hspace{4em}%
  \subcaptionbox{第二个小图形，注意这个图略矮些。如果标题很长的话，它会自动换行\label{fig:subfig2}}
      {\includegraphics[height=2cm]{njuname0.pdf}}
  \caption{包含子图形的大图形(subcaptionbox示例)}
  \label{fig:big1-subcaptionbox}
\end{figure}
\begin{figure}[h]
  \centering%
  \begin{subfigure}{3cm}
    \includegraphics[height=3cm]{njulogo0.pdf}
    \caption{第一个小图形}
  \end{subfigure}%
  \hspace{4em}%
  \begin{subfigure}{0.5\textwidth}
    \includegraphics[height=2cm]{njuname0.pdf}
    \caption{第二个小图形，注意这个图略矮些。subfigure中同一行的子图在顶端对齐。}
  \end{subfigure}
  \caption{包含子图形的大图形(subfigure示例)}
  \label{fig:big1-subfigure}
\end{figure}

古之学者必有师。师者，所以传道受业解惑也。人非生而知之者，孰能无惑？惑而不从师，
其为惑也，终不解矣。生乎吾前，其闻道也固先乎吾，吾从而师之；生乎吾後，其闻道也亦
先乎吾，吾从而师之。吾师道也，夫庸知其年之先後生於吾乎！是故无贵无贱无长无少，道
之所存，师之所存也。

嗟乎！师道之不传也久矣，欲人之无惑也难矣。古之圣人，其出人也远矣，犹且从师而问焉；
今之众人，其下圣人也亦远矣，而耻学於师。是故圣益圣，愚益愚。圣人之所以为圣，愚
人之所以为愚，其皆出於此乎？爱其子，择师而教之，於其身也，则耻师焉，惑焉。彼童子
之师，授之书而习其句读者，非吾所谓传其道、解其惑者也。句读之不知，惑之不解，或师
焉，或不焉，小学而大遗，吾未见其明也。巫医、乐师、百工之人不耻相师，  士大夫之族
曰“师”曰“弟子”之云者，则群聚而笑之。问之，则曰：彼与彼年相若也，道相似也，位
卑则足羞，官盛则近谀。呜呼！师道之不复，可知矣。巫医、乐师、百工之人。吾子不齿，
今其智乃反不能及，其可怪也欤！圣人无常师。孔子师郯子、苌子、师襄、老聃。郯子之徒，
其贤不及孔子。孔子曰：“三人行，必有我师。”是故弟子不必不如师，师不必贤於弟子。
闻道有先後，术业有专攻，如是而已。

如果要把编号的两个图形并排，那么小页就非常有用了：
\begin{figure}
\begin{minipage}{0.48\textwidth}
  \centering
  \includegraphics[height=2cm]{njulogo0.pdf}
  \caption{并排第一个图}
  \label{fig:parallel1}
\end{minipage}\hfill
\begin{minipage}{0.48\textwidth}
  \centering
  \includegraphics[height=2cm]{njulogo0.pdf}
  \caption{并排第二个图}
  \label{fig:parallel2}
\end{minipage}
\end{figure}

李氏子蟠，年十七，好古文、六艺，经传皆通习之，不拘於时，学於余。余嘉其能行古
道，作师说以贻之。

\hfill —— 韩愈（唐）

% \chapter{实现功能}
\section{硬件部分}
\begin{enumerate}
    \item \uline{五级流水线CPU}，实现除特殊访存指令外所有的逻辑、算术、移位、访存、跳转分支指令
    \item \uline{解决数据相关、控制相关问题}。实现跳转指令仅延迟槽中内容一定被执行，第二条以后的指令不受影响；实现访存指令无需延迟槽。
    \item 128KB内存同步读写
    \item VGA，键盘，音频\uline{多IO的调度}
    \item 代码段、数据段、栈段和IO段的\uline{内存划分}
    \item \uline{MMIO}实现多IO管理
\end{enumerate}
\section{软件部分}
\begin{enumerate}
    \item 实现了\uline{750行汇编代码构成的操作系统}，包含三个被调用40次左右的系统函数。以 \verb|printanasciicode| 为例，仅适用90条指令实现了软件控制写入,节省了大量存储器空间，
    \item \uline{代码实现显示字符功能}
    \item 解析./gdb ./fib ./hello ./mu四条用户指令
    \item \uline{gdb单步调试功能}，仿造实际控制流的跳转并输出
    \item 计算斐波那契数列，实现由\uline{用户控制运算第几项}
    \item 通过./mu开启电子琴使能，进入电子琴模式
    \item 开机欢迎动画
\end{enumerate}
\section{其它}
\begin{enumerate}
    \item \uline{搭建MIPS32运行时环境(AM)}，可将C程序编译到MIPS32上，并支持字符输出。结合ICS课程的库函数和测试程序，可提供对大型程序的支持接口
    \item 项目地址：\href{https://github.com/zhengzangw/NJU_MIPS}{https://github.com/zhengzangw/NJU\_MIPS}，提供了项目的说明文档、测试样例、开发工具(Makefile)和硬件实现的git记录。
    \item 由于文件内容多，我们特别为老师批改方便准备了一份文件目录\ref{sec:wjml}结构于本报告的最后。老师可根据该结构验收项目文件，减小验收负担。
\end{enumerate}
\chapter{组内分工}
\section{郑奘巍同学}
\begin{itemize}
    \item 工作量>35小时
    \item 完成了整个五级流水线CPU的编写，实现大量指令，解决数据相关和结构相关
    \item 实现了RAM，ROM和可读写ROM
    \item 实现了VGA，键盘，音频等IO的连接和MMIO
    \item 制作CPU测试样例和初始化IP核文件，完成硬件测试
    \item 运行时环境的搭建，包括Turing Machine接口, ld文件和Makefile文件的编写
\end{itemize}
\section{沈天琪同学}
\begin{itemize}
    \item 工作量约30小时
    \item 完成了整个操作系统的编写
    \item 前期资料收集与系统结构设计，编写小部分测试样例
    \item 欢迎动画的初始化文件编写，键盘模块与显存模块的设计
    \item 共同完成最终硬件测试
\end{itemize}


\chapter{硬件部分}
\section{五级流水线CPU的结构}
    五级流水线CPU分为：取指，译码，执行，访存，回写五个阶段，使得一个周期可以同时进行五条指令。
    %\fig{scale=0.4}{cpu1.jpg}{CPU基本结构}
    %\fig{scale=0.4}{cpu3.png}{五级流水线}
    基本结构如图\ref{fig:cpu1.jpg}所示，五级流水线执行流程如图\ref{fig:cpu3.png}所示。

    可以看到，整个流水线由时钟CLOCK控制，各个模块之间设置传输模块，在每个时钟周期的上升沿将前一阶段的数据传递到下一个模块。本图中还省略了从执行阶段将跳转地址传递到pc、为解决数据冲突从执行和访存连接到译码模块等连线。

    \subsection{取指模块} \label{sec:qzmk}
    取指模块中，ce控制整个CPU的开关，pc为当前指令位置。jmp\_flag\_i 和 stall分别是跳转地址和控制流水线暂停的变量，将在后文讨论。
    \begin{verilog}
    always @ (posedge clk) begin
		if (rst == `RSTENABLE) begin
			ce <= `CHIPDISABLE;
		end else begin
			ce <= `CHIPENABLE;
		end
	end
		
	always @ (posedge clk) begin
		if (ce == `CHIPDISABLE) begin
			pc <= 32'h00000000;
		end else if (stall[0] == `NOSTOP) begin
			if (jmp_flag_i == `JMP) begin //跳转
				pc <= jmp_target_address_i;
			end else begin //地址增加
				pc <= pc + 4'h4;
			end
		end
    end
    \end{verilog}

    \subsection{译码模块} \label{sec:ymmk}
    译码模块中，如果是R类型指令，我们使用如下的划分
    \begin{verilog}
        wire[5:0] op = inst_i[31:26]; //Instruction Code
        wire[4:0] sa= inst_i[10:6];   //sa
        wire[5:0] func = inst_i[5:0]; //Function Code
        wire[4:0] rt = inst_i[20:16]; //rt
        wire[4:0] rs = inst_i[25:21]; //rs
        wire[4:0] rd = inst_i[15:11]; //rd
    \end{verilog}
    如果是I类型指令，我们将immediate的值无符号扩展后放入变量imm中；如果是J型指令，我们将前26位放入jmp\_target\_address中传入下个模块。整个译码模块为一个选择器，先更根据op确定第一级，如果为EXE\_SPECIAL,再根据func确定具体指令。如果为EXE\_REGIMM,再根据rt确定具体指令。
    
    对于每条指令，通过设置wreg, reg1\_read, reg2\_read控制寄存器读写开关，如果读入无效，则将立即数传递给第二个操作数。以addu指令和第二操作数的加载为例，代码如下。
    \begin{verilog}
    case (op) 
        `EXE_SPECIAL: begin
            case (func)
                `EXE_ADDU: begin
                    wreg_o <= `WRITEENABLE; //需要写入，默认为rd
                    aluop_o <= `EXE_ADDU_OP; //操作为无符号加法
                    alusel_o <= `EXE_RES_ARITH; //取算术运算结果
                    reg1_read_o <= `READENABLE; //需要读第一个寄存器，默认为rs
                    reg2_read_o <= `READENABLE; //需要读第二个寄存器，默认为rt
                    instvalid <= `INSTVALID; //指令有效
                end
                ...
            endcase
        ...
    endcase

    always @(*) begin
        stallreq_for_reg2_loadrelate <= `NOSTOP;
        if (rst==`RSTENABLE) begin
            reg2_o <= `ZEROWORD;
        // Data Hazzard
        end else if ( pre_inst_is_load == 1'b1 && ex_wd_i == reg2_addr_o
                       && reg2_read_o == 1'b1 ) begin
            stallreq_for_reg2_loadrelate <= `STOP; //load相关
        end else if ((reg2_read_o==`READENABLE)&&(ex_wreg_i==`WRITEENABLE)&&       (ex_wd_i==reg2_addr_o)) begin
            reg2_o <= ex_wdata_i; //来自执行模块的数据
        end else if ((reg2_read_o==`READENABLE)&&(mem_wreg_i==`WRITEENABLE)&&      (mem_wd_i==reg2_addr_o)) begin
            reg2_o <= mem_wdata_i; //来自访存模块的数据
        end else if (reg2_read_o == `READENABLE) begin
            reg2_o <= reg2_data_i; //需要读数据
        end else if (reg2_read_o == `READDISABLE) begin
            reg2_o <= imm; //无需读寄存器数据，默认提供立即数
        end else begin
            reg2_o <= `ZEROWORD;
        end 
    end
    \end{verilog}

    \subsection{执行模块}
    执行模块由逻辑运算、移位运算、转移、算术模块、乘法模块、除法模块和跳转六个部分并行组成，根据取指模块传递进来的指令类型，选择相应的输出。其中乘法和除法向hi,lo输出，这里没有给出。具体的指令实现在\ref{sec:zlsx}节中给出。同时，由于除法实现的特殊性，我们专门设置了div模块来实现除法指令。
    \begin{verilog}
        always @(*) begin 
            wd_o <= wd_i;
		    if (((aluop_i == `EXE_ADD_OP)||(aluop_i == `EXE_SUB_OP))&&(overflow)  ) begin
			    wreg_o <= `WRITEDISABLE; //无需写入
		    end else begin
				wreg_o <= wreg_i; //写入与译码模块要求相同
		    end
        case (alusel_i)
            `EXE_RES_LOGIC: begin 
                wdata_o <= logicout; //逻辑模块
            end
            `EXE_RES_SHIFT: begin
                wdata_o <= shiftres; //移位模块
            end
            `EXE_RES_MOVE: begin
                wdata_o <= moveres; //转移模块
            end
            `EXE_RES_ARITH: begin
                wdata_o <= arithres; //计算模块
            end
            `EXE_RES_JUMP: begin
                wdata_o <= link_address_i; //跳转模块
            end
            default: begin
                wdata_o <= `ZEROWORD;
            end
        endcase
    end 
    \end{verilog}
    \subsection{访存模块}
        对于非访存指令，在访存模块只是进行数据的传递。对于访存指令，则通过译码模块传递进来的参数进行处理。该模块将向外连接到RAM中，以4字节为单位，向RAM提供地址，是否写入，写入数据等信息。其中，mem\_sel将控制更改的内容属于以4字节为单位的RAM中的哪几位。
    \subsection{寄存器与回写模块}
        回写模块没有单独实现，而是直接写在寄存器内。寄存器模块regfile实现了0-31号寄存器，hilo\_reg则实现了hi和lo寄存器。
    
\section{指令实现} \label{sec:zlsx}
\subsection{逻辑指令}
\subsubsection{基本实现}
    逻辑指令是最简单的部分，共十七条指令，列举如下（uext为对立即数进行零扩展，逻辑位移与算术位移在表中没有体现，可自行区分）
    \tabpc{4}{
        and & rd $\leftarrow$ rs\&rt & or & rd $\leftarrow$ rs|rt \\ \hline
        xor & rd $\leftarrow$ rs xor rt & nor & rd $\leftarrow$ \!(rs xor rt) \\ \hline
        andi& rt $\leftarrow$ rs\&uext(imm) & xori & rt $\leftarrow$ rs xor uext(imm) \\ \hline 
        lui & rt $\leftarrow$ \{imm, 16'b0\} & sll & rd $\leftarrow$ rt<<sa \\ \hline 
        srl & rd $\leftarrow$ t>>sa & sra & rd $\leftarrow$ st>>sa \\ \hline
        sllv & rd $\leftarrow$ rt<<rs[4:0] & rd $\leftarrow$ srlv & rt>>rs[4:0] \\ \hline
        srav & rd $\leftarrow$ rt >> rs[4:0] & pref & nop 
    }
    实现了这些指令的同时，也实现了nop，ssnop，move等指令（语法糖）。这些指令的解耦方式基本上就是经历上述了五个模块，在取码后译码执行并回写。
    \subsubsection{数据相关问题}
    数据冲突指的是流水线中执行的几条指令中，一条指令依赖于前面指令的执行结果，具体又分为RAW,WAR,WAW三种情况。以如下代码为例:
    \begin{verbatim}
        ori $1, $0, 0x1100
        ori $2, $1, 0x0020
    \end{verbatim}
    则在第二条指令的译码阶段，第一条指令还没到达回写阶段导致取得的值不正确。译码阶段和回写阶段相差了两个时钟周期，如果使用插入空指令的方法解决，则浪费时间与空间。我们通过数据前推，将计算结果从执行和访存阶段直接更新到后面指令的译码阶段，从而使其获得正确的结果。代码实现可见\ref{sec:ymmk}节，原理如下图所示。
    %\fig{scale=0.4}{cpu5.png}{数据相关}

\subsection{移位指令}
    移位指令共6条如下。实现时，在译码阶段确定是否需要转移。对于hi,lo的移位操作，也需要考虑上述的数据相关问题。
    \tabpc{4}{
        movn & rd $\leftarrow$ rs if rt$\neq$0 & movz & rd $\leftarrow$ rs if rt=0  \\ \hline
        mfhi & rd $\leftarrow$ hi & mflo & rd $\leftarrow$ lo \\ \hline
        mthi& hi $\leftarrow$ rs & mtlo & lo $\leftarrow$ rs 
    }
\subsection{算术指令}
\subsubsection{基本实现}
    简单的算术指令共15条，实现时注意注意手册对指令的定义。add指令在加法溢出时不保存结果（溢出异常没有实现），addu则不进行溢出检查。其它指令也类似如此。比较特别的指令有slt保存大小比较结果以及clo保存高位1的个数，clz保存高位0的个数。
    \tabpc{4}{
        add(u) & rd $\leftarrow$ rs+rt & sub(u) & rd $\leftarrow$ rs-rt \\ \hline
        slt(u) & rd $\leftarrow$ rs<rt & addi(u) & rd $\leftarrow$ rs+uext(imm) \\ \hline
        slti(u)& rt $\leftarrow$ rs<uext(imm) & clo & 高位连续1的个数 \\ \hline 
        clz & 高位连续0的个数 & mul & rd $\leftarrow$ rs*rt \\ \hline 
        mult(u) & \{hi, lo\} $\leftarrow$ rs*rt & div(u) & \{hi, lo\} $\leftarrow$ rs/rt
    }
\subsubsection{除法}
    我们使用了一种朴素的方法来实现除法。试商法需要32个周期进行除法运算，流程如下。为实现除法模块，我们参考了试商法的状态机设计，并添加流水线控制模块，当进入除法模块后停止整个流水线的运行，直到除法完成再恢复。可见本设计中除法的代价是很大的。
    %\fig{scale=0.6}{div.png}{试商法流程}

\subsection{跳转指令}
\subsubsection{基本实现}
    跳转指令共4条，分支指令10条这里不具体列出，指令的区别在于分支条件不同以及是否保存返回地址。在取指阶段判断是否跳转，如果需要则暂停流水线的运行，以解决结构冲突。以bltzal指令为例，该指令发生跳转当且仅当$rs<0$，且保存指令后地址为返回地址到第31号寄存器。跳转部分代码可见\ref{sec:qzmk}节。
    \begin{verilog}
        `EXE_BLTZAL: begin
            wreg_o <= `WRITEENABLE; //需要写入寄存器
            aluop_o <= `EXE_BLTZAL_OP; //操作指令
            alusel_o <= `EXE_RES_JUMP; //操作类型
            reg1_read_o <= `READENABLE; //需要读一个寄存器
            reg2_read_o <= `READDISABLE;
            wd_o <= 5'b11111; //需要将返回地址写入31号寄存器
            link_addr_o <= pc_plus_8;  //返回地址
            instvalid <= `INSTVALID; //指令有效
            if (reg1_o[31]==1'b1) begin
                jmp_target_address_o <= pc_plus_4 + imm_sll2_signedext; //跳转地址和
                jmp_flag_o <= `JMP; //需要跳转
                next_inst_in_delayslot_o <= `INDELAYSLOT; //下一条指令处于延迟槽
            end
        end
    \end{verilog}
    \begin{verilog}
    always @(*) begin
        if (rst == `RSTENABLE) begin
            stall <= 6'b000000;
        end else if (stallreq_ex == `STOP) begin 
            stall <= 6'b001111; //来自执行模块的停止流水线请求
        end else if (stallreq_id == `STOP) begin 
            stall <= 6'b000111; //来自译码模块的停止流水线请求
        end else begin
            stall <= 6'b000000;
        end
    end
    \end{verilog}

\subsection{访存指令}
    访存过程中，我们使用IP核作为储存器，让IP核的频率为CPU的两倍，可实现同步访存。访存指令共8条如下：
    \tabpc{4}{
        lb & 字节加载 & lbu & 无符号字节加载 \\ \hline
        lh & 半字加载 & lhu & 半字无符号加载 \\ \hline
        lw & 字加载 & sw & 字储存 \\ \hline 
        sb & 字节储存 & sh & 半字储存
    }
    以字节加载指令为例，其代码如下
    \begin{verilog}
        case (aluop_i)
			`EXE_LB_OP: begin
                mem_addr_o <= mem_addr_i;
                mem_we <= `WRITEDISABLE;
                mem_ce_o <= `CHIPENABLE;
                case (mem_addr_i[1:0]) //根据最后两位选择写入值和选择变量
                    2'b00: begin
                        wdata_o <= {{24{mem_data_i[31]}}, mem_data_i[31:24]};
                        mem_sel_o <= 4'b1000;
                    end
                    2'b01: begin
                        wdata_o <= {{24{mem_data_i[23]}}, mem_data_i[23:16]};
                        mem_sel_o <= 4'b0100;
                    end
                    2'b10: begin
                        wdata_o <= {{24{mem_data_i[15]}}, mem_data_i[15:8]};
                        mem_sel_o <= 4'b0010;
                    end
                    2'b11: begin
                        wdata_o <= {{24{mem_data_i[7]}}, mem_data_i[7:0]};
                        mem_sel_o <= 4'b0001;
                    end
                    default: begin
                        wdata_o <= `ZEROWORD;
                    end
                ...
			endcase
        end
    \end{verilog}

    在访存指令中，也有一类数据相关，称为load相关，当需要使用的数据是由访存阶段读入时，就会出现此类问题。为此，当下一条指令需要用到上一条加载指令的结果时，我们也暂停流水线。

    最终，我们的CPU的数据转移流大致如下。
    %\fig{scale=0.4}{cpu6.jpg}{最终CPU的数据转移}

\section{指令测试}
    为了测试指令的正确性，我们利用仿真文件sopc\_tst.v对文件夹program/tst下的汇编测试样例进行仿真，并编写了showreg.do的脚步文件显示需要观测的变量值。

\section{ROM, RAM与IO}
\subsection{储存空间构建与IO配置}
    我们利用IP核构建储存空间。储存设备如下：
    \begin{table}[H]
        \centering
        \begin{tabular}{|c|c|c|c|}
            \hline
            储存设备 & 大小 & 功能 & 初始内容 \\ \hline 
            ROM & 4KB & 指令存储器 & 指令(.text节) \\
            Writable ROM & 1KB & 可读写指令存储器（支持gdb） & 无 \\
            RAM & 128KB & 内存 & 无 \\
            GRAM & 1KB & 显存（存储字符）& 开机画面\\
            \hline 
        \end{tabular}
    \end{table}

    为实现Writable ROM，我们直接使用了reg型变量。对于其它储存器，我们利用了IP核进行构成。类似于内存条的位扩展，我们使用四个IP核构成ROM。每个ROM中相邻的两位在地址上相差4位。

    IO设备复用了之前实验的代码，VGA以字符为单位显示，键盘支持全键位、组合键和大小写，电子琴支持和声。将储存器与CPU打包成sopc，我们构筑整个系统结构如下。
    %\fig{scale=0.4}{cpu4.png}{系统架构}

\subsection{MMIO}
    为了提供对设备的访问，我们使用了内存映射管理所有IO。通过一系列判断语句，模拟了一个MMIO管理器，使得可以通过访存指令获得键盘输入、向VGA写入字符、控制电子琴。具体分配如下。
    \begin{table}[H]
        \centering
        \begin{tabular}{|c|c|c|}
            \toprule
            ADDRESS & DESCRIPTION & FUNCTION\\ \midrule
            0x000000 - 0x003900 & ROM & for instruction executable; lw only \\
            0x003900 - 0x004000 & RAM & for instruction	executable; lw,sw only \\
            0x004000 - 0x005000 & GRAM & for video	readable, writable; access by vga(640*480)\\
            0x005000 & keyboard enable & lb only \\
            0x005004 & keyboard ascii & lb only \\
            0x005008 & audio enable & wb only \\
            - 0x006000 & Reserve & \\
            0x006000 - 0x008000	& RAM:Data Section & \\
            0x008000 - 0x010000 & RAM:Heap & \\
            \bottomrule
        \end{tabular}
    \end{table}
\section{运行时环境AM}
    AM的概念是在拔尖班余子濠老师在计算机系统课程的PA作业中介绍的，我们摘录PA讲义部分如下：
    %\fig{scale=0.5}{PA.png}{AM的概念}

    在PA中，我们补充完成了x86-nemu的AM，并完成了虚拟机上的虚拟机x86-navy的AM。我们设想，通过完成相应的MIPS32的AM的设计，我们可以把PA中的小程序搬运到我们的CPU上（甚至仙剑奇侠传）。这个想法让我激动不已。然而由于没有实现中断以及时间问题，最终我们没有完成全部的设计。

    我们实现了AM的Turing Machine的接口。
    \begin{verilog}
        int video = 0x4001; //显存起始位置
        void _putc(char ch) {
            asm volatile(
                "move $24, %0\n\t"
                "sb %1, 0x0($24)\n\t"
                :
                :"r"(video), "r"(ch)
            );
            video++;
        }

        void _trm_init() {
            asm volatile("ori $sp,$0,0x9000"); //设置栈区
            _halt(main());
        }

        void _halt(int code) {
        if (code) {
            _putc('W'); _putc('A'); //HIT THE BAD TRAP
        } else {
            _putc('A'); _putc('C'); //HIT THE GOOD TRAP
        }
            while (1);
        }
    \end{verilog}

    在loader.ld文件中，将代码节，数据节设置到与上节表格相对应的位置，将起始符号改为\_trm\_init，修改Makefile文件和bash脚本使其支持mips32 AM的编译，编译cputest下的dummy.c文件或sum.c文件，可以看到屏幕上出现了AC。这就相当与PA实验中的HIT THE GOOD TRAP. 由于已经有了\_putc函数，我们可以调用在PA中自己写的klib中的stdio.h中的printf来输出我们想要的格式了。
    
    理论上如果实现了系统调用，那么就可以把操作系统Nanos-lite和仙剑程序跑在我们的CPU上。可惜的是，余老师告诉我们由于没有cache和一些其它原因，这个过程还需要额外调用板上的设备以及及时跑起来可能也只能显示第一张图片就卡死了。考虑时间因素在内，我们最后放弃了这个实现。

    如此一来，借助AM的力量，我们便可以把这两门课程的内容打通，从而更好的理解计算机的体系结构。

\section{遇到的问题与解决方案}
    \begin{itemize}
        \item 实现CPU时的困难：流水线停止，div指令，累乘指令实现等，这些问题都有一定难度，通过阅读资料可以看到前人优秀的实现方法。通过学习可以用到我们的CPU中。
        \item gdb的支持：原先的设计中代码区和数据区分类，代码区无法读写。为了实现对gdb的支持，强行用reg类型实现一个支持读写的代码段，以支持操作系统的设计。
        \item IP核接口不足，无法同时支持CPU读写和VGA访问：由于显存本身较小，可以同时保存两份IP核的储存空间。虽然不是好方法，但是可以work。
        \item MMIO的调试：MMIO的调试比其它指令的调试要复杂，因为此时仿真能提供的帮助很小。通过设计对于不同情况有不同反应的汇编代码和将一些信息直接从LED灯中输出可以帮助调试。
    \end{itemize}
\section{实验启示}
    \begin{itemize}
        \item 项目可以提升自己的多方面能力：通过完成本次实验，不仅提高了自己的Verilog代码能力，加深对CPU和计算机体系的理解，还提高了自己的团队合作能力，项目维护能力，Makefile,loader文件甚至是报告制作的能力。相比与老师提供的往年优秀实验报告的内容，我们两人组队不仅完成了其中三个人队伍的所有内容，还将CPU从单周期换位流水线并在各方面进行了拓展。
        \item RTFM，STFW：遇到问题时，读帮助文档和上网搜索问题资料往往能够解决。
        \item 课程的融汇贯通十分重要：将ICS课程中AM引入到数电实验中来，开启了新世界的大门。虽然我们并没有实现太多具体的程序，但是在以后学弟学妹的制作中，可以尝试把仙剑搬入数电实验之中。
    \end{itemize}

\chapter{软件部分}
\section{操作系统的定义}
    操作系统（英语：operating system，缩写作 OS）是管理计算机硬件与软件资源的计算机程序，同时也是计算机系统的内核与基石。操作系统需要处理如管理与配置内存、决定系统资源供需的优先次序、控制输入与输出设备、操作网络与管理文件系统等基本事务。操作系统也提供一个让用户与系统交互的操作界面。
\section{操作系统中的约定}
    \hspace{-3cm}
    \begin{tabular}{|*{8}{c|}}
    \hline
    \$0 & 恒为0 & \$8 & Fib temp data & \$16 & gdb temp data & \$24 & 列标号 \\ \hline
    \$1	& 留作编译器使用 & \$9 & Fib temp data & \$17 & gdb/testcase ret addr &	\$25 & 行标号\\ \hline
    \$2	& 上一个ASCII输入 &	\$10 & Fib temp data &\$18 & gdb temp data &\$26 & Mu temp data \\ \hline
    \$3	& sys temp data	& \$11 & Fib temp data & \$19 & gdb temp data & \$27 & 未处理行指针寄存器 \\ \hline
    \$4 & sys temp data	& \$12 & Fib temp data &\$20 & gdb temp data & \$28 & ENTER信号\\ \hline
    \$5	 & sys temp data & \$13 & Fib temp data(10) & \$21 &	Testcase use & \$29 & 恒为4\\ \hline
    \$6 & sys temp data & \$14 & gdb temp data(addr) & \$22 & Testcase use & \$30 & 栈指针\\ \hline
    \$7  & 在用户程序中保存\$ra & \$15 & gdb temp data(addr) & \$23 & ret data & \$31 & ret addr(\$ra)\\ \hline
    \end{tabular}
\section{屏幕显示功能}
    最主要的是实现控制流的可控，保证其在主循环的轨道内运行。
    %\fig{scale=1.0}{Picture1.png}{设计方案}
    %\fig{scale=1.0}{Picture4.png}{输出逻辑}
    我们将输入分为三种情形：1.不以./开头，我们认为这是普通文本，不执行比对。2.某一个命令，则此时执行对应命令。3.错误命令，则此时输出错误信息。最终均会回到主循环。\par
    代码如下（这里省略了有关具体比对的细节）。
    %\fig{scale=1.0}{Picture2.png}{代码}
    值得注意的是，我们系统中没有一个通常意义上的键盘缓冲区，我们是借显存作为键盘缓冲区的，这就是为什么我们在寄存器堆中需要一个未处理行指针寄存器了。\par 
    %\fig{scale=0.8}{Picture3.png}{读取函数}
    读取键盘输入的函数十分简单，如图\ref{fig:Picture3.png}所示。而输出函数逻辑较为复杂，如\ref{fig:Picture4.png}所示。
    函数主体约为如上。这部分的功能几乎与实验11等价。

\clearpage
\section{斐波那契数}
    逻辑十分简单，首先接受输入，检测到ENTER信号之后从显存中读取数据，之后迭代运算，最后将结果除10，余数放入栈中，结果参与下一次除法直到为0，最后从栈中获得数据输出，形成十进制输出的情形。
    第一部分，接收数据，等待ENTER输入，如图\ref{fig:Picture13.png}\ref{fig:Picture14.png}。
    %\fig{scale=0.6}{Picture13.png}{输入部分}
    %\fig{scale=0.6}{Picture14.png}{输入部分}
    %\fig{scale=0.6}{Picture15.png}{第二部分}
    %\fig{scale=0.6}{Picture16.png}{第三部分}
    第二部分，ENTER信号产生并被接收后跳转于此，此时数据已经完整，进入迭代，如图\ref{fig:Picture15.png}。第三部分，另一个将数据写入栈中的循环，为输出做准备，如图\ref{fig:Picture16.png}。
    %\fig{scale=0.6}{Picture17.png}{第四部分}
    最后一部分，从栈中得到数据并不断将其写入显存（调用系统函数）由于最后写入了ENTER信号，这意味着我们最终要将这一个额外的信号清除，将未处理行指针寄存器写入正确的值。
    %\fig{scale=0.4}{Picture18.png}{斐波那契数效果演示}
\section{Hello World}
    Hello world开机界面和命令展示如\ref{fig:Picture19.png}\ref{fig:Picture20.png}
    %\fig{scale=0.5}{Picture19.png}{开机界面效果演示}
    %\fig{scale=0.4}{Picture20.png}{hello命令效果演示}
\section{电子琴}
    电子琴程序的功能就是打开电子琴使能，并构建一个检测到backspace输入退出的程序
    %\fig{scale=0.6}{Picture21.png}{电子琴代码}
\section{单步调试}
    我们的单步调试实现受到了操作系统异常处理程序的启发。
    %\fig{scale=0.6}{Picture22.png}{Linux异常处理} 

    而在我们的实现中，由于没有完成异常处理相关的协处理器与指令，这使得我们不得不采用一种模拟的手段来实现相关机制。在实现中，有testcase存放在0x2f00的位置，大体流程如\ref{fig:Picture23.png}
    %\fig{scale=1.0}{Picture23.png}{单步调试流程} 

    \clearpage
    这三条特殊的指令为：
    \begin{verbatim}
    .org 0x2ef0
        nop
        jr  $ra             # jump to $ra
        ori $17,$ra,0x0
    \end{verbatim}
    nop的位置就是将来会被替代的位置。\par
    gdb相当于一个监视者，一次只允许一条指令的运行，当然这样实现是不完美的，他无法处理程序中有跳转的情况，在顺序执行的指令中，它等价于图\ref{fig:Picture24.png}\par 
    
    %\fig{scale=1.0}{Picture24.png}{gdb流程}
\clearpage
\section{遇到的问题与解决方案}
    MARS不支持高级编译选项，所以我们搭建了一个交叉编译工具链。运用makefile实现了流程化的操作，当然，我们自己也写了一些小工具用以将各种初始化文件格式转化为了mif，见图\ref{fig:Picture34.png}(这部分请参考我写的mips交叉编译工具链的配置与使用——基于\verb|Ubuntu18.04|)\par
    有时寄存器在使用时的未初始化会造成相当令人意外的结果，这时就要严格执行所用寄存器的保存于初始化。\par
    MARS不支持高级编译选项，gdb无法在此调试：这使得我们只好在实机上调试，唯一问题就是每一次编译都十分缓慢，这就要求我们预先写好testcase，从简单开始，这样可以大大降低复杂度，有效排除硬件问题。
    %\fig{scale=0.6}{Picture35.png}{测试文件}
    %\fig{scale=1.0}{Picture34.png}{Makefile}
\section{实验启示}
    一定要多交流，有时你想不到的，或许别人就有一个很好的解决方案，互相查漏补缺，效率是很高的。

\chapter{文件结构} \label{sec:wjml}
    以下列出了重要文件及其内容。在主文件夹中大量文件属于IP核，为自动生成的文件。各部分文件夹下有Verilog文件和mif初始化文件，cpu文件夹下为CPU模块，具体见硬件内容描述。os.s为操作系代码。
    \dirtree{%
        .1 NJU\_MIPS.
        .2 README.md.
        .2 macro.v\DTcomment{宏定义}.
        .2 NJU\_MIPS.v\DTcomment{顶层模块}.
        .2 sopc.v\DTcomment{CPU和储存器模块}.
        .2 cpu\DTcomment{MIPS五级流水线CPU代码}.
        .3 cpu.v\DTcomment{CPU顶层模块}.
        .2 data\_RAM\DTcomment{128K RAM代码}.
        .2 inst\_RAM\DTcomment{指令储存器代码}.
        .2 keyboard\DTcomment{键盘代码}.
        .2 vga\DTcomment{VGA代码}.
        .2 audio\DTcomment{电子琴代码}.
        .2 simulation\DTcomment{仿真文件}.
        .3 modelsim.
        .4 sopc\_tst.v\DTcomment{对sopc.v的测试文件}.
        .2 program\DTcomment{初始化代码及生成工具}.
        .3 tst\DTcomment{测试样例}.
        .3 os.s\DTcomment{操作系统代码}.
        .3 inst\_rom.mif\DTcomment{初始化文件}.
        .2 AM\DTcomment{运行时环境（修改自PA实验）}.
        .3 am.
        .4 arch.
        .5 mips32\DTcomment{MIPS32运行时环境}.
        .3 tests\DTcomment{测试文件}.
    }
% \chapter{代码样式}
\begin{algorithm}
\caption{test}
\begin{algorithmic}
    \If {$i\geq maxval$}
        \State $i\gets 0$
    \Else
        \If {$i+k\leq maxval$}
            \State $i\gets i+k$
        \EndIf
    \EndIf
\end{algorithmic}
\end{algorithm}

\begin{shell}
    >> hello
    >> python
    shell code
\end{shell}

\begin{commandshell}
    ls -al
    cd /usr/lib
\end{commandshell}

\begin{latex}
    \setvar{\myname}{\name}
    \newcommand\myaddress{}
    \newcommand{\address}[1]{\renewcommand{\myaddress}{#1}}
\end{latex}

\begin{cdisplay}
    for (int k=0;k<n;++k){
        for (int i=0;i<n;++i){
            for (int j=0;j<n;++j){
                a[i][j][k] = 0;
            }
        }
    }
    // 这个不是Floyed
\end{cdisplay}

\begin{algorithm}
    \caption{代码标题}
    \begin{cpseudo}
        for (int k=0;k<n;++k){
            for (int i=0;i<n;++i){
                for (int j=0;j<n;++j){
                    a[i][j][k] = 0;
                }
            }
        }
        // 这个不是Floyed
    \end{cpseudo}
\end{algorithm}

\begin{cplus}
    #include <bits/stdc++.h>
    using namespace std;
    int main(){
        cout << "Hello World!" << endl;
        return 0;
    }
\end{cplus}

\begin{code}[language=Assembler]
    mov %eax, $ebx
    sub 3, %eax
\end{code}

\begin{python}
    import pandas as pd
    pd.read_csv("name.csv")
\end{python}

\begin{verilog}
    always @ (posedge clk) begin
		if (rst == `RSTENABLE) begin
			ce <= `CHIPDISABLE;
		end else begin
			ce <= `CHIPENABLE;
		end
    end
\end{verilog}
% \backmatter
% \listoffigures
% \listoftables
% \listofequations
% \bibliographystyle{ref/numeric} % ref/numeric,ref/author-year,plainnat,IEEEtran
% \bibliography{ref/refs}
% % 如果使用声明扫描页，将可选参数指定为扫描后的 PDF 文件名，例如：
% \begin{acknowledgement}[scan-statement.pdf]
\begin{acknowledgement}
  衷心感谢导师 xxx 教授和物理系 xxx 副教授对本人的精心指导。他们的言传身教将使
  我终生受益。

  在美国麻省理工学院化学系进行九个月的合作研究期间，承蒙 xxx 教授热心指导与帮助，不
  胜感激。感谢 xx 实验室主任 xx 教授，以及实验室全体老师和同学们的热情帮助和支
  持！本课题承蒙国家自然科学基金资助，特此致谢。

  感谢 \LaTeX 和 thuthesis\cite{thuthesis}，帮我节省了不少时间。
\end{acknowledgement}

% %% 附录
% \begin{appendix}
  % \chapter{外文资料原文}
\label{cha:engorg}

\title{The title of the English paper}

\textbf{Abstract:} As one of the most widely used techniques in operations
research, \emph{ mathematical programming} is defined as a means of maximizing a
quantity known as \emph{bjective function}, subject to a set of constraints
represented by equations and inequalities. Some known subtopics of mathematical
programming are linear programming, nonlinear programming, multiobjective
programming, goal programming, dynamic programming, and multilevel
programming$^{[1]}$.

It is impossible to cover in a single chapter every concept of mathematical
programming. This chapter introduces only the basic concepts and techniques of
mathematical programming such that readers gain an understanding of them
throughout the book$^{[2,3]}$.


\section{Single-Objective Programming}
The general form of single-objective programming (SOP) is written
as follows,
\begin{equation}\tag*{(123)} % 如果附录中的公式不想让它出现在公式索引中，那就请
                             % 用 \tag*{xxxx}
\left\{\begin{array}{l}
\max \,\,f(x)\\[0.1 cm]
\mbox{subject to:} \\ [0.1 cm]
\qquad g_j(x)\le 0,\quad j=1,2,\cdots,p
\end{array}\right.
\end{equation}
which maximizes a real-valued function $f$ of
$x=(x_1,x_2,\cdots,x_n)$ subject to a set of constraints.

\newcommand\Real{\mathbf{R}}
\newtheorem{mpdef}{Definition}[chapter]
\begin{mpdef}
In SOP, we call $x$ a decision vector, and
$x_1,x_2,\cdots,x_n$ decision variables. The function
$f$ is called the objective function. The set
\begin{equation}\tag*{(456)} % 这里同理，其它不再一一指定。
S=\left\{x\in\Real^n\bigm|g_j(x)\le 0,\,j=1,2,\cdots,p\right\}
\end{equation}
is called the feasible set. An element $x$ in $S$ is called a
feasible solution.
\end{mpdef}

\newtheorem{mpdefop}[mpdef]{Definition}
\begin{mpdefop}
A feasible solution $x^*$ is called the optimal
solution of SOP if and only if
\begin{equation}
f(x^*)\ge f(x)
\end{equation}
for any feasible solution $x$.
\end{mpdefop}

One of the outstanding contributions to mathematical programming was known as
the Kuhn-Tucker conditions\ref{eq:ktc}. In order to introduce them, let us give
some definitions. An inequality constraint $g_j(x)\le 0$ is said to be active at
a point $x^*$ if $g_j(x^*)=0$. A point $x^*$ satisfying $g_j(x^*)\le 0$ is said
to be regular if the gradient vectors $\nabla g_j(x)$ of all active constraints
are linearly independent.

Let $x^*$ be a regular point of the constraints of SOP and assume that all the
functions $f(x)$ and $g_j(x),j=1,2,\cdots,p$ are differentiable. If $x^*$ is a
local optimal solution, then there exist Lagrange multipliers
$\lambda_j,j=1,2,\cdots,p$ such that the following Kuhn-Tucker conditions hold,
\begin{equation}
\label{eq:ktc}
\left\{\begin{array}{l}
    \nabla f(x^*)-\sum\limits_{j=1}^p\lambda_j\nabla g_j(x^*)=0\\[0.3cm]
    \lambda_jg_j(x^*)=0,\quad j=1,2,\cdots,p\\[0.2cm]
    \lambda_j\ge 0,\quad j=1,2,\cdots,p.
\end{array}\right.
\end{equation}
If all the functions $f(x)$ and $g_j(x),j=1,2,\cdots,p$ are convex and
differentiable, and the point $x^*$ satisfies the Kuhn-Tucker conditions
(\ref{eq:ktc}), then it has been proved that the point $x^*$ is a global optimal
solution of SOP.

\subsection{Linear Programming}
\label{sec:lp}

If the functions $f(x),g_j(x),j=1,2,\cdots,p$ are all linear, then SOP is called
a {\em linear programming}.

The feasible set of linear is always convex. A point $x$ is called an extreme
point of convex set $S$ if $x\in S$ and $x$ cannot be expressed as a convex
combination of two points in $S$. It has been shown that the optimal solution to
linear programming corresponds to an extreme point of its feasible set provided
that the feasible set $S$ is bounded. This fact is the basis of the {\em simplex
  algorithm} which was developed by Dantzig as a very efficient method for
solving linear programming.
\begin{table}[ht]
\centering
  \centering
  \caption*{Table~1\hskip1em This is an example for manually numbered table, which
    would not appear in the list of tables}
  \label{tab:badtabular2}
  \begin{tabular}[c]{|m{1.5cm}|c|c|c|c|c|c|}\hline
    \multicolumn{2}{|c|}{Network Topology} & \# of nodes &
    \multicolumn{3}{c|}{\# of clients} & Server \\\hline
    GT-ITM & Waxman Transit-Stub & 600 &
    \multirow{2}{2em}{2\%}&
    \multirow{2}{2em}{10\%}&
    \multirow{2}{2em}{50\%}&
    \multirow{2}{1.2in}{Max. Connectivity}\\\cline{1-3}
    \multicolumn{2}{|c|}{Inet-2.1} & 6000 & & & &\\\hline
    \multirow{2}{1.5cm}{Xue} & Rui  & Ni &\multicolumn{4}{c|}{\multirow{2}*{\thuthesis}}\\\cline{2-3}
    & \multicolumn{2}{c|}{ABCDEF} &\multicolumn{4}{c|}{} \\\hline
\end{tabular}
\end{table}

Roughly speaking, the simplex algorithm examines only the extreme points of the
feasible set, rather than all feasible points. At first, the simplex algorithm
selects an extreme point as the initial point. The successive extreme point is
selected so as to improve the objective function value. The procedure is
repeated until no improvement in objective function value can be made. The last
extreme point is the optimal solution.

\subsection{Nonlinear Programming}

If at least one of the functions $f(x),g_j(x),j=1,2,\cdots,p$ is nonlinear, then
SOP is called a {\em nonlinear programming}.

A large number of classical optimization methods have been developed to treat
special-structural nonlinear programming based on the mathematical theory
concerned with analyzing the structure of problems.
\begin{figure}[h]
  \centering
  \includegraphics{njulogo0.pdf}
  \caption*{Figure~1\quad This is an example for manually numbered figure,
    which would not appear in the list of figures}
  \label{tab:badfigure2}
\end{figure}

Now we consider a nonlinear programming which is confronted solely with
maximizing a real-valued function with domain $\Real^n$.  Whether derivatives are
available or not, the usual strategy is first to select a point in $\Real^n$ which
is thought to be the most likely place where the maximum exists. If there is no
information available on which to base such a selection, a point is chosen at
random. From this first point an attempt is made to construct a sequence of
points, each of which yields an improved objective function value over its
predecessor. The next point to be added to the sequence is chosen by analyzing
the behavior of the function at the previous points. This construction continues
until some termination criterion is met. Methods based upon this strategy are
called {\em ascent methods}, which can be classified as {\em direct methods},
{\em gradient methods}, and {\em Hessian methods} according to the information
about the behavior of objective function $f$. Direct methods require only that
the function can be evaluated at each point. Gradient methods require the
evaluation of first derivatives of $f$. Hessian methods require the evaluation
of second derivatives. In fact, there is no superior method for all
problems. The efficiency of a method is very much dependent upon the objective
function.

\subsection{Integer Programming}

{\em Integer programming} is a special mathematical programming in which all of
the variables are assumed to be only integer values. When there are not only
integer variables but also conventional continuous variables, we call it {\em
  mixed integer programming}. If all the variables are assumed either 0 or 1,
then the problem is termed a {\em zero-one programming}. Although integer
programming can be solved by an {\em exhaustive enumeration} theoretically, it
is impractical to solve realistically sized integer programming problems. The
most successful algorithm so far found to solve integer programming is called
the {\em branch-and-bound enumeration} developed by Balas (1965) and Dakin
(1965). The other technique to integer programming is the {\em cutting plane
  method} developed by Gomory (1959).

\hfill\textit{Uncertain Programming\/}\quad(\textsl{BaoDing Liu, 2006.2})

\section*{References}
\noindent{\itshape NOTE: These references are only for demonstration. They are
  not real citations in the original text.}

\begin{translationbib}
\item Donald E. Knuth. The \TeX book. Addison-Wesley, 1984. ISBN: 0-201-13448-9
\item Paul W. Abrahams, Karl Berry and Kathryn A. Hargreaves. \TeX\ for the
  Impatient. Addison-Wesley, 1990. ISBN: 0-201-51375-7
\item David Salomon. The advanced \TeX book.  New York : Springer, 1995. ISBN:0-387-94556-3
\end{translationbib}

\chapter{外文资料的调研阅读报告或书面翻译}

\title{英文资料的中文标题}

{\heiti 摘要：} 本章为外文资料翻译内容。如果有摘要可以直接写上来，这部分好像没有
明确的规定。

\section{单目标规划}
北冥有鱼，其名为鲲。鲲之大，不知其几千里也。化而为鸟，其名为鹏。鹏之背，不知其几
千里也。怒而飞，其翼若垂天之云。是鸟也，海运则将徙于南冥。南冥者，天池也。
\begin{equation}\tag*{(123)}
 p(y|\mathbf{x}) = \frac{p(\mathbf{x},y)}{p(\mathbf{x})}=
\frac{p(\mathbf{x}|y)p(y)}{p(\mathbf{x})}
\end{equation}

吾生也有涯，而知也无涯。以有涯随无涯，殆已！已而为知者，殆而已矣！为善无近名，为
恶无近刑，缘督以为经，可以保身，可以全生，可以养亲，可以尽年。

\subsection{线性规划}
庖丁为文惠君解牛，手之所触，肩之所倚，足之所履，膝之所倚，砉然响然，奏刀騞然，莫
不中音，合于桑林之舞，乃中经首之会。
\begin{table}[ht]
\centering
  \centering
  \caption*{表~1\hskip1em 这是手动编号但不出现在索引中的一个表格例子}
  \label{tab:badtabular3}
  \begin{tabular}[c]{|m{1.5cm}|c|c|c|c|c|c|}\hline
    \multicolumn{2}{|c|}{Network Topology} & \# of nodes &
    \multicolumn{3}{c|}{\# of clients} & Server \\\hline
    GT-ITM & Waxman Transit-Stub & 600 &
    \multirow{2}{2em}{2\%}&
    \multirow{2}{2em}{10\%}&
    \multirow{2}{2em}{50\%}&
    \multirow{2}{1.2in}{Max. Connectivity}\\\cline{1-3}
    \multicolumn{2}{|c|}{Inet-2.1} & 6000 & & & &\\\hline
    \multirow{2}{1.5cm}{Xue} & Rui  & Ni &\multicolumn{4}{c|}{\multirow{2}*{\thuthesis}}\\\cline{2-3}
    & \multicolumn{2}{c|}{ABCDEF} &\multicolumn{4}{c|}{} \\\hline
\end{tabular}
\end{table}

文惠君曰：“嘻，善哉！技盖至此乎？”庖丁释刀对曰：“臣之所好者道也，进乎技矣。始臣之
解牛之时，所见无非全牛者；三年之后，未尝见全牛也；方今之时，臣以神遇而不以目视，
官知止而神欲行。依乎天理，批大郤，导大窾，因其固然。技经肯綮之未尝，而况大坬乎！
良庖岁更刀，割也；族庖月更刀，折也；今臣之刀十九年矣，所解数千牛矣，而刀刃若新发
于硎。彼节者有间而刀刃者无厚，以无厚入有间，恢恢乎其于游刃必有余地矣。是以十九年
而刀刃若新发于硎。虽然，每至于族，吾见其难为，怵然为戒，视为止，行为迟，动刀甚微，
謋然已解，如土委地。提刀而立，为之而四顾，为之踌躇满志，善刀而藏之。”

文惠君曰：“善哉！吾闻庖丁之言，得养生焉。”


\subsection{非线性规划}
孔子与柳下季为友，柳下季之弟名曰盗跖。盗跖从卒九千人，横行天下，侵暴诸侯。穴室枢
户，驱人牛马，取人妇女。贪得忘亲，不顾父母兄弟，不祭先祖。所过之邑，大国守城，小
国入保，万民苦之。孔子谓柳下季曰：“夫为人父者，必能诏其子；为人兄者，必能教其弟。
若父不能诏其子，兄不能教其弟，则无贵父子兄弟之亲矣。今先生，世之才士也，弟为盗
跖，为天下害，而弗能教也，丘窃为先生羞之。丘请为先生往说之。”
\begin{figure}[h]
  \centering
  \includegraphics{njulogo0.pdf}
  \caption*{图~1\hskip1em 这是手动编号但不出现索引中的图片的例子}
  \label{tab:badfigure3}
\end{figure}

柳下季曰：“先生言为人父者必能诏其子，为人兄者必能教其弟，若子不听父之诏，弟不受
兄之教，虽今先生之辩，将奈之何哉？且跖之为人也，心如涌泉，意如飘风，强足以距敌，
辩足以饰非。顺其心则喜，逆其心则怒，易辱人以言。先生必无往。”

孔子不听，颜回为驭，子贡为右，往见盗跖。

\subsection{整数规划}
盗跖乃方休卒徒大山之阳，脍人肝而餔之。孔子下车而前，见谒者曰：“鲁人孔丘，闻将军
高义，敬再拜谒者。”谒者入通。盗跖闻之大怒，目如明星，发上指冠，曰：“此夫鲁国之
巧伪人孔丘非邪？为我告之：尔作言造语，妄称文、武，冠枝木之冠，带死牛之胁，多辞缪
说，不耕而食，不织而衣，摇唇鼓舌，擅生是非，以迷天下之主，使天下学士不反其本，妄
作孝弟，而侥幸于封侯富贵者也。子之罪大极重，疾走归！不然，我将以子肝益昼餔之膳。”


\chapter{其它附录}
前面两个附录主要是给本科生做例子。其它附录的内容可以放到这里，当然如果你愿意，可
以把这部分也放到独立的文件中，然后将其 \cs{input} 到主文件中。

% \end{appendix}
% \end{document}
% \end{latex}
%
% \subsection{选项}
% \label{sec:option}
% \DescribeOption{language}
% 论文的主要语言（默认：中文）。可选：\option{chinese}，\option{english}。决定了封面、标题、定理的语言。
% \DescribeOption{open}
% 正规出版物的章节出现在奇数页，也就是右手边的页面，这就是 \option{right}，。在这种情况下，如果前一章的最后一页也是奇数，那么模板会自动生成一个纯粹的空白页。
% 提交的作业如果是电子稿的话，可以使用连续页，即使用\option{any}
% \DescribeOption{wide}
% 是否使用宽页面。如果生成作业的话，宽页面或许好看。
% \DescribeOption{draft}
% 是否生成水印。生成的水印为 Draft 表示此文档尚为草稿
%
% \subsection{字体配置}
% \label{sec:font-config}
% 使用\CTeX\ 默认字体配置
% \subsubsection{字体命令}
% \label{sec:fontcmds}
% \myentry{字体}
% \DescribeMacro{\songti}
% \DescribeMacro{\fangsong}
% \DescribeMacro{\heiti}
% \DescribeMacro{\kaishu}
% 用来切换宋体、仿宋、黑体、楷体四种基本字体。
% \myentry{字号}
% \DescribeMacro{\chuhao}
% \DescribeMacro{\xiaochu}
% \DescribeMacro{\yihao}
% \DescribeMacro{\xiaoyi}
% \DescribeMacro{\bahao}
% 定义字体大小，分别为
% \begin{center}
% \begin{tabular}{llllll}
% \toprule
% \cs{chuhao} & \cs{xiaochu} & \cs{yihao}  & \cs{xiaoyi}    & \cs{erhao}  & \cs{xiaoer}\\
% \cs{sanhao} & \cs{xiaosan} & \cs{sihao}  & \cs{banxiaosi} & \cs{xiaosi} & \cs{dawu}\\
% \cs{wuhao}  & \cs{xiaowu}  & \cs{liuhao} & \cs{xiaoliu}   & \cs{qihao}  & \cs{bahao}\\\bottomrule
% \end{tabular}
% \end{center}
% 使用方法为：\cs{command}\oarg{num}，其中 command 为字号命令，num 为行距。比
% 如 \cs{xiaosi}[1.5] 表示选择小四字体，行距 1.5 倍。写作指南要求表格中的字体
% 是 \cs{dawu}，模板已经设置好了。
%
% \subsection{封面信息}
% 仿照parts/examples/cover.tex 进行设置
% \subsection{表格}
% \DescribeMacro{\figoptadd}
% \DescribeMacro{\figoptaddcap}
% 定义了两个简单的表格命令
% \begin{latex}
%  \figoptadd{option}{address}
%  \figoptaddcap{option}{address}{caption}
% \end{latex}
% \subsection{图片}
% \DescribeMacro{\tabncc}
% \DescribeMacro{\tabnc}
% 定义了两个简单的图片命令
% \begin{latex}
%  \tabncc{number-of-columns}{content}{caption}
%  \tabnc{number-of-columns}{content}
% \end{latex}
% \subsection{代码}
% \njurepo 预设了如下的代码环境
% \begin{longtable}{ccccc}
% \toprule
% code & cpseudo & cplus & shell & commandshell \\
% verilog & python & latex & &\\  
% \bottomrule
% \end{longtable}
% \subsection{文字}
% \begin{latex}
% \href{link}{words} # 插入链接
% \magenta{品红色字}
% \CJKunderline{下划线字}
% \end{latex}
% 更多关于预置宏包的内容，可见 README.md 以及 \ref{sec:loadpkg}
%
%
% \section{致谢}
% 感谢以下宏包的作者，本宏包从中使用了部分代码和借鉴：
% \begin{itemize}
%  \item 清华大学\thuthesis https://github.com/xueruini/thuthesis
% \end{itemize}
% 
% \StopEventually{}
%
% \section{实现细节}
% \subsection{基本信息}
%    \begin{macrocode}
%<*cls>
\NeedsTeXFormat{LaTeX2e}
\ProvidesClass{njurepo}[2019/01/25 1.0.0 Nanjing University Report Template]
%    \end{macrocode}
%
% \subsection{定义选项}
% \label{sec:defoption}
% 使用kvoptions宏包进行选项设置
%    \begin{macrocode}
\hyphenation{NJU-repo}
\def\njurepo{\textsc{NJU}\-\textsc{repo}}
\def\thuthesis{\textsc{Thu}\-\textsc{Thesis}}
\def\version{1.1.2}
\RequirePackage{kvoptions}
\SetupKeyvalOptions{
    family=nju,
    prefix=nju@,
    setkeys=\kvsetkeys
}
\DeclareStringOption[chinese]{language}[chinese]
\DeclareStringOption[any]{open}[any]
\DeclareBoolOption{wide}
\DeclareBoolOption{color}
\DeclareBoolOption{draft}
\DeclareDefaultOption{\PassOptionsToClass{\CurrentOption}{ctexbook}}

\ProcessKeyvalOptions*
%    \end{macrocode}
%
% 检测选项是否合法
%    \begin{macrocode}
\newcommand\nju@validate@key[1]{%
  \@ifundefined{nju@\csname nju@#1\endcsname true}{%
    \ClassError{njurepo}{Invalid value '\csname nju#1\endcsname'}{}
    }{%
      \csname nju@\csname nju@#1\endcsname true\endcsname
    }
}
\newif\ifnju@chinese
\newif\ifnju@english
\nju@validate@key{language}
\newif\ifnju@any
\newif\ifnju@right
\nju@validate@key{open}
%    \end{macrocode}
% 
% 使用ctexbook宏包
%    \begin{macrocode}
\LoadClass[a4paper,openany,UTF8,zihao=-4,scheme=plain]{ctexbook}
%    \end{macrocode}
%
% \subsection{加载宏包}
% \label{sec:loadpkg}
% 用于开发的宏包
%    \begin{macrocode}
\RequirePackage{etoolbox}
\RequirePackage{ifxetex}
\RequirePackage{xparse}
%    \end{macrocode}
% 用于图片的宏包
%    \begin{macrocode}
\RequirePackage{graphicx}
\graphicspath{{resources/}}
\graphicspath{{resources/logo/}}
\RequirePackage[labelformat=simple]{subcaption}
\RequirePackage{pdfpages}
\includepdfset{fitpaper=true}
\RequirePackage{tikz}
\usetikzlibrary{decorations.pathmorphing,graphs,calc}
\RequirePackage{dirtree}
%    \end{macrocode}
% 用于表格的宏包
%    \begin{macrocode}
\RequirePackage{array}
\RequirePackage{longtable}
\RequirePackage{booktabs}
\RequirePackage{multirow}
\RequirePackage{tabularx}
\RequirePackage{diagbox}
\RequirePackage{makecell}
\RequirePackage{float}
%    \end{macrocode}
% 用于符号或数学的宏包
%    \begin{macrocode}
\RequirePackage{CJKfntef}
\RequirePackage{amsmath}
\RequirePackage[amsmath, thmmarks, hyperref]{ntheorem}
\RequirePackage{physics}
\RequirePackage{bbding,stmaryrd}
%\RequirePackage{fontawesome}
%    \end{macrocode}
% 其它宏包
%    \begin{macrocode}
\RequirePackage[sort&compress]{natbib}
%    \end{macrocode}
%
% 超链接
%    \begin{macrocode}
\RequirePackage{hyperref}
\ifxetex
  \hypersetup{%
    CJKbookmarks=true}
\else
  \hypersetup{%
    unicode=true,
    CJKbookmarks=false}
\fi
\hypersetup{%
  linktoc=all,
  bookmarksnumbered=true,
  bookmarksopen=true,
  bookmarksopenlevel=1,
  breaklinks=true,
  colorlinks=false,
  plainpages=false,
  pdfborder=0 0 0}	
\urlstyle{same}
\def\UrlBreaks{%
  \do\/%
  \do\a\do\b\do\c\do\d\do\e\do\f\do\g\do\h\do\i\do\j\do\k\do\l%
     \do\m\do\n\do\o\do\p\do\q\do\r\do\s\do\t\do\u\do\v\do\w\do\x\do\y\do\z%
  \do\A\do\B\do\C\do\D\do\E\do\F\do\G\do\H\do\I\do\J\do\K\do\L%
     \do\M\do\N\do\O\do\P\do\Q\do\R\do\S\do\T\do\U\do\V\do\W\do\X\do\Y\do\Z%
  \do0\do1\do2\do3\do4\do5\do6\do7\do8\do9\do=\do/\do.\do:%
  \do\*\do\-\do\~\do\'\do\"\do\-}
\Urlmuskip=0mu plus 0.1mu
%    \end{macrocode}
%
% 页眉页脚设置
%    \begin{macrocode}
\RequirePackage{fancyhdr}
\RequirePackage{notoccite}	
%    \end{macrocode}
%
% \subsection{页面设置}
% 使用了thuthesis的非本科生默认配置。
%    \begin{macrocode}
\RequirePackage{geometry}
\ifnju@wide 
\geometry{
    a4paper, %210*297mm
    hcentering,
    ignoreall,
    nomarginpar,
    left=10mm,
    headheight=5mm,
    headsep=5mm,
    textheight=237mm,
    bottom=29mm,
    footskip=6mm
}\else
\geometry{
    a4paper, %210*297mm
    hcentering,
    ignoreall,
    nomarginpar,
    left=30mm,
    headheight=5mm,
    headsep=5mm,
    textheight=237mm,
    bottom=29mm,
    footskip=6mm
}
\fi
%    \end{macrocode}
%
% \subsection{主文档格式}
% \label{sec:mainbody}
%
% \begin{macro}{\cleardoublepage}
%    \begin{macrocode}
\let\nju@cleardoublepage\cleardoublepage
\newcommand{\nju@clearemptydoublepage}{%
  \clearpage{\pagestyle{nju@empty}\nju@cleardoublepage}}
\let\cleardoublepage\nju@clearemptydoublepage
%    \end{macrocode}
% \end{macro}
%
% \begin{macro}{\frontmatter}
% \begin{macro}{\mainmatter}
% \begin{macro}{\backmatter}
%    \begin{macrocode}
\renewcommand\frontmatter{%
    \ifnju@right\cleardoublepage\else\clearpage\fi
    \@mainmatterfalse
    \pagenumbering{Roman}
    \pagestyle{nju@empty}}
\renewcommand\mainmatter{%
    \ifnju@right\cleardoublepage\else\clearpage\fi
    \@mainmattertrue
    \pagenumbering{arabic}
    \pagestyle{nju@headings}}
\renewcommand\backmatter{%
    \ifnju@right\cleardoublepage\else\clearpage\fi
    \@mainmattertrue}
%    \end{macrocode}
% \end{macro}
% \end{macro}
% \end{macro}
%
% \subsection{字体与字号}
% \label{sec:font}
% \subsubsection{英文字体}
% 配置英文字体。
%    \begin{macrocode}
\newcommand\nju@fontset{\csname g__ctex_fontset_tl\endcsname}
\ifthenelse{\equal{\nju@fontset}{fandol}}{
  \setmainfont[
    Extension      = .otf,
    UprightFont    = *-regular,
    BoldFont       = *-bold,
    ItalicFont     = *-italic,
    BoldItalicFont = *-bolditalic,
  ]{texgyretermes}
  \setsansfont[
    Extension      = .otf,
    UprightFont    = *-regular,
    BoldFont       = *-bold,
    ItalicFont     = *-italic,
    BoldItalicFont = *-bolditalic,
  ]{texgyreheros}
  \setmonofont[
    Extension      = .otf,
    UprightFont    = *-regular,
    BoldFont       = *-bold,
    ItalicFont     = *-italic,
    BoldItalicFont = *-bolditalic,
    Scale          = MatchLowercase,
  ]{texgyrecursor}
}{
  \setmainfont{Times New Roman}
  \setsansfont{Arial}
  \ifthenelse{\equal{\nju@fontset}{mac}}{
    \setmonofont[Scale=MatchLowercase]{Menlo}
  }{
    \setmonofont[Scale=MatchLowercase]{Courier New}
  }
}
%    \end{macrocode}
%
% \subsubsection{数学环境字体}
% 配置数学字体（使用unicode-math）
%    \begin{macrocode}
\RequirePackage{unicode-math}
\unimathsetup{
  math-style = ISO,
  bold-style = ISO,
  nabla      = upright,
  partial    = upright,
}
\IfFontExistsTF{STIX2Math.otf}{
  \setmathfont[StylisticSet=8]{STIX2Math.otf}
  \setmathfont[range={scr,bfscr},StylisticSet=1]{STIX2Math.otf}
  \IfFontExistsTF{XITSMath-Regular.otf}{
    \setmathfont[range={\partial,\lbrace,\rbrace}]{XITSMath-Regular.otf}
  }{
    \setmathfont[range={\partial,\lbrace,\rbrace}]{xits-math.otf}
  }
}{
  \setmathfont[
    Extension    = .otf,
    BoldFont     = *bold,
    StylisticSet = 8,
  ]{xits-math}
  \setmathfont[range={cal,bfcal},StylisticSet=1]{xits-math.otf}
}
%    \end{macrocode}
%
% \subsubsection{数学环境符号}
% \begin{macro}{\ldots}
% 省略号一律居中，所以 \cs{ldots} 不再居于底部。
%    \begin{macrocode}
\ifnju@chinese
  \def\mathellipsis{\cdots}
\fi
%    \end{macrocode}
% \end{macro}
%
% \begin{macro}{\le}
% \begin{macro}{\ge}
% \begin{macro}{\leq}
% \begin{macro}{\geq}
% 小于等于号要使用倾斜的形式。
%    \begin{macrocode}
\protected\def\le{\leqslant}
\protected\def\ge{\geqslant}
\AtBeginDocument{%
  \renewcommand\leq{\leqslant}%
  \renewcommand\geq{\geqslant}%
}
%    \end{macrocode}
% \end{macro}
% \end{macro}
% \end{macro}
% \end{macro}
%
% \begin{macro}{\int}
% 积分号 \cs{int} 使用正体，并且上下限默认置于积分号上下两侧。
%    \begin{macrocode}
\removenolimits{%
  \int\iint\iiint\iiiint\oint\oiint\oiiint
  \intclockwise\varointclockwise\ointctrclockwise\sumint
  \intbar\intBar\fint\cirfnint\awint\rppolint
  \scpolint\npolint\pointint\sqint\intlarhk\intx
  \intcap\intcup\upint\lowint
}
%    \end{macrocode}
% \end{macro}
%
% \begin{macro}{\Re}
% \begin{macro}{\Im}
% \begin{macro}{\nabla}
% 实部、虚部操作符使用罗马体 $\mathrm{Re}$、$\mathrm{Im}$ 而不是 fraktur 体
% $\Re$、$\Im$。\cs{nabla} 使用粗正体。
%    \begin{macrocode}
\AtBeginDocument{%
  \renewcommand{\Re}{\operatorname{Re}}%
  \renewcommand{\Im}{\operatorname{Im}}%
  \renewcommand\nabla{\mbfnabla}%
}
%    \end{macrocode}
% \end{macro}
% \end{macro}
% \end{macro}
%
% \begin{macro}{\bm}
% \begin{macro}{\boldsymbol}
% 兼容旧的粗体命令：\pkg{bm} 的 \cs{bm} 和 \pkg{amsmath} 的 \cs{boldsymbol}。
%    \begin{macrocode}
\newcommand\bm{\symbf}
\renewcommand\boldsymbol{\symbf}
%    \end{macrocode}
% \end{macro}
% \end{macro}
%
% \begin{macro}{\square}
% 兼容 \pkg{amssymb} 中的命令。
%    \begin{macrocode}
\newcommand\square{\mdlgwhtsquare}
%    \end{macrocode}
% \end{macro}
%
% 允许太长的公式断行、分页等。
%    \begin{macrocode}
\allowdisplaybreaks[4]
\renewcommand\theequation{\ifnum \c@chapter>\z@ \thechapter-\fi\@arabic\c@equation}
%    \end{macrocode}
%
% 公式距前后文的距离由 4 个参数控制，参见 \cs{normalsize} 的定义。
%    \begin{macrocode}
\def\make@df@tag{\@ifstar\nju@make@df@tag@@\make@df@tag@@@}
\def\nju@make@df@tag@@#1{\gdef\df@tag{\nju@maketag{#1}\def\@currentlabel{#1}}}
\def\nju@maketag#1{\maketag@@@{(\ignorespaces #1\unskip\@@italiccorr)}}
\def\tagform@#1{\maketag@@@{(\ignorespaces #1\unskip\@@italiccorr)\equcaption{#1}}}
%    \end{macrocode}
% 修改 \cs{tagform} 会影响 \cs{eqref}。
%    \begin{macrocode}
\renewcommand{\eqref}[1]{\textup{(\ref{#1})}}
%    \end{macrocode}
%
% \subsubsection{中文字体}
% \pkg{ctex}在微软下使用雅黑字体，在macOS下使用苹方字体。这里不做更改。
%
% \subsubsection{字号}
% WORD 中的字号对应该关系如下（1bp = 72.27/72 pt）:
% \begin{center}
% \begin{tabular}{llll}
% \toprule
% 初号 & 42bp & 14.82mm & 42.1575pt \\
% 小初 & 36bp & 12.70mm & 36.135 pt \\
% 一号 & 26bp & 9.17mm & 26.0975pt \\
% 小一 & 24bp & 8.47mm & 24.09pt \\
% 二号 & 22bp & 7.76mm & 22.0825pt \\
% 小二 & 18bp & 6.35mm & 18.0675pt \\
% 三号 & 16bp & 5.64mm & 16.06pt \\
% 小三 & 15bp & 5.29mm & 15.05625pt \\
% 四号 & 14bp & 4.94mm & 14.0525pt \\
% 小四 & 12bp & 4.23mm & 12.045pt \\
% 五号 & 10.5bp & 3.70mm & 10.59375pt \\
% 小五 & 9bp & 3.18mm & 9.03375pt \\
% 六号 & 7.5bp & 2.56mm & \\
% 小六 & 6.5bp & 2.29mm & \\
% 七号 & 5.5bp & 1.94mm & \\
% 八号 & 5bp & 1.76mm & \\\bottomrule
% \end{tabular}
% \end{center}
%
% \begin{macro}{\normalsize}
% 默认正文小四号 (12bp) 字，行距为固定值 20 bp。
%    \begin{macrocode}
\renewcommand\normalsize{%
  \@setfontsize\normalsize{12bp}{20bp}%
  \abovedisplayskip=12bp \@plus 2bp \@minus 2bp
  \abovedisplayshortskip=12bp \@plus 2bp \@minus 2bp
  \belowdisplayskip=\abovedisplayskip
  \belowdisplayshortskip=\abovedisplayshortskip}
%    \end{macrocode}
% \end{macro}
%
% \begin{macro}{\nju@def@fontsize}
% 根据习惯定义字号。用法：
% \cs{nju@def@fontsize}\marg{字号名称}\marg{磅数}
%
% 避免了字号选择和行距的紧耦合。所有字号定义时为单倍行距，并提供选项指定行距倍数。
%    \begin{macrocode}
\def\nju@def@fontsize#1#2{%
  \expandafter\newcommand\csname #1\endcsname[1][1.3]{%
    \fontsize{#2}{##1\dimexpr #2}\selectfont}}
%    \end{macrocode}
% \end{macro}
%
% \begin{macro}{\chuhao}
% \begin{macro}{\xiaochu}
% \begin{macro}{\yihao}
% \begin{macro}{\xiaoyi}
% \begin{macro}{\erhao}
% \begin{macro}{\xiaoer}
% \begin{macro}{\sanhao}
% \begin{macro}{\xiaosan}
% \begin{macro}{\sihao}
% \begin{macro}{\banxiaosi}
% \begin{macro}{\xiaosi}
% \begin{macro}{\dawu}
% \begin{macro}{\wuhao}
% \begin{macro}{\xiaowu}
% \begin{macro}{\liuhao}
% \begin{macro}{\xiaoliu}
% \begin{macro}{\qihao}
% \begin{macro}{\bahao}
% 一组字号定义。
%    \begin{macrocode}
\nju@def@fontsize{chuhao}{42bp}
\nju@def@fontsize{xiaochu}{36bp}
\nju@def@fontsize{yihao}{26bp}
\nju@def@fontsize{xiaoyi}{24bp}
\nju@def@fontsize{erhao}{22bp}
\nju@def@fontsize{xiaoer}{18bp}
\nju@def@fontsize{sanhao}{16bp}
\nju@def@fontsize{xiaosan}{15bp}
\nju@def@fontsize{sihao}{14bp}
\nju@def@fontsize{banxiaosi}{13bp}
\nju@def@fontsize{xiaosi}{12bp}
\nju@def@fontsize{dawu}{11bp}
\nju@def@fontsize{wuhao}{10.5bp}
\nju@def@fontsize{xiaowu}{9bp}
\nju@def@fontsize{liuhao}{7.5bp}
\nju@def@fontsize{xiaoliu}{6.5bp}
\nju@def@fontsize{qihao}{5.5bp}
\nju@def@fontsize{bahao}{5bp}
%    \end{macrocode}
% \end{macro}
% \end{macro}
% \end{macro}
% \end{macro}
% \end{macro}
% \end{macro}
% \end{macro}
% \end{macro}
% \end{macro}
% \end{macro}
% \end{macro}
% \end{macro}
% \end{macro}
% \end{macro}
% \end{macro}
% \end{macro}
% \end{macro}
% \end{macro}
%
%
% \subsubsection{中文标点}
%
% \newcommand\unicodechar[1]{U+#1（\symbol{"#1}）}
% 由于 Unicode 的一些标点符号是中西文混用的：
% \unicodechar{00B7}、
% \unicodechar{2013}、
% \unicodechar{2014}、
% \unicodechar{2018}、
% \unicodechar{2019}、
% \unicodechar{201C}、
% \unicodechar{201D}、
% \unicodechar{2025}、
% \unicodechar{2026}、
% \unicodechar{2E3A}，
% 所以要根据语言设置正确的字体。
% \footnote{\url{https://github.com/CTeX-org/ctex-kit/issues/389}}
% 所以要根据语言设置正确的字体。
%    \begin{macrocode}
\newcommand\nju@setchinese{%
  \xeCJKResetPunctClass
}
\newcommand\nju@setenglish{%
  \xeCJKDeclareCharClass{HalfLeft}{"2018, "201C}%
  \xeCJKDeclareCharClass{HalfRight}{
    "00B7, "2019, "201D, "2013, "2014, "2025, "2026, "2E3A,
  }%
}
\newcommand\nju@setdefaultlanguage{%
  \ifnju@chinese
    \nju@setchinese
  \else
    \nju@setenglish
  \fi
}
%    \end{macrocode}
%
% \subsection{局部设置}
% \subsubsection{页眉页脚}
% \label{sec:headerfooter}
%
% 定义页眉和页脚样式。
% \begin{macro}{\ps@nju@empty}
% \begin{macro}{\ps@nju@plain}
% \begin{macro}{\ps@nju@headings}
% \begin{itemize}
% \item \texttt{nju@empty}：页眉页脚都没有
% \item \texttt{nju@plain}：只显示页脚的页码。\cs{chapter} 自动调用
% \cs{thispagestyle\{nju@plain\}}。
% \item \texttt{nju@headings}：页眉页脚同时显示
% \end{itemize}
%    \begin{macrocode}
\fancypagestyle{nju@empty}{%
  \fancyhf{}
  \renewcommand{\headrulewidth}{0pt}
  \renewcommand{\footrulewidth}{0pt}}
\fancypagestyle{nju@plain}{%
  \fancyhead{}
  \fancyfoot[C]{\xiaowu\thepage}
  \renewcommand{\headrulewidth}{0pt}
  \renewcommand{\footrulewidth}{0pt}}
\fancypagestyle{nju@headings}{%
  \fancyhead{}
  \fancyhead[C]{\wuhao\normalfont\leftmark}
  \fancyfoot{}
  \fancyfoot[C]{\wuhao\thepage}
  \renewcommand{\headrulewidth}{0.4pt}
  \renewcommand{\footrulewidth}{0pt}}
%    \end{macrocode}
% \end{macro}
% \end{macro}
% \end{macro}
%
% \subsubsection{段落}
% \label{sec:paragraph}
%
% 全文首行缩进 2 字符，标点符号用全角
%    \begin{macrocode}
\ctexset{%
  punct=quanjiao,
  space=auto,
  autoindent=true}
%    \end{macrocode}
%
% \subsubsection{列表}
% 利用 \pkg{enumitem} 命令调整默认列表环境间的距离，以符合中文习惯。
%    \begin{macrocode}
\RequirePackage[shortlabels]{enumitem}
\RequirePackage{environ}
\setlist{nosep}
%    \end{macrocode}
%
%
% \subsubsection{脚注}
% \label{sec:footnote}
% 脚注符合中文习惯，数字带圈。
%    \begin{macrocode}
\ifthenelse{\equal{\nju@fontset}{mac}}{
  \newfontfamily\nju@circlefont{Songti SC Light}
}{
  \ifthenelse{\equal{\nju@fontset}{windows}}{
    \newfontfamily\nju@circlefont{SimSun}
  }{
    \IfFontExistsTF{XITS-Regular.otf}{
      \newfontfamily\nju@circlefont{XITS-Regular.otf}
    }{
      \newfontfamily\nju@circlefont{xits-regular.otf}
    }
  }
}
\def\nju@textcircled#1{%
  \ifnum\value{#1} >9%
    \ClassError{njurepo}%
      {Too many footnotes in this page.}{Keep footnote less than 10.}%
  \fi
  {\nju@circlefont\symbol{\numexpr\value{#1}+"245F\relax}}%
}
\renewcommand{\thefootnote}{\nju@textcircled{footnote}}
\renewcommand{\thempfootnote}{\nju@textcircled{mpfootnote}}
%    \end{macrocode}
%
% 定义脚注分割线，字号（宋体小五），以及悬挂缩进（1.5字符）。
%    \begin{macrocode}
\def\footnoterule{\vskip-3\p@\hrule\@width0.3\textwidth\@height0.4\p@\vskip2.6\p@}
\let\nju@footnotesize\footnotesize
\renewcommand\footnotesize{\nju@footnotesize\xiaowu[1.5]}
%\footnotemargin1.5em\relax
%    \end{macrocode}
%
% \cs{@makefnmark} 默认是上标样式，而在脚注部分要求为正文大小。利用\cs{patchcmd} 动态调整 \cs{@makefnmark} 的定义。
%    \begin{macrocode}
\let\nju@makefnmark\@makefnmark
\def\nju@@makefnmark{\hbox{{\normalfont\@thefnmark}}}
\pretocmd{\@makefntext}{\let\@makefnmark\nju@@makefnmark}{}{}
\apptocmd{\@makefntext}{\let\@makefnmark\nju@makefnmark}{}{}
%    \end{macrocode}
%
%
% \subsubsection{定理环境}
% \label{sec:equation}
%
% 定理标题使用黑体，正文使用宋体，冒号隔开。
%    \begin{macrocode}
\theorembodyfont{\normalfont}
\theoremheaderfont{\normalfont\heiti}
\theoremsymbol{\ensuremath{\square}}
\newtheorem*{proof}{证明}
\theoremsymbol{}
\theoremseparator{：}
\ifnju@chinese
  \theoremstyle{plain}
  \newcommand\nju@assumption@name{假设}
  \newcommand\nju@definition@name{定义}
  \newcommand\nju@proposition@name{命题}
  \newcommand\nju@lemma@name{引理}
  \newcommand\nju@theorem@name{定理}
  \newcommand\nju@axiom@name{公理}
  \newcommand\nju@corollary@name{推论}
  \newcommand\nju@exercise@name{练习}
  \newcommand\nju@example@name{例}
  \newcommand\nju@remark@name{注释}
  \newcommand\nju@conjecture@name{猜想}
  \theoremstyle{break}
  \newcommand\nju@problem@name{问题}
  \newcommand\nju@solution@name{解}
\else
  \theoremstyle{plain}
  \newcommand\nju@assumption@name{Assumption}
  \newcommand\nju@definition@name{Definition}
  \newcommand\nju@proposition@name{Proposition}
  \newcommand\nju@lemma@name{Lemma}
  \newcommand\nju@theorem@name{Theorem}
  \newcommand\nju@axiom@name{Axiom}
  \newcommand\nju@corollary@name{Corollary}
  \newcommand\nju@exercise@name{Exercise}
  \newcommand\nju@example@name{Example}
  \newcommand\nju@remark@name{Remark}
  \newcommand\nju@conjecture@name{Conjecture}
  \theoremstyle{break}
  \newcommand\nju@problem@name{Problem}
  \newcommand\nju@solution@name{Solution}
\fi
\theoremheaderfont{\bfseries}
\newtheorem{assumption}{\nju@assumption@name}[section]
\newtheorem{definition}{\nju@definition@name}[section]
\newtheorem{proposition}{\nju@proposition@name}[section]
\newtheorem{lemma}{\nju@lemma@name}[section]
\newtheorem{theorem}{\nju@theorem@name}[section]
\newtheorem{axiom}{\nju@axiom@name}[section]
\newtheorem{corollary}{\nju@corollary@name}[section]
\newtheorem{exercise}{\nju@exercise@name}[section]
\newtheorem{example}{\nju@example@name}[section]
\newtheorem{remark}{\nju@remark@name}[section]
\newtheorem{problem}{\nju@problem@name}[section]
\newtheorem{conjecture}{\nju@conjecture@name}[section]
\newtheorem{solution}{\nju@solution@name}[section]

%\RequirePackage{microtype}
\ifnju@chinese
\newcommand{\promisewords}{请独立完成作业，不得抄袭。\\若参考了其它资料，请给出引用。\\鼓励讨论，但需独立书写解题过程。}
\else
\newcommand{\promisewords}{I promise this work is done on my own with no plagiarism.}
\fi
\newcommand{\Hrule}{\noindent\rule{\linewidth}{0.5mm}}

\theorempostwork{\vspace{-0.5cm}\Hrule}
\newtheorem*{csolution}{\PencilRightDown\nju@solution@name}
\newtheorem*{nsolution}{\PencilRightDown\nju@solution@name}
\RequirePackage[listings]{tcolorbox}
\newtcolorbox{ps@problem}[1]{fonttitle=\bfseries,title=#1,before skip=0.5cm, after skip=-0.5cm}
\newenvironment{cproblem}[1][]{
    \begin{ps@problem}{\Checkmark\ \nju@problem@name\ #1}
}{
    \end{ps@problem}
}
\theoremstyle{plain}
\newtheorem*{nproblem}{\Checkmark\nju@problem@name}[section]
% \subsubsection{浮动对象}
% \label{sec:float}
% 设置浮动对象和文字之间的距离
%    \begin{macrocode}
\setlength{\floatsep}{20bp \@plus4pt \@minus1pt}
\setlength{\intextsep}{20bp \@plus4pt \@minus2pt}
\setlength{\textfloatsep}{20bp \@plus4pt \@minus2pt}
\setlength{\@fptop}{0bp \@plus1.0fil}
\setlength{\@fpsep}{12bp \@plus2.0fil}
\setlength{\@fpbot}{0bp \@plus1.0fil}
%    \end{macrocode}
%
% 下面这组命令使浮动对象的缺省值稍微宽松一点，从而防止幅度对象占据过多的文本页面，
% 也可以防止在很大空白的浮动页上放置很小的图形。
%    \begin{macrocode}
\renewcommand{\textfraction}{0.15}
\renewcommand{\topfraction}{0.85}
\renewcommand{\bottomfraction}{0.65}
\renewcommand{\floatpagefraction}{0.60}
%    \end{macrocode}
%
% 定制浮动图形和表格标题样式
% \begin{itemize}
%   \item 图表标题字体为 11pt， 这里写作大五号
%   \item 去掉图表号后面的冒号。图序与图名文字之间空一个汉字符宽度。
%   \item 图：caption 在下，段前空 6 磅，段后空 12 磅
%   \item 表：caption 在上，段前空 12 磅，段后空 6 磅
% \end{itemize}
%
%    \begin{macrocode}
\let\old@tabular\@tabular
\def\nju@tabular{\dawu[1.5]\old@tabular}
\DeclareCaptionLabelFormat{nju}{{\dawu[1.5]\normalfont #1~#2}}
\DeclareCaptionLabelSeparator{nju}{\hspace{1em}}
\DeclareCaptionFont{nju}{\dawu[1.5]}
\captionsetup{labelformat=nju,labelsep=nju,font=nju,skip=6bp}
\captionsetup[table]{position=top}
\captionsetup[figure]{position=bottom}
\captionsetup[sub]{font=nju}
\renewcommand{\thesubfigure}{(\alph{subfigure})}
\renewcommand{\thesubtable}{(\alph{subtable})}
% \renewcommand{\p@subfigure}{:}
%    \end{macrocode}
% 我们采用 \pkg{longtable} 来处理跨页的表格。同样我们需要设置其默认字体为五号。
%    \begin{macrocode}
\let\nju@LT@array\LT@array
\def\LT@array{\dawu[1.5]\nju@LT@array} % set default font size
%    \end{macrocode}
%
% \begin{macro}{\hlinewd}
% 简单的表格使用三线表推荐用 \cs{hlinewd}。如果表格比较复杂还是用 \pkg{booktabs} 的命令好一些。
%    \begin{macrocode}
\def\hlinewd#1{%
  \noalign{\ifnum0=`}\fi\hrule \@height #1 \futurelet
    \reserved@a\@xhline}
%    \end{macrocode}
% \end{macro}
%
%
% \subsubsection{章节标题}
% \label{sec:theor}
%    \begin{macrocode}
\ifnju@chinese
  \ctexset{%
    chapter/name={第,章},
    appendixname=附录,
    contentsname={目\hspace{\ccwd}录},
    listfigurename=插图索引,
    listtablename=表格索引,
    figurename=图,
    tablename=表,
    bibname=参考文献,
    indexname=索引,
  }
  \newcommand\listequationname{公式索引}
  \newcommand\equationname{公式}
\else
  \newcommand\listequationname{List of Equations}
  \newcommand\equationname{Equation}
\fi
\newcommand{\cabstractname}{摘\hspace{\ccwd}要}
\newcommand{\eabstractname}{Abstract}
\let\CJK@todaysave=\today
\def\CJK@todaysmall@short{\the\year 年 \the\month 月}
\def\CJK@todaysmall{\the\year 年 \the\month 月 \the\day 日}
\def\CJK@todaybig@short{\zhdigits{\the\year}年\zhnumber{\the\month}月}
\def\CJK@todaybig{\zhdigits{\the\year}年\zhnumber{\the\month}月\zhnumber{\the\day}日}
\def\CJK@today{\CJK@todaysmall}
\renewcommand\today{\CJK@today}
\newcommand\CJKtoday[1][1]{%
  \ifcase#1\def\CJK@today{\CJK@todaysave}
    \or\def\CJK@today{\CJK@todaysmall}
    \or\def\CJK@today{\CJK@todaybig}
  \fi}
%    \end{macrocode}
%
% \pkg{fancyhdr} 定义页眉页脚很方便，但是有一个非常隐蔽的坑。通过 \pkg{fancyhdr}
% 定义的样式在第一次被调用时会修改 \cs{chaptermark}，这会导致页眉信息错误（多余
% 章号并且英文大写）。这是因为在原始的 \file{book.cls} 中定义如下（大意）：
% \begin{latex}
% \newcommand\chaptername{Chapter}
% \newcommand\@chapapp{\chaptername}
% \def\chaptermark#1{
%   \markboth{\MakeUppercase{\@chapapp\ \thechapter}}{}}
% \end{latex}
% 很显然这个 \cs{\@chapapp} 不适合中文，因此我们使用\cs{CTEXthechapter}（
% 如，“第 x 章”），同时会将 \cs{MakeUppercase} 去掉。也就是说我们会做如下动作：
% \begin{latex}
% \renewcommand{\chaptermark}[1]{\@mkboth{\CTEXthechapter\hskip\ccwd#1}{}}
% \end{latex}
% 但，\pkg{fancyhdr} 不知何故在 \cs{ps@fancy} 中对 \cs{chaptermark} 进行重定义
% （其实一模一样），而这个 \cs{ps@fancy} 会在 \cs{fancypagestyle} 中使用，如下：
% \begin{latex}
% \newcommand{\fancypagestyle}[2]{%
%   \@namedef{ps@#1}{\let\fancy@gbl\relax#2\relax\ps@fancy}}
% \end{latex}
% 这样的话，\cs{ps@fancy} 会在 \pkg{fancyhdr} 定义的任何样式首次样被激活时调用，从
% 而覆盖我们的 \cs{chaptermark} 定义（后续样式再激活不会重复覆盖）。所以我们采用如下
% 方法解决：
%    \begin{macrocode}
\AtBeginDocument{%
  \pagestyle{nju@empty}
  \renewcommand{\chaptermark}[1]{\@mkboth{\CTEXthechapter\hskip\ccwd#1}{}}}
%    \end{macrocode}
%
% 各级标题格式设置。
% \begin{description}
% \item[section] 章序号与章名之间空一个汉字符 黑体三号字，居中书写，单倍行距，段
%   前空 24 磅，段后空 18 磅。本科要求：段前段后间距 30/20 pt，行距 20pt。但正文
%   章节 30pt 的话和样例效果不一致。
% \item[section] 一级节标题，例如：\fbox{2.1 实验装置与实验方法}。节标题序号与标
%   题名之间空一个汉字符（下同）。采用黑体四号（14pt）字居左书写，行距为固定
%   值 20 磅，段前空 24 磅，段后空 6 磅。本科：25/12 pt，行距 18pt。
% \item[subsection] 二级节标题，例如：\fbox{2.1.1 实验装置}。采用黑体 13pt 字居左
%   书写，行距为固定值 20 磅，段前空 12 磅，段后空 6 磅。本科：中文黑体 12pt 字，
%   英文 13pt 字，段间距 12/6 pt，行距 15pt。
% \item[subsubsection] 三级节标题，例如：\fbox{2.1.2.1 归纳法}。采用黑体小四号
%   （12pt）字居左书写，行距为固定值 20 磅，段前空 12 磅，段后空 6 磅。
%
% \end{description}
%    \begin{macrocode}
\newcommand\nju@chapter@titleformat[1]{%
    \ifthenelse%
      {\equal{#1}{\eabstractname}}%
      {\bfseries #1}%
      {#1}%
  }
\ctexset{%
  chapter={
    afterindent=true,
    pagestyle={nju@headings},
    beforeskip={9bp},
    aftername=\hskip\ccwd,
    afterskip={24bp},
    format={\centering\sffamily\sanhao[1]},
    nameformat=\relax,
    numberformat=\relax,
    titleformat=\nju@chapter@titleformat,
    lofskip=0pt,
    lotskip=0pt,
  },
  section={
    afterindent=true,
    beforeskip={24bp\@plus 1ex \@minus .2ex},
    afterskip={6bp\@plus .2ex},
    format={\sffamily\sihao[1.429]},
  },
  subsection={
    afterindent=true,
    beforeskip={16bp\@plus 1ex \@minus .2ex},
    afterskip={6bp \@plus .2ex},
    format={\sffamily\banxiaosi[1.538]},
    numberformat={\sffamily\banxiaosi[1.538]},
  },
  subsubsection={
    afterindent=true,
    beforeskip={16bp\@plus 1ex \@minus .2ex},
    afterskip={6bp \@plus .2ex},
    format={\sffamily\xiaosi[1.667]},
  },
  paragraph/afterindent=true,
  subparagraph/afterindent=true}
%    \end{macrocode}
%
% \begin{macro}{\nju@chapter*}
% 默认的 \cs{chapter*} 很难同时满足研究生院和本科生的论文要求。本科论文要求所有的
% 章都出现在目录里，比如摘要、Abstract、主要符号表等，所以可以简单的扩展默
% 认\cs{chapter*} 实现这个目的。但是研究生又不要这些出现在目录中，而且致谢和声明
% 部分的章名、页眉和目录都不同，所以定义一个灵活的 \cs{nju@chapter*} 专门处理这些
% 要求。
%
% \cs{nju@chapter*}\oarg{tocline}\marg{title}\oarg{header}: tocline 是出现在目录
% 中的条目，如果为空则此 chapter 不出现在目录中，如果省略表示目录出现 title；
% title 是章标题；header 是页眉出现的标题，如果忽略则取 title。通过这个宏我才真
% 正体会到 \TeX\ macro 的力量！
%    \begin{macrocode}
\newcounter{nju@bookmark}
\NewDocumentCommand\nju@chapter{s o m o}{
  \IfBooleanF{#1}{%
    \ClassError{njurepo}{You have to use the star form: \string\nju@chapter*}{}
  }%
  \ifnju@right\cleardoublepage\else\clearpage\fi\phantomsection%
  \IfValueTF{#2}{%
    \ifthenelse{\equal{#2}{}}{%
      \addtocounter{nju@bookmark}\@ne
      \pdfbookmark[0]{#3}{njuchapter.\thenju@bookmark}
    }{%
      \addcontentsline{toc}{chapter}{#3}
    }
  }{%
    \addcontentsline{toc}{chapter}{#3}
  }%
  \ctexset{chapter/beforeskip=25bp}
  \chapter*{#3}%
  \ctexset{chapter/beforeskip=15bp}
  \IfValueTF{#4}{%
    \ifthenelse{\equal{#4}{}}
    {\@mkboth{}{}}
    {\@mkboth{#4}{#4}}
  }{%
    \@mkboth{#3}{#3}
  }
}
%    \end{macrocode}
% \end{macro}
%
%
% \subsubsection{目录}
% \label{sec:toc}
% 最多 4 层，即: x.x.x.x，对应的命令和层序号分别是：
% \cs{chapter}(0), \cs{section}(1), \cs{subsection}(2), \cs{subsubsection}(3)。
%    \begin{macrocode}
\setcounter{secnumdepth}{3}
\setcounter{tocdepth}{2}
%    \end{macrocode}
%
% 每章标题行前空 6 磅，后空 0 磅。章节名中英文用 Arial 字体，页码仍用 Times。
% \begin{macro}{\tableofcontents}
%    \begin{macrocode}
\renewcommand\tableofcontents{%
  \nju@chapter*[]{\contentsname}
  \xiaosi[1.65]\@starttoc{toc}\normalsize}
%    \end{macrocode}
% 调整目录样式
%    \begin{macrocode}
\def\@pnumwidth{2em}
\def\@tocrmarg{\@pnumwidth}
\def\@dotsep{1}
\renewcommand*\l@chapter[2]{%
  \ifnum \c@tocdepth >\m@ne
    \addpenalty{-\@highpenalty}%
    \vskip 4bp \@plus\p@
    \setlength\@tempdima{4em}%
    \begingroup
      \parindent \z@ \rightskip \@pnumwidth
      \parfillskip -\@pnumwidth
      \leavevmode
      \advance\leftskip\@tempdima
      \hskip -\leftskip
      {#1}%
      \leaders\hbox{$\m@th\mkern \@dotsep mu\hbox{.}\mkern \@dotsep mu$}\hfill%
      \nobreak{#2}\par
      \penalty\@highpenalty
    \endgroup
  \fi}

\patchcmd{\@dottedtocline}{\hb@xt@\@pnumwidth}{\hbox}{}{}
\renewcommand*\l@section{%
  \@dottedtocline{1}{\ccwd}{2.1em}}
\renewcommand*\l@subsection{%
  \@dottedtocline{2}{2\ccwd}{3em}}
\renewcommand*\l@subsubsection{%
  \@dottedtocline{3}{3.5em}{3.8em}}
%    \end{macrocode}
% \end{macro}
%
% \subsection{附加页面}
% \label{sec:etc}
%
% \subsubsection{封面}
% \label{sec:cover}
% 定义封面参数。
%    \begin{macrocode}
\def\nju@def@term#1{%
  \define@key{nju}{#1}{\csname #1\endcsname{##1}}
  \expandafter\gdef\csname #1\endcsname##1{%
    \expandafter\gdef\csname nju@#1\endcsname{##1}}
  \csname #1\endcsname{}}
\nju@def@term{ctitle}
\nju@def@term{csubtitle}
\nju@def@term{csubsubtitle}
\nju@def@term{etitle}
\nju@def@term{esubtitle}
\nju@def@term{esubsubtitle}
\nju@def@term{cauthor}
\nju@def@term{csupervisor}
\nju@def@term{cassosupervisor}
\nju@def@term{ccosupervisor}
\nju@def@term{eauthor}
\nju@def@term{esupervisor}
\nju@def@term{eassosupervisor}
\nju@def@term{ecosupervisor}
\nju@def@term{cdegree}
\nju@def@term{edegree}
\nju@def@term{cdepartment}
\nju@def@term{edepartment}
\nju@def@term{cmajor}
\nju@def@term{emajor}
\nju@def@term{cdate}
\nju@def@term{edate}
\nju@def@term{stdid}
\nju@def@term{mail}
\cdate{\CJK@todaybig@short}
\edate{\ifcase \month \or January\or February\or March\or April\or May%
       \or June\or July \or August\or September\or October\or November
       \or December\fi\unskip,\ \ \the\year}
%    \end{macrocode}
%
% \begin{environment}{cabstract}
% \begin{environment}{eabstract}
% 摘要最好以环境的形式出现（否则命令的形式会导致开始结束的括号距离太远，我不喜
% 欢），这就必须让环境能够自己保存内容留待以后使用。使用 \pkg{environ} 的
% \cs{Collect@Body} 来实现。
%    \begin{macrocode}
\newcommand{\nju@@cabstract}[1]{\long\gdef\nju@cabstract{#1}}
\newenvironment{cabstract}{\Collect@Body\nju@@cabstract}{}
\newcommand{\nju@@eabstract}[1]{\long\gdef\nju@eabstract{#1}}
\newenvironment{eabstract}{\Collect@Body\nju@@eabstract}{}
%    \end{macrocode}
% \end{environment}
% \end{environment}
%
% \begin{macro}{\nju@parse@keywords}
%   不同论文格式关键词之间的分割不太相同，我们用 \cs{ckeywords} 和
%    \cs{ekeywords} 来收集关键词列表，然后用本命令来生成符合要求的格式。
%    \begin{macrocode}
\def\nju@parse@keywords#1{
  \define@key{nju}{#1}{\csname #1\endcsname{##1}}
  \expandafter\gdef\csname nju@#1\endcsname{}
  \expandafter\gdef\csname #1\endcsname##1{
    \@for\reserved@a:=##1\do{
      \expandafter\ifx\csname nju@#1\endcsname\@empty\else
        \expandafter\g@addto@macro\csname nju@#1\endcsname{%
          \ignorespaces\csname nju@#1@separator\endcsname}
      \fi
      \expandafter\expandafter\expandafter\g@addto@macro%
        \expandafter\csname nju@#1\expandafter\endcsname\expandafter{\reserved@a}}}}
%    \end{macrocode}
% \end{macro}
% \begin{macro}{\ckeywords}
% \begin{macro}{\ekeywords}
% 利用 \cs{nju@parse@keywords} 来定义，内部通过 \cs{nju@ckeywords} 和
% \cs{nju@ekeywords} 来引用。
%    \begin{macrocode}
\nju@parse@keywords{ckeywords}
\nju@parse@keywords{ekeywords}
%    \end{macrocode}
% \end{macro}
% \end{macro}
%
% \begin{macro}{\njusetup}
% 由上可见，封面和封底有一大堆信息需要设置，为了简化操作界面，提供一
% 个 \cs{njusetup} 命令支持 key/value 的方式来设置。key 就是前面各个设置项的
% 名字。\note[说明：]{只能设置普通项，不支持环境项，
% 如 \texttt{cabstract} 和 \texttt{eabstract}。} 由于这些设置项被 \cs{makecover}
% 调用，所以此命令需要在 \cs{makecover} 之前被调用。
%    \begin{macrocode}
\def\njusetup{\kvsetkeys{nju}}
%    \end{macrocode}
% \end{macro}
%
% 定义封面用到的各种文字。
%    \begin{macrocode}
\def\nju@ckeywords@separator{；}
\def\nju@ekeywords@separator{;}
\def\nju@catalog@number@title{分类号}
\def\nju@id@title{编号}
\def\nju@title@sep{：}
\def\nju@schoolname{南京大学}
\def\nju@author@title{姓名}
\def\nju@department@title{系别}
\def\nju@major@title{专业}
\def\nju@supervisor@title{指导教师}
\def\nju@assosuper@title{辅导教师}
\def\nju@studentid@title{学号}
\def\nju@date@title{日期}
\def\nju@mail@title{邮箱}
\newcommand{\nju@ckeywords@title}{关键词：}
\def\nju@title@pre{}

\def\nju@eng@title@sep{:}
\def\nju@eng@author@title{Name}
\def\nju@eng@studentid@title{StdID}
\def\nju@eng@date@title{Date}
\def\nju@eng@mail@title{E-mail}
%    \end{macrocode}
%
% 中文小型标题
%    \begin{macrocode}
\renewcommand{\maketitle}{
  \nju@setup@pdfinfo
  \begin{center} {\LARGE \ifnju@chinese\nju@ctitle\else\nju@etitle\fi}
  \end{center}
  \hspace*{\fill}
  \ifnju@chinese
    \nju@author@title\nju@title@sep\CJKunderline{\nju@cauthor}
  \else
    \nju@eng@author@title\nju@eng@title@sep\underline{\nju@eauthor}
  \fi
  \hspace*{\fill}
  \ifx\nju@stdid\@empty\relax
  \else
    \ifnju@chinese
      \nju@studentid@title\nju@title@sep\CJKunderline{\nju@stdid}
    \else
      \nju@eng@studentid@title\nju@eng@title@sep\underline{\nju@stdid}
    \fi
  \fi
  \hspace*{\fill}
  \ifnju@chinese
    \nju@date@title\nju@title@sep\CJKunderline{\today}
  \else
    \nju@eng@date@title\nju@eng@title@sep\CJKunderline{\nju@edate}
  \fi
  \hspace*{\fill}\\
}
%    \end{macrocode}
%
% 别样封面
%    \begin{macrocode}
\newcommand{\maketitlepage}{
  \nju@setup@pdfinfo
  \begin{titlepage}
    \begin{center}
    \ifx\nju@esubsubtitle\@empty\relax  {\LARGE\sffamily\scshape\ifnju@chinese\nju@csubsubtitle\else\nju@esubsubtitle\fi\ }\\[1.5cm]
    \else
    {\LARGE\sffamily\scshape \ifnju@chinese\nju@csubsubtitle\else\nju@esubsubtitle\fi}\\[1.5cm]
    \fi
		{\Large\sffamily\scshape \ifnju@chinese\nju@csubtitle\else\nju@esubtitle\fi}\\
    \rule{\linewidth}{0.5mm} \\[0.4cm]
		{\huge\sffamily\bfseries \ifnju@chinese\nju@ctitle\else\nju@etitle\fi}\\
		\rule{\linewidth}{0.5mm} \\[1.5cm]
			
		\begin{center}
			\begin{tabular}{@{\hspace{0.5cm}}l@{\hspace{0.5cm}}l}
				\nju@eauthor & \nju@stdid\\
			\end{tabular}
		\end{center}
		\vfill
		{\large \nju@edate}
    \end{center}
    \ifnju@right\cleardoublepage\else\clearpage\fi
  \end{titlepage}
}
%    \end{macrocode}
%
% \myentry{封面第一页}
% \begin{macro}{\nju@first@titlepage}
% 题名使用一号黑体字，一行写不下时可分两行写，并采用 1.25 倍行距。
% 申请学位的学科门类: 小二号宋体字。
% 中文封面页边距：
%  上- 6.0 厘米，下- 5.5 厘米，左- 4.0 厘米，右- 4.0 厘米，装订线 0 厘米；
%
%    \begin{macrocode}
\newcommand\nju@underline[2][6em]{\hskip1pt\underline{\hb@xt@ #1{\hss#2\hss}}\hskip3pt}
\newlength{\nju@title@width}
\ifxetex % todo: ugly codes
  \newcommand{\nju@put@title}[2][\nju@title@width]{%
  \begin{CJKfilltwosides}[b]{#1}#2\end{CJKfilltwosides}}
\else
  \newcommand{\nju@put@title}[2][\nju@title@width]{%
  \begin{CJKfilltwosides}{#1}#2\end{CJKfilltwosides}}
\fi
\newcommand{\nju@first@titlepage}{
  \begin{center}
    \vspace*{-1.6cm}
    \parbox[b][2.4cm][t]{\textwidth}{%
      \rule{1cm}{0cm}}
      \vskip0.65cm
      {\includegraphics[width=0.3\textwidth]{njuname0.pdf}}
      \par\vskip1.5cm
      {\xiaochu\heiti\ziju{0.5}\textbf\nju@csubtitle}
      \vskip2.2cm\hskip0.8cm
      \noindent\heiti\xiaoer\nju@title@pre
      \parbox[t]{12cm}{%
      \ignorespaces\yihao[1.51]%
      \renewcommand{\CJKunderlinebasesep}{0.25cm}%
      \renewcommand{\ULthickness}{1.3pt}%
      \ifxetex
        \xeCJKsetup{underline/format=\color{black}}%
      \else
        \def\CJKunderlinecolor{\color{black}}%
      \fi
      \centering\CJKunderline*{\nju@ctitle}
      
    }%
      \vskip1.3cm
%    \end{macrocode}
%
% 作者及导师信息部分使用三号仿宋字
%    \begin{macrocode}
      \vskip0.75cm
      \ifx\nju@cassosupervisor\@empty%
        \def\nju@tempa{7.15cm}
      \else%
        \def\nju@tempa{8.15cm}
      \fi%
      \parbox[t][\nju@tempa][t]{\textwidth}{%
        {\fangsong\sanhao[1.95]%
         \hspace*{1.9cm}
         \setlength{\nju@title@width}{4em}
         \setlength{\extrarowheight}{6pt}
         \ifxetex % todo: ugly codes
           \begin{tabular}{p{\nju@title@width}@{}l@{\extracolsep{8pt}}l}
         \else
           \begin{tabular}{p{\nju@title@width}l@{}l}
         \fi
             \nju@put@title{\nju@department@title} & \nju@title@sep
               & \nju@cdepartment\\
             \nju@put@title{\nju@major@title}      & \nju@title@sep
               & \nju@cmajor\\
             \nju@put@title{\nju@author@title}     & \nju@title@sep
               & \nju@cauthor \\
             \nju@put@title{\nju@supervisor@title} & \nju@title@sep
               & \nju@csupervisor\\
             \ifx\nju@cassosupervisor\@empty\else%
               \nju@put@title{\nju@assosuper@title} & \nju@title@sep
               & \nju@cassosupervisor\\
             \fi
           \end{tabular}
        }}
%    \end{macrocode}
%
% 论文成文打印的日期，用三号宋体汉字，不用阿拉伯数字
% 本科：论文成文打印的日期用阿拉伯数字，采用小四号宋体
%    \begin{macrocode}
     \begin{center}
       {\vskip-1.0cm\xiaosi
         \songti\nju@cdate}
     \end{center}
    \end{center}} % end of titlepage
%    \end{macrocode}
% \end{macro}
%
% \myentry{英文封面}
% \begin{macro}{\nju@engcover}
%    \begin{macrocode}
\newcommand{\nju@engcover}{%
  \begin{center}
    \vspace*{-5pt}
    \parbox[t][5.2cm][t]{\paperwidth-7.2cm}{
      \renewcommand{\baselinestretch}{1.5}
      \begin{center}
        \erhao[1.1]\bfseries\sffamily\nju@etitle%
      \end{center}}
    \parbox[t][][b]{\paperwidth-7.2cm}{
      \renewcommand{\baselinestretch}{1.3}
      \begin{center}
        \sanhao\sffamily by\\[3bp]
        \bfseries\nju@eauthor%
        \ifx\nju@emajor\empty\relax\else
          \\(~\nju@emajor~)%
        \fi
      \end{center}}
    \par\vspace{0.9cm}
    \parbox[t][2.1cm][t]{\paperwidth-7.2cm}{
      \renewcommand{\baselinestretch}{1.2}
      \xiaosan\centering
      \begin{tabular}{rl}
        Supervisor : & \nju@esupervisor\\
        \ifx\nju@eassosupervisor\@empty
          \else Associate Supervisor : & \nju@eassosupervisor\\\fi
        \ifx\nju@ecosupervisor\@empty
          \else Cooperate Supervisor : & \nju@ecosupervisor\\\fi
      \end{tabular}}
    \parbox[t][2cm][b]{\paperwidth-7.2cm}{
    \begin{center}
      \sanhao\bfseries\sffamily\nju@edate
    \end{center}}
  \end{center}}
%    \end{macrocode}
% \end{macro}
%
% \begin{macro}{\makecover}
% 生成封面总命令。
%    \begin{macrocode}
\def\makecover{%
  \nju@setup@pdfinfo\nju@makecover}
\def\nju@setup@pdfinfo{%
  \ifnju@chinese
    \hypersetup{
      pdftitle    = \nju@ctitle,
      pdfauthor   = \nju@cauthor,
      pdfsubject  = \nju@cdegree,
      pdfkeywords = \nju@ckeywords,
    }%
  \else
    \hypersetup{
      pdftitle    = \nju@etitle,
      pdfauthor   = \nju@eauthor,
      pdfsubject  = \nju@edegree,
      pdfkeywords = \nju@ekeywords,
    }%
  \fi
  \hypersetup{
    pdfcreator={\njurepo-v\version}}}
\NewDocumentCommand{\nju@makecover}{o}{
  \phantomsection
  \pdfbookmark[-1]{\nju@ctitle}{ctitle}
  \normalsize%
  \begin{titlepage}
    \ifnju@chinese
      \nju@first@titlepage
    \else
      \nju@engcover
    \fi
    \ifnju@right\cleardoublepage\else\clearpage\fi
  \end{titlepage}
}
\newcommand{\makeabstract}{
  \normalsize
  \nju@makeabstract
  \let\@tabular\nju@tabular
}
%    \end{macrocode}
% \end{macro}
%
% \subsubsection{摘要}
% \label{sec:abstractformat}
%
% \begin{macro}{\nju@put@keywords}
% 排版关键字。
%    \begin{macrocode}
\newbox\nju@kw
\newcommand\nju@put@keywords[2]{%
  \begingroup
    \setbox\nju@kw=\hbox{#1}
    \indent%
    \box\nju@kw#2\par
  \endgroup}
%    \end{macrocode}
% \end{macro}
%
% \begin{macro}{\nju@makeabstract}
% 中文摘要部分的标题为“\textbf{摘要}”，用黑体三号字。
%    \begin{macrocode}
\newcommand{\nju@makeabstract}{%
  \clearpage
  \pagestyle{nju@plain}
  \pagenumbering{Roman}
%    \end{macrocode}
%
% 摘要内容用小四号字书写，两端对齐，汉字用宋体，外文字用 Times New Roman 体，
% 标点符号一律用中文输入状态下的标点符号。
%    \begin{macrocode}
  \ifnju@chinese
    \nju@setchinese
    \nju@chapter*[]{\cabstractname} % no tocline
    \nju@cabstract
    \vskip12bp
    \nju@put@keywords{\textbf\nju@ckeywords@title}{\nju@ckeywords}
  \else
  \nju@setenglish
    \nju@chapter*[]{\eabstractname} % no tocline
    \nju@eabstract
    \vskip12bp
    \nju@put@keywords{%
      \textbf{Key Words:\enskip}}{\nju@ekeywords}%
  \fi
  \nju@setdefaultlanguage
}
%    \end{macrocode}
% \end{macro}
%
%
%
% \subsubsection{主要符号表}
% \label{sec:denotationfmt}
% \begin{environment}{denotation}
% 主要符号对照表。
%    \begin{macrocode}
\ifnju@chinese
  \newcommand\nju@denotation@name{主要符号对照表}
\else
  \newcommand\nju@denotation@name{Nomenclature}
\fi
\newenvironment{denotation}[1][2.5cm]{%
  \nju@chapter*[]{\nju@denotation@name} % no tocline
  \vskip-30bp\xiaosi[1.6]\begin{nju@denotation}[labelwidth=#1]
}{%
  \end{nju@denotation}
}
\newlist{nju@denotation}{description}{1}
\setlist[nju@denotation]{%
  nosep,
  font=\normalfont,
  align=left,
  leftmargin=!, % sum of the following 3 lengths
  labelindent=0pt,
  labelwidth=2.5cm,
  labelsep*=0.5cm,
  itemindent=0pt,
}
%    \end{macrocode}
% \end{environment}
%
% \subsubsection{致谢与声明}
% \label{sec:ackanddeclare}
%
% \begin{environment}{acknowledgement}
% 支持扫描文件替换。
%    \begin{macrocode}
\ifnju@chinese
  \newcommand\nju@ack@name{致\hspace{\ccwd}谢}
\else
  \newcommand\nju@ack@name{Acknowledgments}
\fi
\newcommand\nju@declarename{声\hspace{\ccwd}明}
\newcommand{\nju@declaretext}{本人郑重声明：所呈交的学位论文，是本人在导师指导下
  ，独立进行研究工作所取得的成果。尽我所知，除文中已经注明引用的内容外，本学位论
  文的研究成果不包含任何他人享有著作权的内容。对本论文所涉及的研究工作做出贡献的
  其他个人和集体，均已在文中以明确方式标明。}
\newcommand{\nju@signature}{签\hspace{1em}名：}
\newcommand{\nju@backdate}{日\hspace{1em}期：}
%    \end{macrocode}
%
%  \cs{cleardoublepage}，使新开章节的页码到达正确的状态。否则会因为 \cs{addcontentsline}
% 在 chapter 之前而导致目录页码错误。
% 定义致谢与声明环境。
%    \begin{macrocode}
\NewDocumentEnvironment{acknowledgement}{o}{%
    \nju@chapter*{\nju@ack@name}
  }
%    \end{macrocode}
%
% 声明部分
%    \begin{macrocode}
  {
    \ifnju@english\relax\else%
      \IfNoValueTF{#1}{%
        \nju@chapter*{\nju@declarename}
        \par{\xiaosi\parindent2em\nju@declaretext}\vskip2cm
        {\xiaosi\hfill\nju@signature\nju@underline[2.5cm]\relax%
         \nju@backdate\nju@underline[2.5cm]\relax}%
      }{%
        \includepdf[pagecommand={\thispagestyle{nju@empty}%
          \phantomsection\addcontentsline{toc}{chapter}{\nju@declarename}%
        }]{#1}%
      }%
    \fi
  }
%    \end{macrocode}
% \end{environment}
%
% \subsubsection{图表索引}
% \label{sec:threeindex}
% \begin{macro}{\listoffigures}
% \begin{macro}{\listoffigures*}
% \begin{macro}{\listoftables}
% \begin{macro}{\listoftables*}
% 定义图表以及公式目录样式。
%    \begin{macrocode}
\def\nju@starttoc#1{% #1: float type, prepend type name in \listof*** entry.
  \let\oldnumberline\numberline
  \def\numberline##1{\oldnumberline{\csname #1name\endcsname\hskip.4em ##1}}
  \@starttoc{\csname ext@#1\endcsname}
  \let\numberline\oldnumberline}
\def\nju@listof#1{% #1: float type
  \@ifstar
    {\nju@chapter*[]{\csname list#1name\endcsname}\nju@starttoc{#1}}
    {\nju@chapter*{\csname list#1name\endcsname}\nju@starttoc{#1}}}
\renewcommand\listoffigures{\nju@listof{figure}}
\renewcommand*\l@figure{\addvspace{6bp}\@dottedtocline{1}{0em}{4em}}
\renewcommand\listoftables{\nju@listof{table}}
\let\l@table\l@figure
%    \end{macrocode}
% \end{macro}
% \end{macro}
% \end{macro}
% \end{macro}
%
% \begin{macro}{\equcaption}
%   本命令只是为了生成公式列表，所以这个 caption 是假的。如果要编号最好用
%    equation 环境，如果是其它编号环境，请手动添加 \cs{equcaption}。
% 用法如下：
%
% \cs{equcaption}\marg{counter}
%
% \marg{counter} 指定出现在索引中的编号，一般取 \cs{theequation}，如果你是用
%  \pkg{amsmath} 的 \cs{tag}，那么默认是 \cs{tag} 的参数；除此之外可能需要你
% 手工指定。
%
%    \begin{macrocode}
\def\ext@equation{loe}
\def\equcaption#1{%
  \addcontentsline{\ext@equation}{equation}%
                  {\protect\numberline{#1}}}
%    \end{macrocode}
% \end{macro}
%
% \begin{macro}{\listofequations}
% \begin{macro}{\listofequations*}
% \LaTeX\ 默认没有公式索引，此处定义自己的 \cs{listofequations}。
%    \begin{macrocode}
\newcommand\listofequations{\nju@listof{equation}}
\let\l@equation\l@figure
%    \end{macrocode}
% \end{macro}
% \end{macro}
%
% \subsection{参考文献}
% \label{sec:ref}
%
% \begin{macro}{\inlinecite}
% 依赖于 \pkg{natbib} 宏包，修改其中的命令。 旧命令 \cs{onlinecite} 依然可用。
%    \begin{macrocode}
\newcommand\bibstyle@inline{\bibpunct{[}{]}{,}{n}{,}{,}}
\DeclareRobustCommand\inlinecite{\@inlinecite}
\def\@inlinecite#1{\begingroup\let\@cite\NAT@citenum\citep{#1}\endgroup}
\let\onlinecite\inlinecite
%    \end{macrocode}
% \end{macro}
%
% 参考文献的正文部分用五号字。
% 行距采用固定值 16 磅，段前空 3 磅，段后空 0 磅。
%
% 复用 \pkg{natbib} 的 \texttt{thebibliography} 环境，调整距离。
%    \begin{macrocode}
\renewcommand\bibsection{\nju@chapter*{\bibname}}
\renewcommand\bibfont{\wuhao[1.5]}
\setlength\bibhang{2\ccwd}
\addtolength{\bibsep}{-0.7em}
\setlength{\labelsep}{0.4em}
\def\@biblabel#1{[#1]\hfill}
%    \end{macrocode}
%
% 两种引用样式：
%    \begin{macrocode}
\expandafter\newcommand\csname bibstyle@numeric\endcsname{%
  \bibpunct{[}{]}{,}{s}{,}{\textsuperscript{,}}}
\expandafter\newcommand\csname bibstyle@author-year\endcsname{%
  \bibpunct{(}{)}{;}{a}{,}{,}}
%    \end{macrocode}
%
% 下面修改 \pkg{natbib} 的引用格式，主要是将页码写在上标位置。
% numeric 模式的 \cs{citet} 的页码：
%    \begin{macrocode}
\patchcmd\NAT@citexnum{%
  \@ifnum{\NAT@ctype=\z@}{%
    \if*#2*\else\NAT@cmt#2\fi
  }{}%
  \NAT@mbox{\NAT@@close}%
}{%
  \NAT@mbox{\NAT@@close}%
  \@ifnum{\NAT@ctype=\z@}{%
    \if*#2*\else\textsuperscript{#2}\fi
  }{}%
}{}{}
%    \end{macrocode}
%
% Numeric 模式的 \cs{citep} 的页码：
%    \begin{macrocode}
\renewcommand\NAT@citesuper[3]{\ifNAT@swa
  \if*#2*\else#2\NAT@spacechar\fi
\unskip\kern\p@\textsuperscript{\NAT@@open#1\NAT@@close\if*#3*\else#3\fi}%
   \else #1\fi\endgroup}
%    \end{macrocode}
%
% Author-year 模式的 \cs{citet} 的页码：
%    \begin{macrocode}
\patchcmd{\NAT@citex}{%
  \if*#2*\else\NAT@cmt#2\fi
  \if\relax\NAT@date\relax\else\NAT@@close\fi
}{%
  \if\relax\NAT@date\relax\else\NAT@@close\fi
  \if*#2*\else\textsuperscript{#2}\fi
}{}{}
%    \end{macrocode}
%
% Author-year 模式的 \cs{citep} 的页码：
%    \begin{macrocode}
\renewcommand\NAT@citesuper[3]{\ifNAT@swa
  \if*#2*\else#2\NAT@spacechar\fi
\unskip\kern\p@\textsuperscript{\NAT@@open#1\NAT@@close\if*#3*\else#3\fi}%
   \else #1\fi\endgroup}
%    \end{macrocode}
%
% 在顺序编码制下，\pkg{natbib} 只有在三个以上连续文献引用才会使用连接号，
% 这里修改为允许两个引用使用连接号。
%    \begin{macrocode}
\patchcmd{\NAT@citexnum}{%
  \ifx\NAT@last@yr\relax
    \def@NAT@last@yr{\@citea}%
  \else
    \def@NAT@last@yr{--\NAT@penalty}%
  \fi
}{%
  \def@NAT@last@yr{-\NAT@penalty}%
}{}{}
%    \end{macrocode}
%
% \subsection{附录}
% \label{sec:appendix}
% \begin{environment}{appendix}
% 主要给本科做外文翻译用。
%    \begin{macrocode}
\let\nju@appendix\appendix
\renewenvironment{appendix}{%
  \let\title\nju@appendix@title
  \nju@appendix}{%
  \let\title\@gobble}
%    \end{macrocode}
% \end{environment}
%
% \begin{macro}{\title}
% 本科外文翻译文章的标题，用法：\cs{title}\marg{资料标题}。这个命令只能在附录环
% 境下使用。
%    \begin{macrocode}
\let\title\@gobble
\newcommand{\nju@appendix@title}[1]{%
  \begin{center}
    \xiaosi[1.667] #1
  \end{center}}
%    \end{macrocode}
% \end{macro}
%
% \begin{environment}{translationbib}
% 外文资料的参考文用宋体五号字，取固定行距17pt，段前后3pt。
%    \begin{macrocode}
\newlist{translationbib}{enumerate}{1}
\setlist[translationbib]{label=[\arabic*],align=left,nosep,itemsep=6bp,
  leftmargin=10mm,labelsep=!,before=\vspace{0.5\baselineskip}\wuhao[1.3]}
%    \end{macrocode}
% \end{environment}
%\marginpar{这是边注}
%
%\subsection{颜色}
%    \begin{macrocode}
\RequirePackage{xcolor}
\definecolor{codegreen}{rgb}{0,0.6,0}
\definecolor{codegray}{rgb}{0.5,0.5,0.5}
\definecolor{codepurple}{rgb}{0.58,0,0.82}
\definecolor{backcolour}{rgb}{0.95,0.95,0.92}
\newcommand{\red}[1]{\textcolor{red}{#1}}
\newcommand{\redoverlay}[2]{\textcolor<#2>{red}{#1}}
\newcommand{\green}[1]{\textcolor{green}{#1}}
\newcommand{\greenoverlay}[2]{\textcolor<#2>{green}{#1}}
\newcommand{\blue}[1]{\textcolor{blue}{#1}}
\newcommand{\blueoverlay}[2]{\textcolor<#2>{blue}{#1}}
\newcommand{\purple}[1]{\textcolor{purple}{#1}}
\newcommand{\cyan}[1]{\textcolor{cyan}{#1}}
\newcommand{\teal}[1]{\textcolor{teal}{#1}}
\newcommand{\magenta}[1]{{\color{magenta}#1}}
\newcommand{\note}[2][Note]{{%
  \color{magenta}{\bfseries #1}\emph{#2}}}
%    \end{macrocode}
%
%\subsection{代码}
%    \begin{macrocode}
\RequirePackage{verbatim}
\RequirePackage{algorithm}
\RequirePackage{algpseudocode}
\newcommand{\pseduo}[2]{
\begin{algorithm}
	\caption{\textsc{#1}}
	\label{alg:#1}
	\begin{algorithmic}[1]
		#2
	\end{algorithmic}
\end{algorithm}
}
\RequirePackage{listings}
\lstdefinestyle{lstStyleBase}{%
   basicstyle=\small\ttfamily,
   aboveskip=\medskipamount,
   belowskip=\medskipamount,
   lineskip=0pt,
   boxpos=c,
   showlines=false,
   extendedchars=true,
   upquote=true,
   tabsize=2,
   showtabs=false,
   showspaces=false,
   showstringspaces=false,
   numbers=none,
   linewidth=\linewidth,
   xleftmargin=4pt,
   xrightmargin=0pt,
   resetmargins=false,
   breaklines=true,
   breakatwhitespace=false,
   breakindent=0pt,
   breakautoindent=true,
   columns=flexible,
   keepspaces=true,
   gobble=2,
   framesep=3pt,
   rulesep=1pt,
   framerule=1pt,
   backgroundcolor=\color{gray!5},
   stringstyle=\color{green!40!black!100},
   keywordstyle=\bfseries\color{blue!50!black},
   commentstyle=\slshape\color{black!60}
}

\newtcblisting{commandshell}{colback=black,colupper=white,colframe=yellow!75!black, listing only,listing options={style=tcblatex,language=sh},
every listing line={\textcolor{red}{\small\ttfamily\bfseries \$>}}}

\lstdefinestyle{lstStyleShell}{%
   style=lstStyleBase,
   frame=l,
   rulecolor=\color{purple},
   language=bash}

\lstdefinestyle{lstStyleLaTeX}{%
   style=lstStyleBase,
   frame=l,
   rulecolor=\color{violet},
   language=[LaTeX]TeX}

\lstdefinestyle{lstStylecdisplay}{%
  style=lstStyleBase,
  frame=tb,
  rulecolor=\color{cyan},
  keywordstyle=\color{magenta}\bfseries\ttfamily,
  commentstyle=\color{codegreen}\ttfamily,
	stringstyle=\color{codepurple}\ttfamily\sffamily,
	backgroundcolor=\color{backcolour},
	captionpos=b,
	numbers=left,
	numberstyle=\footnotesize\color{codegray},
	stepnumber=1,
  numbersep=5pt,
  language=C
}

\lstdefinestyle{lstStylecpseudo}{%
  style=lstStyleBase,
  frame=none,
  keywordstyle=\color{magenta}\bfseries\ttfamily,
  commentstyle=\color{codegreen}\ttfamily,
	stringstyle=\color{codepurple}\ttfamily\sffamily,
	captionpos=b,
	numbers=left,
	numberstyle=\footnotesize\color{codegray},
	stepnumber=1,
  numbersep=5pt,
  language=C
}

\lstdefinestyle{lstStylecplus}{%
  style=lstStyleBase,
  frame=l,
  rulecolor=\color{blue},
  language=C++
}

\lstdefinestyle{lstStyleverilog}{%
  style=lstStyleBase,
  frame=l,
  rulecolor=\color{brown},
  language=verilog
}

\lstdefinestyle{lstStylepython}{%
  style=lstStyleBase,
  frame=l,
  rulecolor=\color{pink},
  language=python
}

\lstnewenvironment{code}{\lstset{style=lstStyleBase}}{}
\lstnewenvironment{latex}{\lstset{style=lstStyleLaTeX}}{}
\lstnewenvironment{shell}{\lstset{style=lstStyleShell}}{}
\lstnewenvironment{cdisplay}{\lstset{style=lstStylecdisplay}}{}
\lstnewenvironment{cplus}{\lstset{style=lstStylecplus}}{}
\lstnewenvironment{verilog}{\lstset{style=lstStyleverilog}}{}
\lstnewenvironment{python}{\lstset{style=lstStylepython}}{}
\lstnewenvironment{cpseudo}{\lstset{style=lstStylecpseudo}}{}
%    \end{macrocode}
%
% \subsection{快速插入图片或图表}
%    \begin{macrocode}
\newcommand{\figoptadd}[2]{
	\begin{figure}[H]
		\centering
		\includegraphics[#1]{figs/#2}
	\end{figure}
}

%%%%%%%%%%%%%%%%%%%%
\newcommand{\figoptaddcap}[3]{
	\begin{figure}[H]
		\centering
		\includegraphics[#1]{figs/#2}
		\caption{#3}
		\label{fig:#2}
	\end{figure}
}
%%%%%%%%%%%%%%%%%%%
\newcommand{\tabncc}[3]{
	\begin{table}[H]
		\centering
		\begin{tabular}{|*{#1}{c|}}
		\toprule
		#2\\
		\bottomrule
	\end{tabular}
	\caption{#3}
	\label{form:#3}
\end{table}}
%%%%%%%%%%%%%%%%%%%
\newcommand{\tabnc}[2]{
	\begin{table}[H]
		\centering
		\begin{tabular}{|*{#1}{c|}}
		\toprule
		#2\\
		\bottomrule
	\end{tabular}
\end{table}}
\newcommand{\tnl}{\tabularnewline\midrule}
%    \end{macrocode}
%
% \subsection{借用dtx文件代码}
%    \begin{macrocode}
\def\cmd#1{\cs{\expandafter\cmd@to@cs\string#1}}
\def\cmd@to@cs#1#2{\char\number`#2\relax}
\DeclareRobustCommand\cs[1]{\texttt{\char`\\#1}}
\newcommand*{\meta}[1]{{%
  \ensuremath{\langle}\rmfamily\itshape#1\/\ensuremath{\rangle}}}
\providecommand\marg[1]{%
  {\ttfamily\char`\{}\meta{#1}{\ttfamily\char`\}}}
\providecommand\oarg[1]{%
  {\ttfamily[}\meta{#1}{\ttfamily]}}
\providecommand\parg[1]{%
  {\ttfamily(}\meta{#1}{\ttfamily)}}
\providecommand\pkg[1]{{\sffamily#1}}
%    \end{macrocode}
% 
% \subsection{水印}
%    \begin{macrocode}
\RequirePackage{watermark}
\ifnju@draft
\AtEndOfClass{
	\watermark{% 
		\parbox[b][\paperheight]{\paperwidth}{% 
		\vfill 
		\centering% 
		\begin{tikzpicture}[remember picture,overlay] 
			\node [rotate=45,scale=10] at ($(current page.center) +(-1cm,1cm)$) 
			{\textcolor[gray]{0.8}{DRAFT}}; 
			\node [rotate=45,scale=3] at ($(current page.center) +(1cm,-1cm)$) 
			{\textcolor[gray]{0.75}{Compile time: \the\year - \the\month - \the\day}}; 
		\end{tikzpicture}% 
		\vfill 
		}
  }
}
\fi
%    \end{macrocode}
% 补丁
%    \begin{macrocode}
\renewcommand{\thesection}{\arabic{section}}
\newcommand{\nchapter}[1]{
  {\let\clearpage\relax\par\vspace{1cm} \chapter*{\LARGE#1}}
}
\newcommand*{\rom}[1]{\expandafter\@slowromancap\romannumeral #1@}
%    \end{macrocode}
%
% \subsection{自定义代码}
%    \begin{macrocode}

\newcommand{\blankpage}{
	\clearpage
	\begin{titlepage}
		\null\vfil
		\begin{center}
			\textit{This page intentionally left blank.}
		\end{center}
	\end{titlepage}
}
\newcommand{\rmnum}[1]{\romannumeral #1}
\newcommand{\Rmnum}[1]{\expandafter\@slowromancap\romannumeral #1@}
%    \end{macrocode}
% \subsection{结束部分}
% \label{sec:finish}
%    \begin{macrocode}
\AtEndOfClass{\sloppy}
%    \end{macrocode}
%</cls> 
%
%
%
% \iffalse
%    \begin{macrocode}
%<*dtx-style>
\ProvidesPackage{dtx-style}
\RequirePackage{hypdoc}
\RequirePackage{ifthen}
\RequirePackage[UTF8,scheme=chinese]{ctex}
\RequirePackage{newpxtext}
\RequirePackage{newpxmath}
\RequirePackage[
  top=2.5cm, bottom=2.5cm,
  left=4cm, right=2cm,marginparwidth=2.6cm,marginparsep=3mm,
  headsep=3mm]{geometry}
\RequirePackage{array,longtable,booktabs}
\RequirePackage{listings}
\RequirePackage{fancyhdr}
\RequirePackage{xcolor}
\definecolor{codegreen}{rgb}{0,0.6,0}
\definecolor{codegray}{rgb}{0.5,0.5,0.5}
\definecolor{codepurple}{rgb}{0.58,0,0.82}
\definecolor{backcolour}{rgb}{0.95,0.95,0.92}
\newcommand{\red}[1]{\textcolor{red}{#1}}
\newcommand{\redoverlay}[2]{\textcolor<#2>{red}{#1}}
\newcommand{\green}[1]{\textcolor{green}{#1}}
\newcommand{\greenoverlay}[2]{\textcolor<#2>{green}{#1}}
\newcommand{\blue}[1]{\textcolor{blue}{#1}}
\newcommand{\blueoverlay}[2]{\textcolor<#2>{blue}{#1}}
\newcommand{\purple}[1]{\textcolor{purple}{#1}}
\newcommand{\cyan}[1]{\textcolor{cyan}{#1}}
\newcommand{\teal}[1]{\textcolor{teal}{#1}}
\RequirePackage{enumitem}
\RequirePackage{etoolbox}
\RequirePackage{metalogo}
\RequirePackage{mathtools}
\DeclarePairedDelimiter{\ceil}{\lceil}{\rceil}
\DeclarePairedDelimiter{\floor}{\lfloor}{\rfloor}
\DeclareMathOperator{\Hamilton}{\hat{H}} 
\ifthenelse{\equal{\@nameuse{g__ctex_fontset_tl}}{mac}}{%
  \xeCJKsetwidth{‘’“”}{1em}
}{}

\colorlet{nju@macro}{blue!60!black}
\colorlet{nju@env}{blue!70!black}
\colorlet{nju@option}{purple}
\patchcmd{\PrintMacroName}{\MacroFont}{\MacroFont\bfseries\color{nju@macro}}{}{}
\patchcmd{\PrintDescribeMacro}{\MacroFont}{\MacroFont\bfseries\color{nju@macro}}{}{}
\patchcmd{\PrintDescribeEnv}{\MacroFont}{\MacroFont\bfseries\color{nju@env}}{}{}
\patchcmd{\PrintEnvName}{\MacroFont}{\MacroFont\bfseries\color{nju@env}}{}{}

\def\DescribeOption{%
  \leavevmode\@bsphack\begingroup\MakePrivateLetters%
  \Describe@Option}
\def\Describe@Option#1{\endgroup
  \marginpar{\raggedleft\PrintDescribeOption{#1}}%
  \nju@special@index{option}{#1}\@esphack\ignorespaces}
\def\PrintDescribeOption#1{\strut \MacroFont\bfseries\sffamily\color{nju@option} #1\ }
\def\nju@special@index#1#2{\@bsphack
  \begingroup
    \HD@target
    \let\HDorg@encapchar\encapchar
    \edef\encapchar usage{%
      \HDorg@encapchar hdclindex{\the\c@HD@hypercount}{usage}%
    }%
    \index{#2\actualchar{\string\ttfamily\space#2}
           (#1)\encapchar usage}%
    \index{#1:\levelchar#2\actualchar
           {\string\ttfamily\space#2}\encapchar usage}%
  \endgroup
  \@esphack}

\lstdefinestyle{lstStyleBase}{%
   basicstyle=\small\ttfamily,
   aboveskip=\medskipamount,
   belowskip=\medskipamount,
   lineskip=0pt,
   boxpos=c,
   showlines=false,
   extendedchars=true,
   upquote=true,
   tabsize=2,
   showtabs=false,
   showspaces=false,
   showstringspaces=false,
   numbers=none,
   linewidth=\linewidth,
   xleftmargin=4pt,
   xrightmargin=0pt,
   resetmargins=false,
   breaklines=true,
   breakatwhitespace=false,
   breakindent=0pt,
   breakautoindent=true,
   columns=flexible,
   keepspaces=true,
   gobble=2,
   framesep=3pt,
   rulesep=1pt,
   framerule=1pt,
   backgroundcolor=\color{gray!5},
   stringstyle=\color{green!40!black!100},
   keywordstyle=\bfseries\color{blue!50!black},
   commentstyle=\slshape\color{black!60}}

\lstdefinestyle{lstStyleShell}{%
   style=lstStyleBase,
   frame=l,
   rulecolor=\color{purple},
   language=bash}

\lstdefinestyle{lstStyleLaTeX}{%
   style=lstStyleBase,
   frame=l,
   rulecolor=\color{violet},
   language=[LaTeX]TeX}
\lstdefinestyle{lstStylecplus}{%
   style=lstStyleBase,
   frame=l,
   rulecolor=\color{blue},
   language=C++
 }

\lstnewenvironment{latex}{\lstset{style=lstStyleLaTeX}}{}
\lstnewenvironment{shell}{\lstset{style=lstStyleShell}}{}
\lstnewenvironment{cplus}{\lstset{style=lstStylecplus}}{}

\setlist{nosep}

\DeclareDocumentCommand{\option}{m}{\textsf{#1}}
\DeclareDocumentCommand{\env}{m}{\texttt{#1}}
\DeclareDocumentCommand{\pkg}{s m}{%
  \texttt{#2}\IfBooleanF#1{\nju@special@index{package}{#2}}}
\DeclareDocumentCommand{\file}{s m}{%
  \texttt{#2}\IfBooleanF#1{\nju@special@index{file}{#2}}}
\newcommand{\myentry}[1]{%
  \marginpar{\raggedleft\color{purple}\bfseries\strut #1}}
\newcommand{\note}[2][Note]{{%
  \color{magenta}{\bfseries #1}\emph{#2}}}

\def\njurepo{\textsc{NJU}\-\textsc{repo}}
\def\thuthesis{\textsc{Thu}\-\textsc{Thesis}}
%</dtx-style>
%    \end{macrocode}
% \fi
% \Finale
